\section{Reduction to theta lifting and geometry}\label{sec:pf:galred}

\providecommand{\galredWD}{\mathrm{WD}}
\providecommand{\galredrc}{\mathrm{rc}}
\providecommand{\galredIrr}{\operatorname{Irr}}
\providecommand{\galredHT}{\mathrm{HT}}
\providecommand{\galredinf}{\operatorname{inf}}
\providecommand{\galredU}{\mathrm{U}}
\providecommand{\galredsph}{\mathrm{WD}^{\mathrm{ur}}}
\providecommand{\galredSat}{\operatorname{Sat}}
\providecommand{\galredcp}{\operatorname{cp}}
\providecommand{\galredtr}{\operatorname{tr}}
\providecommand{\galredFr}{\operatorname{Fr}}
\providecommand{\galredSp}{\operatorname{Sp}}

This section proves \cref{thm:galois:reduction}: the geometric statement $\sfname{GeoOne}$
(\cref{def:galois:geoone}) and the theta lifting statement $\sfname{TLiftK}$ (\cref{def:galois:tliftk})
together imply the statement $\sfname{Row0}(\mathrm{recT})$, a slight strengthening of
$\sfname{Row0b}(\mathrm{recT})$ (\cref{def:row0b}). Both $\sfname{GeoOne}$ and $\sfname{TLiftK}$ are hypotheses
here; they are proved in the formal development (\cref{thm:geoone,thm:galois:tlift}), and with them the
results of this section give \cref{thm:row0b}.

\paragraph{The two halves of the proof.} The two halves meet in a statement $\sfname{CoreAll}$
(\cref{pf:galred:coreall}). It asserts that every regular algebraic, polarized, cuspidal automorphic
representation $\pi$ of $\GL_a(\A_E)$, $E$ a CM field, has a \emph{weakly attached} $\ell$-adic Galois
representation (one whose Frobenius characteristic polynomials match the Satake parameters of $\pi$ at almost
all places), and that every weakly attached representation $r$ satisfies local--global compatibility at
\emph{every} place $w$ where $\pi$ is unramified: matching Frobenius eigenvalues if $w\nmid\ell$, and if
$w\mid\ell$ the de Rham property with the expected Hodge--Tate weights, crystallinity and the expected
Frobenius on the Weil--Deligne representation.

\begin{enumerate}[label=(\alph*),nosep]
\item \emph{From $\sfname{CoreAll}$ to $\sfname{Row0}$} (\cref{sec:pf:galred:rest}). The weights are read off from
  the infinitesimal characters and the polarization. The remaining work is local--global compatibility at
  ramified places $w$, in the two cases where the local parameter $y=\galredrc(\sigma^\natural)$ of a local
  component $\sigma$ has no monodromy and is either abelian, or tame and squarefree at a place with $E_w\cong\Q_g$.
  In the abelian case $\sigma^\natural$ is an irreducible principal series, and for every Weil group element $g$
  of nonzero degree a solvable CM base change $M/E$ with $g\in G_M$ makes the place unramified; $\sfname{CoreAll}$
  for the base change computes the characteristic polynomial of $r(g)$. In the tame case a counting argument for
  depth-zero supercuspidal representations shows that $\sigma$ has depth zero, and an unramified base change of
  suitable degree reduces this case to the abelian one.
\item \emph{From $\sfname{GeoOne}\wedge\sfname{TLiftK}$ to $\sfname{CoreAll}$}
  (\cref{sec:pf:galred:odd,sec:pf:galred:core}). For automorphic representations of totally definite unitary
  groups of odd rank $a$, $\sfname{GeoOne}$ produces Galois representations, of unknown dimension, whose
  Frobenius elements are annihilated by the Hecke polynomials; a \emph{division} argument extracts
  $a$-dimensional representations (``seeds''). A density argument (``congruences'') shows that every
  automorphic representation of odd rank is an $\ell$-adic limit of seeds, and the limit of the seed
  representations, constructed through pseudo-characters, is attached to it (``limits''). Two theta lifts
  ($\sfname{TLiftK}$) to odd rank and the subtraction of $\ell$-adic characters reach even rank; descent
  passes to $\GL_a$; quadratic base change and patching remove the conditions on the field; twisting passes
  from conjugate self-dual to polarized representations.
\end{enumerate}

\paragraph{What is proved here.} The deductions at the top two or three levels of the formal proof are given in
full: the main theorem and its two halves, the deduction of $\sfname{Row0}$ from $\sfname{CoreAll}$ and two
local--global statements $\sfname{R1}$, $\sfname{R2}$, the local lemmas behind the abelian and tame cases
(including the counting argument for depth-zero representations), the weights, the base change argument for
$\sfname{R1}$, the reduction of $\sfname{R2}$ to $\sfname{R1}$, the extraction of seeds from $\sfname{GeoOne}$
given an abstract division theorem, the limit argument, the passage from odd rank to all ranks, the descent to
$\GL_a$, the patching arguments, and the passage from $\sfname{Core}$ to $\sfname{CoreAll}$. The auxiliary results behind
these deductions (the abstract division theorem, the closedness of crystallinity and of Fontaine's Weil--Deligne
representation under limits, pseudo-characters, the density statements and weight families of the congruence argument,
supports under unramified base change, the counting of depth-zero representations, the construction of CM extensions with
prescribed local behaviour, and the data for the theta lifts) are proved in \cref{sec:pf:galred:aux} down to their main
ingredients, which are stated precisely and cited.

\paragraph{Organization.} \Cref{sec:pf:galred:stmts} fixes notation, makes the statements precise and proves the
main theorem modulo three theorems: \cref{pf:galred:rest} (half (a)), \cref{pf:galred:coreallof} and
\cref{pf:galred:seedcong} (half (b)). \Cref{sec:pf:galred:rest} proves \cref{pf:galred:rest}.
\Cref{sec:pf:galred:odd} treats definite unitary groups of odd rank (seeds, congruences, limits), and
\cref{sec:pf:galred:core} passes from odd rank to $\sfname{Core}$ and to $\sfname{CoreAll}$. \Cref{sec:pf:galred:aux} proves
the auxiliary results used along the way.

% ---------------------------------------------------------------------------
\subsection{Notation and statements}\label{sec:pf:galred:stmts}

\paragraph{Places and Weil groups.} We use the conventions of \cref{sec:galois:local,sec:galois:row}. For a number field
$E$, $\Omega$ is an algebraic closure of $E$, fields between $E$ and $\Omega$ are regarded as subfields of $\Omega$, and
$G_E=\Aut_E(\Omega)$ with the Krull topology. A finite place $w$ of $E$ is a nonzero prime of $\OO_E$, $E_w$ its completion and
$q_w=\Nm w$. A place $A$ of $\Omega$ above $w$ is a valuation ring of $\Omega$ lying over $w$; it has
decomposition, inertia and wild inertia groups $D_A\supset I_A\supset P_A$ in $G_E$ and a Weil group
$W_A\subset D_A$ with degree map $\deg\colon W_A\to\Z$, whose kernel is $I_A$. An \emph{arithmetic Frobenius}
at $A$ is an element $\sigma\in D_A$ acting on the residue field of $A$ by $x\mapsto x^{q_w}$; it lies in $W_A$
and has degree $1$. If $E\subset M\subset\Omega$ and $u$ is the place of $M$ below $A$, then $W_A(M)$, $I_A(M)$ denote the
corresponding subgroups of $G_M$; an element of $G_M$ is identified with the element of $G_E$ with the same action on
$\Omega$.

$\galredWD_A$ is the set of isomorphism classes of Frobenius-semisimple Weil--Deligne representations of
$W_A$ on spaces $\C^n$ (\cref{def:galois:wd}). For $x=(\rho_x,N_x)\in\galredWD_A$ and $g\in W_A$ we write
$P_x(g)=\det(X-\rho_x(g))\in\C[X]$ for the characteristic polynomial of $\rho_x(g)$. For a multiset $S$ we
write $\prod_S(X-s)$ for $\prod_{s\in S}(X-s)$.

\paragraph{Local representations and local parameter maps.} $\galredIrr(E_w)$, supercuspidal supports, segments,
Langlands quotients, depth zero and \emph{good} representations are as in \cref{def:galois:irr}. For
$x\in\galredIrr(E_w)$ we write $\supp(x)$ for its supercuspidal support (a multiset of cuspidal elements of
$\galredIrr(E_w)$); a multiset $s$ of such elements is \emph{unlinked} if $y\ne z\nu$ for all $y,z\in s$
($\nu=|\det|$). Two elements $y,z$ of $\galredIrr(E_w)$ are \emph{inertially equivalent} if $z=y\otimes\mu$ for an
unramified character $\mu$. A \emph{local parameter map} is a map $\galredrc\colon\galredIrr(E_w)\to\galredWD_A$
satisfying the laws (L1)--(L12) of \cref{def:rectexists} for some local Artin map $a$ (we then say that
$\galredrc$ \emph{has Artin map $a$}); in the formal statement of the injectivity law (L10) only $y$ is assumed
to be good: $\galredrc(x)=\galredrc(y)$ with $y$ good implies $x=y$. Local parameter maps exist
(\cref{thm:rectexists}).

\paragraph{Automorphic notions.} Cuspidal automorphic representations $\pi$ of $\GL_a(\A_E)$, their local
components $\sigma$ at finite places, the normalized classes $\sigma^\natural=\sigma\otimes|\det|^{(1-a)/2}$,
infinitesimal characters $\galredinf_\tau(\pi)$, regular algebraicity and polarized pairs $(\pi,\chi)$ are as in
\cref{def:galois:aut}. We say that $\pi$ is \emph{unramified at $w$ with Satake parameter $B$}, and write
$\galredSat_w(\pi)=B$, if the representation of $\GL_a(E_w)$ on $V_\pi$ has nonzero $\GL_a(\OO_{E_w})$-fixed
vectors and every nonzero fixed vector is a spherical vector with Satake parameter $B$, a multiset of $a$ nonzero
complex numbers; $B$ is then unique. For such $B$ we put
\[
  B^{(a)}_w=\{q_w^{(a-1)/2}b:\ b\in B\}.
\]

\paragraph{$\ell$-adic representations.} Fix a prime $\ell$ and a field isomorphism $\iota\colon\Qlbar\to\C$. A
homomorphism $r\colon G_E\to\GL_a(\Qlbar)$ is \emph{semisimple} if every $r$-stable subspace of $\Qlbar^a$ has an
$r$-stable complement. For $w\nmid\ell$, $\galredWD_\iota(r|_{W_A})\in\galredWD_A$ is the Frobenius-semisimplified
Weil--Deligne representation of \cref{def:galois:awayb}(a). For $w\mid\ell$ the clauses below are expressed
relative to \emph{frames at $A$}: a frame consists of a package of period rings with an $A_{\inf}$-structure, an
identification of it with the development's Fontaine rings, and a bridge, that is, an embedding $j$ of $\Omega$
into $\Qlbar$ with $j(E)$ dense in a finite extension $K$ of $\Ql$, together with a topological isomorphism
$D_A\cong\Gal(\Qlbar/K)$ compatible with $j$, such that $A=\{x:|j(x)|\le1\}$ (\cref{def:galois:awayb}(d)). Frames
at $A$ exist only if $w\mid\ell$. Relative to a frame, $r|_{D_A}$ may be crystalline, semistable, de Rham or
potentially semistable; for every embedding $\tau\colon E\to\Qlbar$ \emph{induced by the frame}, i.e.\ of the form
$\tau'\circ j|_E$ for an embedding $\tau'$ of $K$, it has a multiset of $\tau$-Hodge--Tate weights (when de
Rham); and there is Fontaine's Weil--Deligne representation $\galredWD(r|_{D_A})\in\galredWD_A$, transported by
$\iota$ (the trivial class if $r|_{D_A}$ is not potentially semistable). A property of $r|_{D_A}$ relative to a frame
\emph{holds at $w$} if it holds for every place $A$ of $\Omega$ above $w$ and every frame at $A$.

\paragraph{Unitary groups.} $\galredU_n$ denotes the totally definite unitary group of
\cref{sec:galois:arch}; its automorphic representations, weights and Satake parameters at split places are
recalled in \cref{sec:pf:galred:odd}.

\subsubsection{The statements \texorpdfstring{$\sfname{Row0}$}{Row0} and \texorpdfstring{$\sfname{Row0b}$}{Row0b}}

\begin{definition}[$\sfname{Away}$; $\sfname{Props}$ and $\sfname{Row0}$]\label{pf:galred:row0}
Let $x,y\in\galredWD_A$, $A$ above a finite place $w$ of $E$. We say that $w$ is a \emph{$\Q_g$-place} if
$E_w\cong\Q_g$ as topological fields for a prime $g$. $y$ is \emph{abelian} if every commutator in $W_A$ acts on
$y$ with characteristic polynomial $(X-1)^{\dim y}$, \emph{tame} if every element of $P_A$ does, and
\emph{squarefree} as in \cref{def:galois:awayb}(b). Then $\sfname{Away}(x,y)$ is the conjunction of: $\dim
x=\dim y$; if $y$ is abelian and $N_y=0$, then $x=y$; and if $w$ is a $\Q_g$-place and $y$ is tame, squarefree
and has $N_y=0$, then $N_x=0$ and $P_x(g)=P_y(g)$ for all $g\in W_A$ of degree $0$. Thus $\sfname{Away}$ is
$\sfname{AwayB}$ (\cref{def:galois:awayb}(c)) with ``$w$ is $(\dim y)$-large'' replaced by ``$w$ is a
$\Q_g$-place''.

For a recognition notion $\mathrm{Rc}$, $\sfname{Props}_{\mathrm{Rc}}(\pi,\iota,r,\mathrm{wt},\galredHT)$ is
$\sfname{PropsB}_{\mathrm{Rc}}(\pi,\iota,r,\mathrm{wt},\galredHT)$ (\cref{def:row0b}) with $\sfname{Away}$ in place of
$\sfname{AwayB}$ in (away) and (nice), and with ``$w$ is a $\Q_g$-place'' in place of ``$w$ is $a$-large'' in
(nice). $\sfname{Row0}(\mathrm{Rc})$ is $\sfname{Row0b}(\mathrm{Rc})$ with $\sfname{Props}$ in place of
$\sfname{PropsB}$. The instance used is $\mathrm{Rc}=\mathrm{recT}$: the local parameter maps.
\end{definition}

\begin{lemma}\label{pf:galred:row0b}
For every recognition notion $\mathrm{Rc}$, $\sfname{Row0}(\mathrm{Rc})$ implies $\sfname{Row0b}(\mathrm{Rc})$.
\end{lemma}

\begin{proof}
Keep $r$, $\mathrm{wt}$ and $\galredHT$. An $n$-large place is a $\Q_g$-place, so $\sfname{Away}(x,y)$ implies
$\sfname{AwayB}(x,y)$, and the hypothesis of (nice) in $\sfname{PropsB}$ implies that of (nice) in
$\sfname{Props}$; every clause of $\sfname{PropsB}$ therefore follows from the corresponding clause of
$\sfname{Props}$. (This is \lref{UpperTheory.GalStar.Row0.row0b}.)
\end{proof}

\subsubsection{Weakly attached representations and \texorpdfstring{$\sfname{CoreAll}$}{CoreAll}}

\begin{definition}[Weakly attached representations]\label{pf:galred:weak}
Let $r\colon G_E\to\GL_a(\Qlbar)$, $w$ a finite place of $E$ and $S$ a multiset of complex numbers. We say that
$r$ \emph{has geometric Frobenius multiset $S$ at $w$}, and write $\galredFr_w(r)=S$, if for every place $A$ of
$\Omega$ above $w$, $r$ is trivial on $I_A$ and $\iota\bigl(\det(X-r(\sigma^{-1}))\bigr)=\prod_S(X-s)$ for every
arithmetic Frobenius $\sigma$ at $A$. The representation $r$ is \emph{weakly attached} to a cuspidal
automorphic representation $\pi$ of $\GL_a(\A_E)$ if it is continuous and semisimple and for all but finitely
many finite places $w$, $\pi$ is unramified at $w$ and $\galredFr_w(r)=\galredSat_w(\pi)^{(a)}_w$.
\end{definition}

For a multiset $B$ of $a$ nonzero complex numbers and $A$ above $w$, let $\galredsph_A(B)\in\galredWD_A$ be the
class of $N=0$ and $\rho=\bigoplus_{b\in B}\bigl(g\mapsto(q_w^{(1-a)/2}b^{-1})^{\deg g}\bigr)$; elements of degree
$-1$ (geometric Frobenius elements) act on it with eigenvalues $B^{(a)}_w$.

\begin{definition}[The clauses at unramified places]\label{pf:galred:clauses}
Let $\pi$ be a cuspidal automorphic representation of $\GL_a(\A_E)$, $r\colon G_E\to\GL_a(\Qlbar)$, $w$ a finite
place and $B$ a multiset. The four \emph{clauses} at $(w,B)$ are:
\begin{description}[nosep]
\item[(unr)] if $w\nmid\ell$, then $\galredFr_w(r)=B^{(a)}_w$;
\item[(dR)] at $w$, $r|_{D_A}$ is de Rham, and for every embedding $\tau$ induced by the frame its $\tau$-Hodge--Tate
  multiset is $\{\frac{a-1}2-z:\ z\in\galredinf_{\iota\tau}(\pi)\}$ (whenever this is a multiset of integers);
\item[(crys)] at $w$, $r|_{D_A}$ is crystalline whenever $\pi$ has a spherical local component at $w$;
\item[(phi)] at $w$, $P_{\galredWD(r|_{D_A})}(g)=P_{\galredsph_A(B)}(g)$ for every $g\in W_A$.
\end{description}
\formallabel{GlobalTheory.RoadU.SphAt}
\end{definition}

The clauses (dR), (crys), (phi) are vacuous if $w\nmid\ell$, as there are no frames; (dR) and (crys) do not
involve $B$.

\begin{definition}[Levels; $\sfname{Ex}$, $\sfname{Core}$, $\sfname{CoreAll}$]\label{pf:galred:coreall}
Let $E$ be a CM field with maximal totally real subfield $E^+$, $\ell$ a prime, $\pi$ a cuspidal automorphic
representation of $\GL_a(\A_E)$ and $\chi$ a Hecke character of $E^+$. We say that $(\pi,\chi)$ is \emph{in scope}
\begin{description}[nosep]
\item[$\mathrm{all}$] if $\pi$ is regular algebraic and $(\pi,\chi)$ is polarized;
\item[$\mathrm{good}$] if moreover $E$ contains $\sqrt{-d}$ for some integer $d>0$, $[E^+:\Q]$ is even, and every
  place $w$ of $E$ which is not split over $E^+$ satisfies $w\nmid\ell$, is unramified over $E^+$, and $\pi$ is
  unramified at $w$; we then also say that $\pi$ is \emph{good};
\item[$\mathrm{csd}$] if moreover $\chi=\eta^a$, where $\eta$ is the quadratic Hecke character of $E^+$ attached to
  $E/E^+$; then $\pi^c\cong\pi^\vee$ ($\pi$ is conjugate self-dual).
\end{description}
A place $w$ is \emph{admissible} for the level $\mathrm{all}$ always, and for the levels $\mathrm{good}$,
$\mathrm{csd}$ if it is split over $E^+$. For a level $s$:
\begin{itemize}[nosep]
\item $\sfname{Ex}(s)$: for all $E,\Omega,\ell,\iota$, $a\ge1$ and $(\pi,\chi)$ in scope $s$, some $r$ is weakly
  attached to $\pi$;
\item $\sfname{SLC}(s,c)$, for $c\in\{\mathrm{unr},\mathrm{dR},\mathrm{crys},\mathrm{phi}\}$: for all $(\pi,\chi)$
  in scope $s$, every $r$ weakly attached to $\pi$, every place $w$ admissible for $s$ at which $\pi$ is unramified,
  and $B=\galredSat_w(\pi)$, the clause $(c)$ at $(w,B)$ holds.
\end{itemize}
$\sfname{Core}$ is $\sfname{Ex}(\mathrm{csd})\wedge\forall c\,\sfname{SLC}(\mathrm{csd},c)$, and
$\sfname{CoreAll}$ is $\sfname{Ex}(\mathrm{all})\wedge\forall c\,\sfname{SLC}(\mathrm{all},c)$.
\formallabel{GlobalTheory.RoadU.CoreAll}
\end{definition}

\begin{remark}\label{pf:galred:coreallrem}
This is the precise form of the statement $\sfname{CoreAll}$ of \cref{def:galois:core}. Two features should be
noted. The clauses are required for \emph{every} weakly attached $r$, not only for the one provided by
$\sfname{Ex}$; and the clause (phi) compares characteristic polynomials of all elements of $W_A$, not only of
Frobenius elements. Since two weakly attached representations are isomorphic (\cref{pf:galred:repiso}), the
first feature costs nothing.
\end{remark}

\subsubsection{The main theorem}

The proof of half (b) passes through the following statements about definite unitary groups of odd rank, which
are defined precisely in \cref{sec:pf:galred:oddstmts}: $\sfname{SeedFinOdd}$ (existence of seed
representations), $\sfname{CongFinOdd}$ (density of seeds) and $\sfname{LimFinOdd}$ (existence of limit
representations).

\begin{theorem}[\cref{thm:galois:reduction}]\label{pf:galred:main}
\begin{enumerate}[label=(\alph*),nosep]
\item $\sfname{CoreAll}$ implies $\sfname{Row0}(\mathrm{recT})$, hence $\sfname{Row0b}(\mathrm{recT})$.
\item $\sfname{GeoOne}\wedge\sfname{TLiftK}$ implies $\sfname{CoreAll}$.
\end{enumerate}
Consequently $\sfname{GeoOne}\wedge\sfname{TLiftK}$ implies $\sfname{Row0}(\mathrm{recT})$.
\formallabel{UpperTheory.Gal.Int.galoisLocalGlobal_of_geometry_lift}
\end{theorem}

\begin{corollary}\label{pf:galred:root}
$\sfname{GeoOne}\wedge\sfname{TLiftK}$ implies $\sfname{Row0b}(\mathrm{recT})$.
\formallabel{UpperTheory.Gal.Int.galoisLocalGlobalStrong_of_geometry_lift}
\end{corollary}

\begin{proof}[Proof of \cref{pf:galred:main,pf:galred:root}]
(a) is \cref{pf:galred:rest}, together with \cref{pf:galred:row0b}. For (b), assume $\sfname{GeoOne}$ and
$\sfname{TLiftK}$. By \cref{pf:galred:seedcong}, $\sfname{GeoOne}$ implies $\sfname{SeedFinOdd}$ and
$\sfname{CongFinOdd}$ holds; by \cref{pf:galred:coreallof}, $\sfname{SeedFinOdd}\wedge\sfname{CongFinOdd}\wedge
\sfname{TLiftK}$ implies $\sfname{CoreAll}$. The consequence follows from (a) and (b), and the corollary from
\cref{pf:galred:row0b}.
\end{proof}

The proof of (b) is the chain
\begin{multline}\label{eq:pf:galred:chain}
\sfname{SeedFinOdd}\wedge\sfname{CongFinOdd}\ \Rightarrow\ \sfname{LimFinOdd}\
\xRightarrow{\ \sfname{ClosedPhi}\ }\ \sfname{LimU}\ \Rightarrow\ \sfname{CoreUd}\
\xRightarrow{\ \sfname{TLiftK}\ }\ \sfname{AllUd}\\
\xRightarrow{\ \text{descent}\ }\ \sfname{CoreGLu}\ \xRightarrow{\ \text{quadratic base change}\ }\
\sfname{RowCsd}\ \Rightarrow\ \sfname{Core}\ \Rightarrow\ \sfname{CoreAll},
\end{multline}
whose links are \cref{pf:galred:lim,pf:galred:limU,pf:galred:joinU,pf:galred:joinAll,pf:galred:joinGL,%
pf:galred:joinUp,pf:galred:joinEnd,pf:galred:coreallofcore}.

% ---------------------------------------------------------------------------
\subsection{From \texorpdfstring{$\sfname{CoreAll}$}{CoreAll} to \texorpdfstring{$\sfname{Row0}$}{Row0}}\label{sec:pf:galred:rest}

This subsection proves half (a) of \cref{pf:galred:main}:

\begin{theorem}\label{pf:galred:rest}
$\sfname{CoreAll}$ implies $\sfname{Row0}(\mathrm{recT})$.
\formallabel{UpperTheory.Gal.Int.rest}
\end{theorem}

The proof (\cref{sec:pf:galred:restproof}) combines $\sfname{CoreAll}$ with two local--global statements
$\sfname{R1}$ and $\sfname{R2}$ (\cref{pf:galred:R1R2}), which are deduced from $\sfname{CoreAll}$ in
\cref{sec:pf:galred:R1,sec:pf:galred:R2}.

\subsubsection{Local lemmas}\label{sec:pf:galred:local}

Throughout this paragraph $E$ is a number field, $w$ a finite place of $E$, $F=E_w$ with residue field $k_F$, $A$ a
place of $\Omega$ above $w$, and $\galredrc\colon\galredIrr(F)\to\galredWD_A$ a local parameter map with Artin map
$a$. We fix an element $\varphi\in W_A$ of degree $1$; every $g\in W_A$ is $\varphi^{\deg g}t$ with $t\in I_A$.

\begin{lemma}\label{pf:galred:ine}
Let $x,y\in\galredWD_A$ with $P_x(t)=P_y(t)$ for all $t\in I_A$. Then $\dim x=\dim y$ and $P_x(g)=P_y(g)$ for
every $g\in W_A$ of degree $0$.
\formallabel{UpperTheory.Gal.Int.Row.charPol_of_ine}
\end{lemma}

\begin{proof}
An element of degree $0$ lies in $I_A$, and $\dim x$ is the degree of $P_x(1)$.
\end{proof}

\begin{lemma}\label{pf:galred:degne0}
Let $x,y\in\galredWD_A$ with $P_x(g)=P_y(g)$ for every $g\in W_A$ of nonzero degree. Then $P_x(g)=P_y(g)$ for
every $g\in W_A$.
\formallabel{GlobalTheory.AboveL.charPol_eq_of_deg_ne_zero}
\end{lemma}

\begin{proof}
Comparing degrees at $g=\varphi$ gives $\dim x=\dim y$. Equal characteristic polynomials have equal traces, so
$\galredtr\rho_x(g)=\galredtr\rho_y(g)$ whenever $\deg g\ne0$. Let $g$ have degree $0$ and put $M=\rho_x(g)$,
$\Phi=\rho_x(\varphi)$, $M'=\rho_y(g)$, $\Phi'=\rho_y(\varphi)$. The sequence
$u_i=\galredtr(M\Phi^i)-\galredtr(M'\Phi'^i)$ vanishes for $i\ge1$, as $\deg(g\varphi^i)=i$. By the Cayley--Hamilton
theorem it satisfies the linear recurrence whose characteristic polynomial is $\det(X-\Phi)\det(X-\Phi')$; the
constant term of this polynomial is nonzero because $\Phi,\Phi'$ are invertible, so the recurrence can be run
backwards and $u_0=0$. Hence $\galredtr\rho_x=\galredtr\rho_y$ on all of $W_A$. Applying this to the powers $g^j$,
$j\ge1$, of an arbitrary $g$ gives $\galredtr\rho_x(g)^j=\galredtr\rho_y(g)^j$ for all $j\ge1$, and in characteristic $0$
the power sums of the eigenvalues determine the characteristic polynomial (Newton's identities).
\end{proof}

\begin{lemma}\label{pf:galred:generic}
If $\sigma$ is a local component at $w$ of a cuspidal automorphic representation of $\GL_a(\A_E)$, then
$\sigma^\natural$ is generic.
\formallabel{UpperTheory.Gal.Int.Row.generic_algNorm}
\end{lemma}

\begin{proof}
$\sigma$ is generic by \lref{GlobalTheory.CuspRep.Generic.isGeneric_of_isLocComp}, which states that local
components of cuspidal automorphic representations of $\GL_a$ are generic (\cite{Shalika1974}), and genericity
is preserved by the twist by $|\det|^{(1-a)/2}$.
\end{proof}

\begin{lemma}\label{pf:galred:split}
Let $x\in\galredIrr(F)$ be generic with $N_{\galredrc(x)}=0$. Then $\supp(x)=\{y_1,\dots,y_k\}$ is unlinked and
$\galredrc(x)=\bigoplus_i\galredrc(y_i)$.
\formallabel{UpperTheory.Gal.Int.Row.split}
\end{lemma}

\begin{proof}
By Zelevinsky's classification of generic representations (\cite[\S9]{Zelevinsky1980}), $x=L(\Delta_1,\dots,\Delta_k)$ for pairwise
unlinked segments $\Delta_i=[\rho_i,\rho_i\nu^{m_i-1}]$. By (L3),
$\galredrc(x)=\bigoplus_i\galredrc(\rho_i)(\frac{m_i-1}2)\otimes\galredSp(m_i)$. The monodromy of
$V\otimes\galredSp(m)$ is nonzero if $V\ne0$ and $m\ge2$; as $N_{\galredrc(x)}=0$, all $m_i=1$ and
$\galredrc(x)=\bigoplus_i\galredrc(\rho_i)$. The support of
the Langlands quotient of segments of length one is the multiset of the $\rho_i$
(\lref{GlobalTheory.T.supercuspidalSupport_of_lengthOne_langlandsDatum}), and segments of length one are
unlinked if and only if no $\rho_i$ is $\rho_j\nu$
(\lref{GlobalTheory.T.unlinkedLengthOneSegments_firstTerms_not_nuTwist}).
\end{proof}

A representation $x\in\galredIrr(F)$ is \emph{plain toral} if $\supp(x)$ consists of characters
$\xi_1,\dots,\xi_n$ of $F^\times$ and is unlinked; equivalently, $x$ is an irreducible principal series.

\begin{lemma}\label{pf:galred:plaintor}
Let $x\in\galredIrr(F)$ be generic, with $\galredrc(x)$ abelian and $N_{\galredrc(x)}=0$. Then $x$ is plain toral.
\formallabel{UpperTheory.Gal.Int.Row.plainTor_of_tor}
\end{lemma}

\begin{proof}
By \cref{pf:galred:split}, $\supp(x)=\{y_1,\dots,y_k\}$ is unlinked and $\galredrc(x)=\bigoplus\galredrc(y_i)$.
For $s,t\in W_A$, $P_{\galredrc(x)}([s,t])$ is a power of $X-1$, and $P_{\galredrc(y_i)}([s,t])$ divides it. Each
$\galredrc(y_i)$ is irreducible (L4). An irreducible element of $\galredWD_A$ in which every commutator has
characteristic polynomial dividing a power of $X-1$ is one-dimensional
(\lref{UpperTheory.Zero.LLC.AbC.dim_one}: the commutators act semisimply with all eigenvalues $1$, hence
trivially; so $\rho(I_A)$ is central, consists of scalars by Schur's lemma, and $\rho(W_A)$ is abelian). By (L1),
each $y_i$ is a character.
\end{proof}

\begin{lemma}\label{pf:galred:torfix}
Let $y\in\galredWD_A$ be tame and let $\gamma\in I_A$ generate $I_A$ topologically together with $P_A$. If every
eigenvalue $z$ of $\rho_y(\gamma)$ satisfies $z^{q_w-1}=1$, then $y$ is abelian.
\formallabel{UpperTheory.Gal.Int.Row.tor_of_fix}
\end{lemma}

\begin{proof}
$\rho_y(\gamma)$ is semisimple with eigenvalues of order dividing $q_w-1$, so $\rho_y(\gamma)^{q_w-1}=1$. Since $y$ is
tame, $\rho_y(\varphi\gamma\varphi^{-1})=\rho_y(\gamma)^{q_w}=\rho_y(\gamma)$, so $\rho_y(\varphi)$ commutes with
$\rho_y(\gamma)$. For a matrix $Z$ commuting with $\rho_y(\gamma)$, the set of $t\in I_A$ such that $\rho_y(t)$ commutes
with $Z$ is a closed subgroup ($\rho_y$ is trivial on an open subgroup of $I_A$) containing $\gamma$ and $P_A$ (which acts
trivially), hence equal to $I_A$. Applied to $Z=\rho_y(\varphi)$ and to $Z=\rho_y(t)$, $t\in I_A$, this shows that
$\rho_y(W_A)=\langle\rho_y(\varphi)\rangle\rho_y(I_A)$ is commutative.
\end{proof}

\begin{lemma}[Counting depth-zero representations]\label{pf:galred:dzcusp}
Assume that $w$ is a $\Q_g$-place, and let $x\in\galredIrr(F)$ be cuspidal with $\galredrc(x)$ tame. Then $x$ has
depth zero.
\formallabel{UpperTheory.Gal.Int.Row.dz_of_cusp}
\end{lemma}

\begin{proof}
Let $n=n(x)$ and let $\mathcal L$ be the finite set of monic irreducible polynomials of degree $n$ over $k_F$
different from $X$. Fix $\phi\in W_A$ of degree $-1$; then $\varpi=a(\phi)$ is a uniformizer. Let $c\in\C^\times$ be
the product of the eigenvalues of $\rho_{\galredrc(x)}(\phi)$. We use two counting statements:
\begin{enumerate}[label=(\roman*),nosep]
\item (\lref{PositivityTheory.Rg.DzCount.depth_zero_intermediate_holds}) There is a map $f$ from
  $\galredWD_A$ to $k_F[X]$ which maps the set $\mathcal T$ of irreducible tame $y\in\galredWD_A$ of dimension
  $n$ with $\det\rho_y(\phi)=c$ injectively into $\mathcal L$.
\item (\cref{pf:galred:dzcount}) There is a map $Y\colon\mathcal L\to\galredIrr(F)$,
  injective on $\mathcal L$, such that each $Y(P)$ is a cuspidal representation of $\GL_n(F)$ of depth zero
  whose central character $\omega$ satisfies $\omega(\varpi)=c$.
\end{enumerate}
As $F\cong\Q_g$, every cuspidal representation of depth zero is good. For $P\in\mathcal L$, $\galredrc(Y(P))$ is
irreducible (L4), tame (L11), of dimension $n$ (L1), and $\det\rho_{\galredrc(Y(P))}(\phi)=\omega(a(\phi))=c$ (L8); so
$\galredrc\circ Y$ maps $\mathcal L$ into $\mathcal T$, injectively by (L10). Hence $f\circ\galredrc\circ Y$ maps the
finite set $\mathcal L$ injectively into itself, and is therefore bijective. Now $\galredrc(x)\in\mathcal T$: it is
irreducible by (L4), tame by hypothesis, of dimension $n$ by (L1), and $\det\rho_{\galredrc(x)}(\phi)=c$ by the choice of
$c$. So $f(\galredrc(x))=f(\galredrc(Y(P)))$ for some $P$; injectivity of $f$ on $\mathcal T$ gives
$\galredrc(x)=\galredrc(Y(P))$, and (L10) gives $x=Y(P)$, which has depth zero.
\end{proof}

\begin{lemma}\label{pf:galred:inert}
Let $y,z\in\galredIrr(F)$ with $y$ good and $z=y\otimes\mu$ for an unramified character $\mu$. Then
$P_{\galredrc(z)}(t)=P_{\galredrc(y)}(t)$ for all $t\in I_A$.
\formallabel{UpperTheory.Gal.Int.Row.charPol_of_inert}
\end{lemma}

\begin{proof}
By (L7) and (L2), $\galredrc(z)=\galredrc(y)\otimes\langle\mu\circ a\rangle$, and $\mu(a(t))=1$ for $t\in I_A$
because $a(t)$ is a unit.
\end{proof}

A representation $x\in\galredIrr(F)$ is \emph{apart} if the members of $\supp(x)$ of dimension $\ge2$ are pairwise
(as members of the multiset) not inertially equivalent.

\begin{lemma}\label{pf:galred:mixed}
Assume that $w$ is a $\Q_g$-place, and let $x\in\galredIrr(F)$ be generic such that $y=\galredrc(x)$ is tame,
squarefree and has $N_y=0$. Then $x$ has depth zero, $\supp(x)$ is unlinked, and $x$ is apart.
\formallabel{UpperTheory.Gal.Int.Row.mixed_of_tame}
\end{lemma}

\begin{proof}
By \cref{pf:galred:split}, $\supp(x)=\{y_1,\dots,y_k\}$ is unlinked and $y=\bigoplus\galredrc(y_i)$; a summand of a
tame representation is tame, so each $y_i$ has
depth zero by \cref{pf:galred:dzcusp}. A representation whose supercuspidal support consists of depth-zero
representations has depth zero (\cite{MoyPrasad1996}; in the development
\lref{GlobalTheory.DepthZeroPass.depthZeroInduction}), so $x$ has depth zero. Supercuspidal supports are
unique (\lref{UpperTheory.Zero.LLC.goodSeg}), so every support of $x$ is $\{y_1,\dots,y_k\}$.

Suppose that $y_i,y_j$ ($i\ne j$) have dimension $\ge2$ and $y_j=y_i\otimes\mu$ with $\mu$ unramified. Choose
$\gamma\in I_A$ generating $I_A$ topologically together with $P_A$ (\lref{GlobalTheory.exists_tameGenAt_of_isPlaceOver}).
As $y_i$ is good, \cref{pf:galred:inert} gives $P_{\galredrc(y_j)}(\gamma)=P_{\galredrc(y_i)}(\gamma)$, so
$P_{\galredrc(y_i)}(\gamma)^2$ divides $P_y(\gamma)$, and every root of $P_{\galredrc(y_i)}(\gamma)$ is a multiple root of
$P_y(\gamma)$, hence a $(q_w-1)$-th root of unity because $y$ is squarefree. By \cref{pf:galred:torfix},
$\galredrc(y_i)$ is abelian; being irreducible it is one-dimensional (\lref{UpperTheory.Zero.LLC.AbC.dim_one}),
so $n(y_i)=1$ by (L1), a contradiction.
\end{proof}

\begin{lemma}\label{pf:galred:guard}
Let $\sigma$ be a local component at a $\Q_g$-place $w$ of a cuspidal automorphic representation of
$\GL_a(\A_E)$, $a\ge1$, such that $y=\galredrc(\sigma^\natural)$ is tame, squarefree and has $N_y=0$. Then $\sigma$
has nonzero vectors fixed by $K_1=\ker(\GL_a(\OO_F)\to\GL_a(k_F))$, and $\sigma^\natural$ is apart.
\formallabel{UpperTheory.Gal.Int.Row.guard_tame}
\end{lemma}

\begin{proof}
Apply \cref{pf:galred:mixed} to $x=\sigma^\natural$, which is generic by \cref{pf:galred:generic}. Depth zero is
preserved by unramified twists (\lref{UpperTheory.Zero.LLC.twTame}), in particular by the twist by
$|\det|^{(a-1)/2}$, so $\sigma$ has depth zero.
\end{proof}

\subsubsection{Weights}

\begin{lemma}\label{pf:galred:isweight}
Let $(\pi,\chi)$ be in scope $\mathrm{all}$, $a\ge1$. There is $\mathrm{wt}\in\Z$ such that
$\galredinf_{\tau c}(\pi)=\{a-1-\mathrm{wt}-z:\ z\in\galredinf_\tau(\pi)\}$ for every embedding $\tau\colon E\to\C$.
\formallabel{GlobalTheory.PolarizedWeight.exists_isWeight}
\end{lemma}

\begin{proof}
The Capelli elements commute with all operators on $V_\pi$, so $\pi$ has an infinitesimal character at every $\tau$
(\lref{GlobalTheory.CuspRep.hasInfChar_infChar_of_commute}). Let $\psi=\chi\circ\Nm_{E/E^+}$, and for a Hecke
character $\theta$ of $E$ let $e_\tau(\theta)\in\C$ be its exponent at $\tau$, so that twisting by $\theta$ shifts
$\galredinf_\tau$ by $e_\tau(\theta)$. The polarization is an invariant pairing, with respect to complex conjugation,
between $\pi\otimes\psi^{-1}$ and $\pi$, and such a pairing reverses the signs of the infinitesimal characters
(\lref{GlobalTheory.T.infChar_conjEmbedding_neg_of_invariant_pairing}). Hence
\begin{equation}\label{eq:pf:galred:infconj}
  \galredinf_{\bar\tau}(\pi)=\{-(z+e_\tau(\psi^{-1})):\ z\in\galredinf_\tau(\pi)\}.
\end{equation}
By regular algebraicity, the elements of $\galredinf_\tau(\pi)$ and of $\galredinf_{\bar\tau}(\pi)$ lie in
$\frac{a-1}2+\Z$; comparing an element $z$ of $\galredinf_\tau(\pi)$ with the element $-(z+e_\tau(\psi^{-1}))$ of
$\galredinf_{\bar\tau}(\pi)$ shows that $k(\tau)=e_\tau(\psi)=-e_\tau(\psi^{-1})$ is an integer. Moreover
$e_{\bar\tau}(\psi)=e_\tau(\psi)$ (\lref{GlobalTheory.T.exponent_heckeChar_compNorm_conjEmbedding}), so
$k(\tau)=n(\tau|_{E^+})$ for a function $n$ on the embeddings $\phi$ of $E^+$. For a totally positive $x\in E^{+\times}$,
let $u$ be the archimedean idele of $E$ with coordinates $\exp(\frac12\log\phi(x))$; then $\Nm_{E/E^+}(u)=x$ at
infinity, and $\chi_\infty(x)=\psi_\infty(u)=\prod_\phi\phi(x)^{n(\phi)}$. A Hecke character of a totally real field with
this property has parallel weight, i.e.\ $n$ is constant (\lref{UpperTheory.heckeChar_parallel_weight}). Put
$\mathrm{wt}=a-1-n$. Then $e_\tau(\psi^{-1})=-n$ for all $\tau$, and \eqref{eq:pf:galred:infconj} reads
$\galredinf_{\bar\tau}(\pi)=\{n-z\}=\{a-1-\mathrm{wt}-z\}$. Finally $\bar\tau=\tau\circ c$.
\end{proof}

\begin{corollary}\label{pf:galred:weights}
Let $(\pi,\chi)$ be in scope $\mathrm{all}$, $a\ge1$, and $\iota\colon\Qlbar\to\C$. There are $\mathrm{wt}\in\Z$ and
multisets $\galredHT_\tau$ of integers, $\tau\colon E\to\Qlbar$, such that
$\galredHT_\tau=\{\frac{a-1}2-z:z\in\galredinf_{\iota\tau}(\pi)\}$ and $\galredHT_{\tau c}=\{\mathrm{wt}-h:h\in\galredHT_\tau\}$.
\formallabel{UpperTheory.Gal.Int.Row.exists_weights}
\end{corollary}

\begin{proof}
By regular algebraicity $\frac{a-1}2-z\in\Z$ for $z\in\galredinf_{\iota\tau}(\pi)$, which defines $\galredHT_\tau$.
With $\mathrm{wt}$ from \cref{pf:galred:isweight}, $\galredHT_{\tau c}=\{\frac{a-1}2-(a-1-\mathrm{wt}-z)\}=
\{\mathrm{wt}-(\frac{a-1}2-z)\}$.
\end{proof}

\subsubsection{Two local--global statements and the deduction of \texorpdfstring{$\sfname{Row0}$}{Row0}}
\label{sec:pf:galred:restproof}

\begin{definition}[$\sfname{R1}$, $\sfname{R2}$]\label{pf:galred:R1R2}
Consider $E,\Omega,\ell,\iota$, $n\ge1$, $(\pi,\chi)$ in scope $\mathrm{all}$ with $\pi$ on $\GL_n$, $r$ weakly
attached to $\pi$, a finite place $w$ and a local component $\sigma$ of $\pi$ at $w$.
\begin{itemize}[nosep]
\item $\sfname{R1}$ states: if $\sigma^\natural$ is plain toral, then
  \begin{enumerate}[label=(\roman*),nosep]
  \item if $w\nmid\ell$: for every $A$ above $w$, $x=\galredWD_\iota(r|_{W_A})$ has $N_x=0$ and $x=\galredrc(\sigma^\natural)$
    for every local parameter map $\galredrc\colon\galredIrr(E_w)\to\galredWD_A$;
  \item if $w\mid\ell$: the clause (dR) holds at $w$, and at $w$, $r|_{D_A}$ is potentially semistable,
    $x=\galredWD(r|_{D_A})$ has $N_x=0$, and $P_x(g)=P_{\galredrc(\sigma^\natural)}(g)$ for all $g\in W_A$ and all local
    parameter maps $\galredrc$.
  \end{enumerate}
\item $\sfname{R2}$ states: if $\sigma$ has nonzero $K_1$-fixed vectors, and $\sigma^\natural$ is plain toral or [$w$ is a
  $\Q_g$-place and $\sigma^\natural$ is apart], then for every $A$ above $w$ and every local parameter map $\galredrc$ with
  $N_{\galredrc(\sigma^\natural)}=0$:
  \begin{enumerate}[label=(\roman*),nosep]
  \item if $w\nmid\ell$: $x=\galredWD_\iota(r|_{W_A})$ has $N_x=0$ and $P_x(t)=P_{\galredrc(\sigma^\natural)}(t)$ for
    $t\in I_A$;
  \item if $w\mid\ell$: (dR) holds at $w$, and for every frame at $A$, $r|_{D_A}$ is potentially semistable,
    $x=\galredWD(r|_{D_A})$ has $N_x=0$ and $P_x(t)=P_{\galredrc(\sigma^\natural)}(t)$ for $t\in I_A$.
  \end{enumerate}
\end{itemize}
\formallabel{GlobalTheory.DeRhamOfPlainTorus.R1}
\formallabel{UpperTheory.Gal.ToralGalois.R2.LocalGlobalCompatAtIwahoriFixedPlaces}
\end{definition}

\begin{proposition}\label{pf:galred:deduce}
$\sfname{CoreAll}\wedge\sfname{R1}\wedge\sfname{R2}$ implies $\sfname{Row0}(\mathrm{recT})$.
\formallabel{UpperTheory.Gal.Int.Row.galoisLocalGlobal_of_localGlobalCompatibility}
\end{proposition}

\begin{proof}
Let $E,\Omega,\ell,\iota$, $a\ge1$ and $(\pi,\chi)$ be as in $\sfname{Row0}$; then $(\pi,\chi)$ is in scope
$\mathrm{all}$. By $\sfname{Ex}(\mathrm{all})$ there is $r$ weakly attached to $\pi$; it is continuous and
semisimple. Take $\mathrm{wt}$ and $\galredHT$ from \cref{pf:galred:weights}, which gives (weights). Let $w$ be a
finite place, $A$ above $w$, $\sigma$ a local component of $\pi$ at $w$, $\galredrc$ a local parameter map and
$y=\galredrc(\sigma^\natural)$; recall that $\sigma^\natural$ is generic (\cref{pf:galred:generic}).

\emph{(away)}, $w\nmid\ell$, $x=\galredWD_\iota(r|_{W_A})$. $\dim x=a=\dim y$ by construction of $x$ and (L1). If $y$
is abelian with $N_y=0$, then $\sigma^\natural$ is plain toral (\cref{pf:galred:plaintor}) and $x=y$ by
$\sfname{R1}$(i). If $w$ is a $\Q_g$-place and $y$ is tame, squarefree with $N_y=0$, then by \cref{pf:galred:guard}
$\sigma$ has nonzero $K_1$-fixed vectors and $\sigma^\natural$ is apart, so $\sfname{R2}$(i) gives $N_x=0$ and $P_x=P_y$
on $I_A$, hence on all elements of degree $0$ (\cref{pf:galred:ine}).

\emph{(crys).} If $\sigma$ is spherical, then $\pi$ is unramified at $w$
(a spherical local component has a Satake parameter, which is then the Satake parameter of
$\pi$ at $w$), and $\sfname{SLC}(\mathrm{all},\mathrm{crys})$ and
$\sfname{SLC}(\mathrm{all},\mathrm{dR})$ give, for every frame, that $r|_{D_A}$ is crystalline and de Rham with
$\tau$-Hodge--Tate weights $\galredHT_\tau$.

\emph{(nice).} Assume $N_y=0$ and that $y$ is abelian, or $w$ is a $\Q_g$-place and $y$ is tame and squarefree; fix a
frame at $A$ and put $x=\galredWD(r|_{D_A})$. If $y$ is abelian, $\sigma^\natural$ is plain toral and $\sfname{R1}$(ii)
gives (dR) (hence the de Rham property with weights $\galredHT_\tau$), potential semistability, $N_x=0$ and $P_x=P_y$
on $W_A$; two Frobenius-semisimple Weil--Deligne representations without monodromy and with the same characteristic
polynomials are equal (\lref{GlobalTheory.WeakTheRep.eq_of_mono0_charPol}), so $x=y$ and $\sfname{Away}(x,y)$
holds. Otherwise $y$ is not abelian, and \cref{pf:galred:guard} and $\sfname{R2}$(ii) give (dR), potential
semistability, $N_x=0$ and $P_x=P_y$ on $I_A$; by \cref{pf:galred:ine}, $\dim x=\dim y$ and $P_x=P_y$ in degree $0$,
which is $\sfname{Away}(x,y)$.
\end{proof}

The formal proof states the deduction for an arbitrary class of places in the tame clause; with the class of all
places it gives $\sfname{Row0}$. (The development also contains the variant for the class of $a$-large places, which
gives $\sfname{Row0b}$ directly; it is not used.)

\subsubsection{Compatibility at plain toral places}\label{sec:pf:galred:R1}

The base change arguments use a package of published results. $\sfname{Inp}$ is the conjunction of: Jacquet--Shalika
rigidity (strong multiplicity one for isobaric sums) \cite{JacquetShalika1981}; Gelfand--Kazhdan's description of the
contragredient \cite{GelfandKazhdan1975}; admissibility of irreducible smooth representations of $\GL_n$ and finite
length of standard modules \cite{BernsteinZelevinsky1977}; the Langlands classification and Zelevinsky's description
of the Grothendieck ring \cite{Zelevinsky1980}; the classification of generic unitary representations
\cite{Tadic1986}; the strong form of cyclic base change \cite{ArthurClozel1989}; the existence of $\ell$-adic
characters attached to algebraic Hecke characters, with their local properties \cite{Serre1968}; and the theorem that
de Rham representations are potentially semistable \cite{Berger2002}. All of these are proved in the formal development. $\sfname{Inq}$ is $\sfname{Inp}$ together with the statement that if
$M/K$ is a Galois extension of prime degree of totally complex fields and $P$ is a cuspidal lift of $\pi$ to $M$, then
$P$ is a \emph{character lift} of $\pi$ (\cref{pf:galred:lifts}); this is proved in the development
(\lref{UpperTheory.Gal.GamX.charLiftPrime_holds}), so $\sfname{Inq}$ holds.

\begin{definition}[Lifts]\label{pf:galred:lifts}
Let $M/K$ be an extension of totally complex number fields, $\pi$ cuspidal on $\GL_a(\A_K)$ and $P$ cuspidal on
$\GL_a(\A_M)$. $P$ is a \emph{lift} of $\pi$ if: for every place $u$ of $M$ above a place $v$ of $K$ with $q_u=q_v^f$, if
$\pi$ is unramified at $v$ with Satake parameter $B$, then $P$ is unramified at $u$ with Satake parameter $\{b^f:b\in
B\}$; $\galredinf_\tau(P)=\galredinf_{\tau|_K}(\pi)$ for all $\tau\colon M\to\C$; and whenever $K_v\cong M_u$ continuously
and compatibly with $K\to M$, the local components of $\pi$ at $v$ and of $P$ at $u$ correspond. $P$ is a
\emph{character lift} of $\pi$ if moreover, for all places $u\mid v$ and local components $\tau$ of $\pi$ at $v$ and $T$ of
$P$ at $u$: if $\supp(\tau)$ is the multiset of characters $\xi_1,\dots,\xi_a$, then $\supp(T)$ is the multiset of the
characters $\xi_j\circ\Nm_{M_u/K_v}$.
\formallabel{GlobalTheory.RoadU.CharLift}
\end{definition}

If $P$ is a lift of $\pi$ and $r$ is weakly attached to $\pi$, then $r|_{G_M}$ is weakly attached to $P$
(\lref{GlobalTheory.GalCut.weak_comp_res}), as Frobenius elements at $u$ are the $f$-th powers of Frobenius
elements at $v$.

\begin{lemma}\label{pf:galred:rcchars}
Let $x\in\galredIrr(F)$ with $\supp(x)=\{\xi_1,\dots,\xi_n\}$ an unlinked multiset of characters. Then
$N_{\galredrc(x)}=0$ and $P_{\galredrc(x)}(g)=\prod_j(X-\xi_j(a(g)))$ for all $g\in W_A$.
\formallabel{GalTw.GlobalTheory.RoadU.UnlinkedChars.rc_unlinked_chars}
\end{lemma}

\begin{proof}
By the Langlands classification $x=L(\mathfrak m)$ for a multiset $\mathfrak m$ of segments, and $\supp(x)$ is the
multiset of the terms of the segments of $\mathfrak m$ \cite{Zelevinsky1980}; by uniqueness of supports, every term is
some $\xi_j$. A segment of length $\ge2$ contains two terms $\rho,\rho\nu$, contradicting unlinkedness; so all
segments have length one and $\mathfrak m=\{\xi_1,\dots,\xi_n\}$. By (L3) and (L2),
$\galredrc(x)=\bigoplus_j\langle\xi_j\circ a\rangle$.
\end{proof}

The following lemma carries the base change argument. For a field $E\subset M\subset\Omega$, an element of $W_A$ fixing
$M$ is regarded as an element of $W_A(M)$.

\begin{lemma}[Spherical lifts]\label{pf:galred:sphlift}
Let $(\pi,\chi)$ be in scope $\mathrm{all}$ on $\GL_n$, $\sigma$ a local component of $\pi$ at $w$ with
$\supp(\sigma^\natural)=\{\xi_1,\dots,\xi_n\}$ an unlinked multiset of characters, $a$ a local Artin map at $A$, and
$g\in W_A$ with $\deg g\ne0$. Assume that $\pi$ has a weakly attached representation. There are a CM field $M$, $E\subset
M\subset\Omega$, fixed by $g$, a cuspidal lift $P$ of $\pi$ to $M$ with $(P,\chi')$ in scope $\mathrm{all}$ for some
$\chi'$, and $B$ such that $P$ is unramified at the place $u$ of $M$ below $A$ with Satake parameter $B$ and
$P_{\galredsph_A(B)}(g_0)=\prod_j(X-\xi_j(a(g_0)))$ for every $g_0\in W_A$ fixing $M$.
\formallabel{GalTw.GlobalTheory.RoadU.exists_sph_lift_of_toral}
\end{lemma}

\begin{proof}
Put $\xi_j^0=\xi_j|\cdot|^{(n-1)/2}$; these are the characters of the support of $\sigma$. By \cref{pf:galred:cmlift}
(which uses $\sfname{Inq}$) there are a CM field $M\subset\Omega$ containing $E$ and fixed by $g$, and a cuspidal character
lift $P$ of $\pi$ to $M$, in scope $\mathrm{all}$, such that $\xi^0_j(a(t))=1$ for all $j$ and every $t\in W_A$ of degree $0$
fixing $M$.

By local class field theory, the condition on the $\xi^0_j$ means that the characters $\xi^0_j\circ\Nm_{M_u/E_w}$ are
unramified (\lref{GlobalTheory.T.unramified_comp_norm_of_killed_inertia_fixing}). The local component of $P$ at $u$
has support $\{\xi^0_j\circ\Nm_{M_u/E_w}\}$, so $P$ is unramified at $u$, with Satake parameter $B$ given by the values of
these characters at a uniformizer, and $\galredsph_A(B)$ has the stated characteristic polynomials
(\lref{GalTw.GlobalTheory.T.frobenius_charpoly_eq_satake_of_place_above}).
\end{proof}

\begin{lemma}\label{pf:galred:toralgal}
Assume $\sfname{SLC}(\mathrm{all},\mathrm{unr})$. In the situation of \cref{pf:galred:sphlift}, let $r$ be weakly attached
to $\pi$, $w\nmid\ell$ and $x=\galredWD_\iota(r|_{W_A})$. Then $N_x=0$, $P_x(g)=\prod_j(X-\xi_j(a(g)))$ for every
$g\in W_A$, and $x=\galredrc(\sigma^\natural)$ for every local parameter map $\galredrc$.
\formallabel{GalTw.GlobalTheory.RoadU.toralGal_of_holds}
\end{lemma}

\begin{proof}
Let $g\in W_A$ with $\deg g\ne0$, and take $M,u,P,B$ as in \cref{pf:galred:sphlift}. Then $u\nmid\ell$, $r|_{G_M}$ is
weakly attached to $P$, and $\sfname{SLC}(\mathrm{all},\mathrm{unr})$ applied to $P$ at $u$ gives
$\galredFr_u(r|_{G_M})=B^{(n)}_u$, hence $\galredWD_\iota(r|_{W_A(M)})=\galredsph_A(B)$
(\lref{GlobalTheory.T.weilDeligneParam_eq_unramifiedSatake_of_frobenius}). The Weil--Deligne representation of
$r|_{G_M}$ at $A$ is the restriction of that of $r$ (\lref{GlobalTheory.AboveL.wdOf_restricts}), so $N_x=0$ and
$P_x(g)=P_{\galredsph_A(B)}(g)=\prod_j(X-\xi_j(a(g)))$. By \cref{pf:galred:degne0} this identity extends to all $g\in W_A$.
If $\galredrc$ is a local parameter map with Artin map $a'$, then by \cref{pf:galred:rcchars} and the identity just
proved for $a'$ in place of $a$, $\galredrc(\sigma^\natural)$ has no monodromy and the same characteristic polynomials as
$x$; so $x=\galredrc(\sigma^\natural)$ (\lref{GlobalTheory.WeakTheRep.eq_of_mono0_charPol}).
\end{proof}

\begin{theorem}\label{pf:galred:R1thm}
$\sfname{CoreAll}$ implies $\sfname{R1}$.
\formallabel{UpperTheory.Gal.ToralGalois.galoisCompat_unlinkedSupport_of_spherical}
\end{theorem}

\begin{proof}
Let $\sigma^\natural$ be plain toral, with support $\{\xi_1,\dots,\xi_n\}$
(\lref{GlobalTheory.RoadU.plainTor_exists_chars}). If $w\nmid\ell$, this is \cref{pf:galred:toralgal}. Let $w\mid\ell$.
The clause (dR) at $w$ is obtained from $\sfname{SLC}(\mathrm{all},\mathrm{dR})$ for the lift $P$ of
\cref{pf:galred:sphlift} (with $g$ of degree $1$): the de Rham property and the Hodge--Tate weights descend from the
open subgroup $D_A(M)$ to $D_A$, and the infinitesimal characters of $P$ are those of $\pi$
(\lref{GalTw.GlobalTheory.DeRhamOfPlainTorus.deRham_of_plainTor}). For the remaining assertions fix a frame at $A$
and $g$ of nonzero degree, and take $M,u,P,B$ from \cref{pf:galred:sphlift}. By $\sfname{SLC}(\mathrm{all},\mathrm{crys})$
and $\sfname{SLC}(\mathrm{all},\mathrm{phi})$ for $P$ at $u$, $r|_{D_A(M)}$ is crystalline, hence semistable without
monodromy, and its Weil--Deligne representation has the characteristic polynomials of $\galredsph_A(B)$. Potential
semistability and Fontaine's Weil--Deligne representation are compatible with restriction to the open subgroup $D_A(M)$,
so $r|_{D_A}$ is potentially semistable, $x=\galredWD(r|_{D_A})$ has $N_x=0$ and $P_x(g)=\prod_j(X-\xi_j(a(g)))$
(\lref{GalTw.GlobalTheory.RoadU.toralAbove_of_holds}). By \cref{pf:galred:degne0,pf:galred:rcchars},
$P_x(g)=P_{\galredrc(\sigma^\natural)}(g)$ for all $g\in W_A$ and all local parameter maps $\galredrc$.
\end{proof}

\subsubsection{Compatibility on inertia at depth-zero places}\label{sec:pf:galred:R2}

The plain toral case of $\sfname{R2}$ is contained in $\sfname{R1}$. The other case is reduced to it by an unramified
base change which turns the depth-zero supercuspidal members of the support into characters. A representation $x$ of
$\GL_n(E_w)$ is \emph{mixed} if it has depth zero and $\supp(x)$ is unlinked and apart. The reduction uses the
following statements; (ii) and (v) are proved in \cref{sec:pf:galred:aux}, the others in the development.

\begin{enumerate}[label=(\roman*),nosep]
\item (\lref{UpperTheory.Gal.ToralGalois.R2.mixedOf}) If $\sigma$ is a local component of a cuspidal $\pi$ with
  nonzero $K_1$-fixed vectors, $\sigma^\natural$ is apart and $N_{\galredrc(\sigma^\natural)}=0$ for some local parameter map,
  then $\sigma$ is mixed.
\item (\cref{pf:galred:unrfield}) For a CM field $E$, a split place $w$,
  $f\ge1$ and a finite set $S$ of places, there is an abelian CM extension $X/E$, built as a tower of steps adapted to
  $S$, in which the inertia groups at $w$ and at the non-split places act trivially and every place of $X$ above $w$
  has norm $q_w^f$.
\item (\lref{UpperTheory.Gal.NearKilling.allLift_holds}) If $(\pi,\chi)$ is in scope $\mathrm{all}$, $\pi$ is unramified
  outside $S$ and $X$ is as in (ii), then $\pi$ has a cuspidal lift $P$ to $X$ through a chain of Galois extensions of
  prime degree, $(P,\chi')$ is in scope $\mathrm{all}$ for some $\chi'$, and $r|_{G_X}$ is weakly attached to $P$
  whenever $r$ is weakly attached to $\pi$.
\item (\lref{UpperTheory.Gal.NearKilling.ineNorm}) If the inertia group at $A$ fixes $X$ and $u$ is the place of $X$
  below $A$, then $I_A(X)=I_A$, and local Artin maps of $X_u$ and $E_w$ are compatible with the norm on inertia.
\item (\cref{pf:galred:suppbc}) \emph{Supports under unramified base change.} If $\tau,\tau'$ are local components of $\pi$
  at a $\Q_g$-place $v$ whose inertia groups fix $X$ and of $P$ at the place $u$ of $X$ below $A$, $\tau$ has depth zero
  and $q_u=g^f$ with $f$ divisible by the dimensions of the members of $\supp(\tau)$, then $\supp(\tau')$ is a multiset of
  tame characters $\xi_1,\dots,\xi_n$ of $X_u^\times$ with $P_{\galredrc(\tau)}(t)=\prod_j(X-\xi_j(a'(t)))$ for $t\in I_A$,
  for all local parameter maps $\galredrc$ of $E_v$ and local Artin maps $a'$ of $X_u$, and it is unlinked if $\tau$ is
  mixed.
\item (\lref{UpperTheory.Gal.ToralGalois.R2.recChars}) In the situation of (v), if $\{\xi_j\}$ is unlinked then
  $\tau'^\natural$ is plain toral; and $P_{\galredrc'(\tau'^\natural)}(t)=\prod_j(X-\xi_j(a'(t)))$ for $t\in I_A(X)$, every
  local parameter map $\galredrc'$ of $X_u$ and every local Artin map $a'$ (compare \cref{pf:galred:rcchars}).
\end{enumerate}

\begin{theorem}\label{pf:galred:R2thm}
$\sfname{R1}$ implies $\sfname{R2}$.
\formallabel{UpperTheory.Gal.ToralGalois.R2.localGlobalCompatAtIwahoriFixed_of_unlinked}
\end{theorem}

\begin{proof}
Let the data of $\sfname{R2}$ be given. If $\sigma^\natural$ is plain toral, $\sfname{R1}$ gives more than required.
Otherwise $w$ is a $\Q_g$-place and $\sigma^\natural$ is apart, so $\sigma$ is mixed by (i). A $\Q_g$-place of a CM field
is split over $E^+$: it is unramified over $g$ by \cref{pf:galred:gnotsq} and has residue degree $1$ over $g$, so it is
neither ramified nor inert over $E^+$. Let $f=n!$ and let $S$ be the
finite set of places where $\pi$ is ramified. By (ii)--(iv) there are an abelian CM extension $X/E$ whose inertia
groups at $w$ fix $X$ and whose places above $w$ have norm $q_w^f$, and a cuspidal lift $P$ of $\pi$ to $X$ in scope
$\mathrm{all}$ with $r|_{G_X}$ weakly attached to $P$. Let $u$ be the place of $X$ below $A$ and $\tau'$ a local component
of $P$ at $u$. The members of a support of a representation of $\GL_n$ have dimension at most $n$, hence dividing $n!$;
so by (v) and (vi), $\tau'^\natural$ is plain toral and $P_{\galredrc'(\tau'^\natural)}(t)=P_{\galredrc(\sigma^\natural)}(t)$ for all
$t\in I_A$, for every local parameter map $\galredrc'$ of $X_u$, which exists by \cref{thm:rectexists}. (Here
$\galredrc(\sigma^\natural)=\galredrc(\sigma)(\frac{1-n}2)$ by (L5) has the same characteristic polynomials on inertia as
$\galredrc(\sigma)$.)

Apply $\sfname{R1}$ to $P$, $r|_{G_X}$ and $\tau'$ at $u$. If $w\nmid\ell$, it gives $N=0$ and
$\galredWD_\iota(r|_{W_A(X)})=\galredrc'(\tau'^\natural)$; since $I_A(X)=I_A$ and the Weil--Deligne representation of
$r|_{G_X}$ is the restriction of that of $r$ (\lref{GlobalTheory.AboveL.wdOf_restricts}), $x=\galredWD_\iota(r|_{W_A})$
has $N_x=0$ and $P_x(t)=P_{\galredrc(\sigma^\natural)}(t)$ for $t\in I_A$. If $w\mid\ell$, $\sfname{R1}$(ii) at $u$ gives, for
every frame, potential semistability of $r|_{D_A(X)}$, absence of monodromy and the characteristic polynomials on
inertia, and these descend to $r|_{D_A}$ as in the proof of \cref{pf:galred:R1thm}; the clause (dR) at $w$ descends from
the clauses (dR) at the places of such extensions above $w$, whose lifts have the same infinitesimal characters
(\lref{UpperTheory.Gal.ToralGalois.R2.drDesc}).
\end{proof}

\begin{proof}[Proof of \cref{pf:galred:rest}]
By \cref{pf:galred:R1thm,pf:galred:R2thm}, $\sfname{CoreAll}$ implies $\sfname{R1}$ and $\sfname{R2}$, and
\cref{pf:galred:deduce} gives $\sfname{Row0}(\mathrm{recT})$.
\end{proof}

% ---------------------------------------------------------------------------
\subsection{Definite unitary groups of odd rank}\label{sec:pf:galred:odd}

\subsubsection{Notation and the three statements}\label{sec:pf:galred:oddstmts}

Let $E$ be a CM field. A place $w$ of $E$ is \emph{split} if another place of $E$ lies over the same place of $E^+$;
two places are \emph{apart} if they lie over different places of $E^+$; $w$ is apart from a finite set $S$ if it
is apart from every member of $S$. For a split place $w$ there is $c(w)\ne w$ over the same place of $E^+$.
$\galredU_a(\A_{E^+})=\{g\in\GL_a(\A_E):{}^t\bar g\,g=1\}$, and an \emph{automorphic representation} $U$ of $\galredU_a$
is a nonzero irreducible right-translation-stable space of automorphic forms on
$\galredU_a(E^+)\backslash\galredU_a(\A_{E^+})$ (\cref{sec:galois:arch}). At a split place $w$ the local group is
$\GL_a(E_w)$; a \emph{local component} of $U$ at $w$ is a smooth irreducible $\sigma$ with a nonzero equivariant map to
$U$, and $U$ is \emph{spherical at $w$ with Satake parameter $B$} if a local component at $w$ is spherical with Satake
parameter $B$; we write $\galredSat_w(U)=B$ (it is unique), and $B^{(a)}_w$ as before. At a non-split place, $U$ is
\emph{unramified} if it has a nonzero vector fixed by the hyperspecial subgroup. For $\tau\colon E\to\C$, a dominant
weight $\lambda\in\Z^a$ ($\lambda_0\ge\dots\ge\lambda_{a-1}$) is a \emph{weight of $U$ at $\tau$} if the irreducible algebraic
representation of highest weight $\lambda$ of the factor at $\tau$ occurs in $U$; its \emph{Hodge--Tate multiset} is
$J_\lambda=\{k-\lambda_k:0\le k<a\}=\{\frac{a-1}2-z:z\in\galredinf(\lambda)\}$, where
$\galredinf(\lambda)=\{\lambda_k+\frac{a-1}2-k\}$. An embedding $\tau\colon E\to\C$ \emph{lies over} a place $v$ above
$\ell$ if $\iota^{-1}\tau$ induces $v$.

\begin{definition}[Conditions on fields, places and representations]\label{pf:galred:conds}
Let $\ell$ be a prime, $d>0$ an integer and $a\ge1$.
\begin{enumerate}[label=(\alph*),nosep]
\item $\sfname{FieldOK}(E,\ell,d)$: $E$ contains a square root $x$ of $-d$; $-d$ is a square in $\Q_\ell$; $[E^+:\Q]$ is
  even; every non-split place of $E$ is prime to $\ell$ and unramified over $E^+$. Then all places above $\ell$ are
  split.
\item A place $w$ is \emph{$d$-usable} if it is split and $-d$ is a square in $\Q_q$, $q$ the prime below $w$.
\item $U$ is \emph{unramified off $T$}, $T$ a finite set of non-split places, if it is unramified at all non-split
  places outside $T$.
\item $U$ has a \emph{type of $\ell$-power order} at $w_0$ if there are an open compact torsion-free subgroup $J$ of
  $\GL_a(E_{w_0})$, a surjection $\vartheta\colon J\to\Z/\ell^m$ and a nontrivial character $\chi$ of $\Z/\ell^m$ such that
  every smooth irreducible representation of $\GL_a(E_{w_0})$ with a nonzero $\chi\circ\vartheta$-eigenvector is
  supercuspidal, and $U$ contains a nonzero $\chi\circ\vartheta$-eigenvector for $J$.
\item A multiset $J$ of $a$ integers, sorted as $j_0\le\dots\le j_{a-1}$, is \emph{gapped} if $j_{i+1}-j_0>a(j_i-j_0)$ for
  all $i+1<a$. $U'$ has \emph{lacunary weight at $v'$ relative to $U$} if $U$ and $U'$ have the same weights at every
  $\tau$ such that every place over which $\tau$ lies is apart from $w_0$ and $v'$, and there is $\tau_0$ lying over $v'$
  such that $J_\lambda$ is gapped for every weight $\lambda$ of $U'$ at $\tau_0$.
\item A multiset $B=\{\alpha_1,\dots,\alpha_a\}$ of nonzero complex numbers is \emph{generic} if
  $D_a(\alpha_1,\dots,\alpha_a)\ne0$, where
  \begin{multline*}
    D_a=\prod_{i\ne j}\prod_{c=1}^{a}(X_i^c-X_j^c)\cdot\prod_{j\notin\{i,k\}}(X_iX_k-X_j^2)\\
    \cdot\prod_{m}\Bigl(\prod_iX_i^{m_i}-\prod_iX_i\Bigr)\in\Z[X_1,\dots,X_a],
  \end{multline*}
  the last product over the $m\in\{0,\dots,a\}^a$ with $\sum m_i=a$ and $m\ne(1,\dots,1)$. $U$ is \emph{generic at $w$} if
  every Satake parameter of $U$ at $w$ is generic.
\item For a representation $r\colon G_E\to\GL_a(\Qlbar)$ and $w\mid\ell$: $\sfname{CrysDR}(U,r,w)$ means that at $w$,
  $r|_{D_A}$ is crystalline and de Rham, and its $\tau$-Hodge--Tate multiset is $J_\lambda$ for every $\tau$ induced by the
  frame and every weight $\lambda$ of $U$ at $\iota\tau$; $\sfname{Phi}(r,w,B)$ means that at $w$,
  $P_{\galredWD(r|_{D_A})}(g)=P_{\galredsph_A(B)}(g)$ for all $g\in W_A$ (formally: the power sums of the roots agree for
  every exponent $n\ge0$, which is equivalent).
\end{enumerate}
\end{definition}

The three statements about odd rank are the following. In all of them $E$ is a CM field, $\Omega,\ell,\iota$ are as
before, $a$ is odd, $\sfname{FieldOK}(E,\ell,d)$ holds, $T$ is a finite set of non-split places, and $w_0,v'$ are places
above $\ell$ which are apart.

\begin{definition}[$\sfname{SeedFinOdd}$]\label{pf:galred:seeddef}
For all such data, every automorphic $U_0$ of $\galredU_a$ and all finite sets $S,Y$ of places, there is a finite set
$Z_0$ of $d$-usable places prime to $\ell$ and apart from $w_0$ and $S$, such that every \emph{seed} $U$, i.e.\ every
automorphic representation of $\galredU_a$ which is unramified off $T$, has a type of $\ell$-power order at $w_0$, has
lacunary weight at $v'$ relative to $U_0$, is spherical at every split place apart from $w_0$ and $S$ and at $v'$,
and is generic at the places of $Z_0$, has a continuous semisimple $r\colon G_E\to\GL_a(\Qlbar)$ such that:
\begin{enumerate}[label=(S\arabic*),nosep]
\item there are an open dense $V\subset G_E$ and a finite set $X$ of places such that for every $d$-usable $w\notin X$
  prime to $\ell$ and apart from $S$, every $B=\galredSat_w(U)$ and every $A$ above $w$: $r(I_A)=1$, and
  $\iota\det(X-r(s^{-1}))=\prod_{B^{(a)}_w}(X-b)$ for every arithmetic Frobenius $s\in V$ at $A$;
\item for $w\in Y$ $d$-usable, prime to $\ell$, apart from $S$, with $B=\galredSat_w(U)$ generic: $\galredFr_w(r)=B^{(a)}_w$;
\item for $w\in Y$ above $\ell$, apart from $w_0$, $v'$ and $S$: $\sfname{CrysDR}(U,r,w)$;
\item for $w\in Y$ above $\ell$, apart from $w_0$, $v'$ and $S$, with $B=\galredSat_w(U)$ generic: $\sfname{Phi}(r,w,B)$.
\end{enumerate}
\formallabel{GlobalTheory.CongL.SeedFinOdd}
\end{definition}

\begin{definition}[$\sfname{CongFinOdd}$]\label{pf:galred:congdef}
A \emph{family of local Hecke operators} of level $(K_w)_w$ ($K_w\subset\GL_a(E_w)$ for split $w$) indexed by a set $J$
consists of split places $w_j\nmid\ell$ and finitely supported functions $f_j$ on $\GL_a(E_{w_j})$ with $\ell$-adically
integral values (after $\iota^{-1}$), such that $f_j$ acts on the $K_{w_j}$-fixed vectors of every smooth irreducible
$\sigma$ by a scalar $z_j(\sigma)$. The \emph{eigenvalue system} of $U$ is $\theta_U=(\iota^{-1}z_j(\sigma_j))_j$, $\sigma_j$ the
local component of $U$ at $w_j$. For $\theta\colon J\to\Qlbar$ and a set $\mathcal S$ of such functions, $\theta$ is
\emph{polynomially approximable} by $\mathcal S$ if for every $\varepsilon>0$ there are finitely many
$\eta\in\mathcal S$ and $c_\eta\in\Qlbar$ with $|P(\theta)-\sum c_\eta P(\eta)|\le\varepsilon$ for \emph{all}
$P\in\Z[X_j:j\in J]$.

$\sfname{CongFinOdd}$ states: for all data as above and every $U$ unramified off $T$, there is a level $(K_w)$, with
$K_w$ open compact for split $w$ apart from $w_0$ and $K_w=\GL_a(\OO_{E_w})$ at split places apart from $w_0,v'$ where
$U$ is spherical, such that the local components of $U$ at split places apart from $w_0,v'$ have nonzero
$K_w$-fixed vectors, and for all finite sets $Z$ (of split places apart from $w_0,v'$ where $U$ is spherical) and $W$
(of such places above $\ell$) there are nonzero constants $c(w,n)$, $w\in W$, $n\in\Z$, such that for every family of
local Hecke operators of level $(K_w)$ the extended system
$\bigl(\theta_U,\ (\iota^{-1}(c(w,n)p_n(U,w)))_{w\in W,n\in\Z}\bigr)$, $p_n(U,w)=\sum_{b\in B^{(a)}_w}b^n$, is
polynomially approximable by the extended systems of seeds $U'$ (relative to $U_0=U$) of the same level, with the same
weights as $U$ at the embeddings over places apart from $w_0,v'$, generic at the places of $Z\cup W$; and if the
operators have rational values, then $\theta_U$ takes values in a finite extension of $\Ql$.
\formallabel{GlobalTheory.CongL.CongFinOdd}
\end{definition}

\begin{definition}[$\sfname{LimFinOdd}$]\label{pf:galred:limdef}
For all data as above and every $U$ unramified off $T$ there is a continuous semisimple $r\colon
G_E\to\GL_a(\Qlbar)$ such that:
\begin{enumerate}[label=(L\arabic*),nosep]
\item $\galredFr_w(r)=\galredSat_w(U)^{(a)}_w$ for every $d$-usable $w\nmid\ell$ at which $U$ is spherical;
\item $\sfname{CrysDR}(U,r,w)$ for every $w\mid\ell$ apart from $w_0,v'$ at which $U$ is spherical;
\item for every such $w$ with $B=\galredSat_w(U)$, every frame at a place $A$ above $w$ and every $\sigma\in W_A$, there are
  a set $\mathcal C$ of continuous representations, crystalline (for the frame) with a common Hodge--Tate multiset
  $J_\tau$ for each $\tau$, which \emph{jointly approximates} $r$ and the power sums of $\sigma$ on $\galredsph_A(B)$: for
  every $n$ and $\varepsilon>0$ there are finitely many $\rho\in\mathcal C$ and $c_\rho$ with $|\galredtr r(g)-\sum
  c_\rho\galredtr\rho(g)|\le\varepsilon$ for all $g\in G_E$ and $|\iota^{-1}p_j(\sigma,\galredsph_A(B))-\sum
  c_\rho\iota^{-1}p_j(\sigma,\galredWD(\rho|_{D_A}))|\le\varepsilon$ for $j\le n$, where $p_j(\sigma,x)$ is the $j$-th power
  sum of the roots of $P_x(\sigma)$.
\end{enumerate}
\formallabel{GlobalTheory.CongL.LimFinOdd}
\end{definition}

\begin{theorem}\label{pf:galred:seedcong}
$\sfname{GeoOne}$ implies $\sfname{SeedFinOdd}$, and $\sfname{CongFinOdd}$ holds.
\formallabel{UpperTheory.Gal.Int.Div.seedFinOdd_of_geoOne}
\formallabel{UpperTheory.Gal.Int.Cong.congFinOdd_of_ex}
\end{theorem}

The two assertions are proved in \cref{sec:pf:galred:seeds,sec:pf:galred:cong}; the limit statement in
\cref{sec:pf:galred:lim}.

\subsubsection{Seeds from geometry}\label{sec:pf:galred:seeds}

$\sfname{GeoOne}$ (\cref{def:galois:geoone}) provides representations $\rho$ of unknown dimension whose Frobenius
elements are \emph{annihilated} by the Hecke polynomials. We call a continuous semisimple $\rho\colon
G_E\to\GL_{d'}(\Qlbar)$ \emph{basic} for $U$ (and $d$) if (1) for all but finitely many $d$-usable $w\nmid\ell$ where $U$ is
spherical with $B=\galredSat_w(U)$, $\rho$ is unramified at $w$ and $\prod_{B^{(a)}_w}(X-\iota^{-1}b)$ annihilates
$\rho(s^{-1})$ for every arithmetic Frobenius $s$ at every $A$ above $w$ (``$\rho$ is killed by $B$ at $w$''); and (2) at some place
above $\ell$ and some frame, $\rho$ is de Rham and some $\tau$-Hodge--Tate multiset of $\rho$ contains a \emph{lacunary}
set of $a$ integers, i.e.\ a set of $a$ integers no two disjoint nonempty subsets of which have the same average. This
is the notion of \cref{def:galois:geoone}.

The extraction of $a$-dimensional representations rests on an abstract division theorem. For representations
$r\colon\Gamma\to\GL_a(k)$ and $\rho\colon\Gamma\to\GL_N(k)$ of a group, $r$ is \emph{inside} $\rho$ if there is an
$N\times a$ matrix $M$ of rank $a$ with $\rho(g)M=Mr(g)$ for all $g$, and $r$ \emph{covers} $\rho$ if there are $m$ and an
$(am)\times N$ matrix $M'$ of rank $N$ with $(r(g)\otimes1_m)M'=M'\rho(g)$ for all $g$.

\begin{theorem}[Division]\label{pf:galred:division}
Let $\Gamma$ be a compact topological group, $K$ an algebraically closed topological field of characteristic $0$, $K'$
an algebraically closed extension of $K$, and $\rho\colon\Gamma\to\GL_N(K)$ continuous and semisimple. Let
$S\subset\Gamma$ be dense, and for $s\in S$ let $P_s\in K[X]$ be monic of degree $n$ with $P_s(\rho(s))=0$. Let $Z\in
M_N(K')$ lie in the Lie algebra of the Zariski closure of $\rho(\Gamma)$ (every polynomial vanishing on $\rho(\Gamma)$ has
differential at $1$ vanishing on $Z$), and assume that $\det(X-Z)$ has $n$ distinct roots $e_1,\dots,e_n$ which are
lacunary ($|B|\sum_{i\in A}e_i\ne|A|\sum_{i\in B}e_i$ for disjoint nonempty $A,B$). Let $\psi\colon\Gamma\to K^\times$ be a
continuous character with $\psi(s)=(-1)^nP_s(0)$ for $s\in S$. Then there is a continuous semisimple $r\colon
\Gamma\to\GL_n(K)$ which is inside $\rho$ and covers $\rho$, with $\det r=\psi$, $P_s(r(s))=0$ for $s\in S$, whose
characteristic polynomial annihilates $\rho(g)$ for all $g$, such that the set of $g$ with separable characteristic
polynomial $\det(X-r(g))$ is open and dense, and $\det(X-r(s))=P_s$ for every $s\in S$ for which $\det(X-r(s))$ is
separable or $P_s$ is product-generic (separable, and no multiset of $n$ roots of $P_s$ other than the multiset of all roots
has the same product).
\formallabel{LateralTheory.DivDb.goal_holds}
\end{theorem}

The proof is given in \cref{sec:pf:galred:divproof}. The element $Z$ is provided by Sen's theory \cite{Sen1980}, in which the
Sen operator of $\rho$ at a place above $\ell$ lies in the Lie algebra of the Zariski closure of the image and has the
Hodge--Tate weights as eigenvalues: if $\rho\colon G_E\to\GL_N(\Qlbar)$ is continuous and some $\tau$-Hodge--Tate multiset
of $\rho$ at a place above $\ell$ contains $n$ lacunary integers, then there is $Z$ as in \cref{pf:galred:division} with
$n$ distinct lacunary eigenvalues
(\lref{UpperTheory.Gal.Int.Div.senWLac_holds} and \lref{UpperTheory.Gal.Int.Div.lacBridge_of}).

\begin{lemma}\label{pf:galred:divset}
Let $\rho\colon G_E\to\GL_{d'}(\Qlbar)$ be continuous and semisimple, with a $\tau$-Hodge--Tate multiset containing $a$
lacunary integers. Let $\psi\colon G_E\to\Qlbar^\times$ be a continuous character such that the set of $g$ for which some
monic $P\in\Qlbar[X]$ of degree $a$ satisfies $P(\rho(g))=0$ and $\psi(g)=(-1)^aP(0)$ is dense. Then there is a continuous
semisimple $r\colon G_E\to\GL_a(\Qlbar)$ inside $\rho$ and covering $\rho$, such that the set of $g$ with $\det(X-r(g))$
separable is open and dense.
\formallabel{UpperTheory.Gal.Int.Div.divSet_holds}
\end{lemma}

\begin{proof}
Choose such a $P=P_g$ for each $g$ in the dense set, and apply \cref{pf:galred:division} with $\Gamma=G_E$ (compact),
$K=\Qlbar$ and the element $Z$ given by Sen's theory.
\end{proof}

A continuous character $\psi$ of $G_E$ is a \emph{determinant character} for $U$ if at every split place $w$ where $U$ is
spherical with Satake parameter $B$, the four clauses hold for $\psi$ (as a one-dimensional representation) with
Frobenius multiset $\{\prod_{b\in B^{(a)}_w}b\}$ and some Hodge--Tate multisets. Such a $\psi$ exists for every $U$
(\lref{UpperTheory.Gal.Int.Div.detEx_holds}); it is built from an algebraic Hecke character $\chi$ of $E$ with
$\chi(\varpi_w)=\prod_{b\in B}b$ at the split places $w$ where $U$ is spherical with Satake parameter $B$, which exists by
\lref{UpperTheory.Gal.Int.Div.centU_holds} (central characters). A continuous semisimple $r\colon G_E\to\GL_a(\Qlbar)$ is a \emph{division
representation} for $U$ if it satisfies (S1) of \cref{pf:galred:seeddef} with $S=\emptyset$.

\begin{lemma}\label{pf:galred:divone}
Let $d$ satisfy $\sqrt{-d}\in E$, let $U$ be automorphic, $\psi$ a determinant character for $U$, and $\rho$ basic for $U$. Then
there is a division representation $r$ for $U$ which is inside $\rho$ and covers $\rho$.
\formallabel{UpperTheory.Gal.Int.Div.divOne_holds}
\end{lemma}

\begin{proof}
Let $X_\rho$ be the finite exceptional set of (1). By Chebotarev's density theorem \cite{Chebotarev1926}, the set of
inverses of arithmetic Frobenius elements at $d$-usable places $w\notin X_\rho$, $w\nmid\ell$, where $U$ is spherical, is
dense in $G_E$ (\lref{GlobalTheory.CongL.UnitLim.dense_frob_usable_off}). For such an element $g=s^{-1}$ take
$P=\prod_{B^{(a)}_w}(X-\iota^{-1}b)$: it is monic of degree $a$, annihilates $\rho(g)$ because $\rho$ is basic, and
$(-1)^aP(0)=\iota^{-1}\prod_{B^{(a)}_w}b=\psi(g)$ because $\psi$ is a determinant character. \Cref{pf:galred:divset} gives
$r$ inside and covering $\rho$, with separable characteristic polynomial on an open dense set $V'$. Let $V$ be the set of
$s$ with $s^{-1}\in V'$. Let $M$ be the matrix of rank $a$ with $\rho(g)M=Mr(g)$. If $\rho(g)=1$ then $M=Mr(g)$, so
$r(g)=1$; and if $P(\rho(g))=0$ then $MP(r(g))=P(\rho(g))M=0$, so $P(r(g))=0$. Hence for $w\notin X_\rho$ as in (S1), $A$
above $w$ and an arithmetic Frobenius $s\in V$: $r(I_A)=1$, and the characteristic polynomial of $r(s^{-1})$ is separable, so
it equals the minimal polynomial, which divides $P$; both being monic of degree $a$, it equals $P$.
\end{proof}

A finite multiset $s$ of nonzero complex numbers is \emph{multiplicatively generic} if every multiset of $|s|$ elements of
$s$ (counted with multiplicity, i.e.\ every multiset with the same cardinality, all of whose members belong to $s$) having
the same product as $s$ is equal to $s$.

\begin{lemma}\label{pf:galred:frobcond}
Let $\psi$ be a determinant character for $U$, $r$ inside $\rho$ with $\det r=\psi$, $w\nmid\ell$ split with
$B=\galredSat_w(U)$ multiplicatively generic, and assume $\rho$ is unramified at $w$ and killed by $B$ at $w$. Then
$\galredFr_w(r)=B^{(a)}_w$.
\formallabel{UpperTheory.Gal.Int.Div.frobSeed_of}
\end{lemma}

\begin{proof}
$B^{(a)}_w$ is multiplicatively generic as well, since it is obtained from $B$ by multiplication by the constant
$q_w^{(a-1)/2}$, which multiplies the products of all multisets of cardinality $a$ by the same factor. $r$ is trivial on $I_A$
since $\rho$ is. For an arithmetic Frobenius $s$, the eigenvalues of $r(s^{-1})$ are roots of
$\prod_{B^{(a)}_w}(X-\iota^{-1}b)$ (which annihilates $\rho(s^{-1})$, hence $r(s^{-1})$), and their product is $\det r(s^{-1})=
\psi(s^{-1})=\iota^{-1}\prod_{B^{(a)}_w}b$. By multiplicative genericity, the multiset of eigenvalues is $\iota^{-1}B^{(a)}_w$.
\end{proof}

If $B=\{\alpha_1,\dots,\alpha_a\}$ is generic in the sense of \cref{pf:galred:conds}(f), then $B$ is multiplicatively generic.
Indeed the $\alpha_i$ are distinct (first factor of $D_a$), so a multiset of $a$ elements of $B$ is described by multiplicities
$m_i\in\{0,\dots,a\}$ with $\sum m_i=a$; equality of products means $\prod\alpha_i^{m_i}=\prod\alpha_i$, and the last factor of
$D_a$ forces $m=(1,\dots,1)$.

\begin{proposition}\label{pf:galred:divall}
Let $\sqrt{-d}\in E$ and let $U$ be automorphic with a basic representation. There is a division representation $r$ for
$U$ such that for every basic $\rho$ and every place $w$:
\begin{enumerate}[label=(\roman*),nosep]
\item if $w$ is split, $w\nmid\ell$, $B=\galredSat_w(U)$ is generic and $\rho$ is unramified at $w$ and killed by $B$ at $w$,
  then $\galredFr_w(r)=B^{(a)}_w$;
\item if at $w$ the representation $\rho$ is crystalline and its $\tau$-Hodge--Tate multisets contain $J_\lambda$ for every
  weight $\lambda$ of $U$ at $\iota\tau$ ($\tau$ induced by the frame), then $\sfname{CrysDR}(U,r,w)$;
\item if $w$ is split, $B=\galredSat_w(U)$ is generic, and at $w$ the representation $\rho$ is crystalline and, for some
  element $g$ of degree $-1$, $P_{\galredWD(\rho|_{D_A})}(g)$ divides a power of $P_{\galredsph_A(B)}(g)$, then $r$ is crystalline at $w$
  and $\sfname{Phi}(r,w,B)$.
\end{enumerate}
\formallabel{UpperTheory.Gal.Int.Div.divAll_holds}
\end{proposition}

\begin{proof}
Take a determinant character $\psi$ (\lref{UpperTheory.Gal.Int.Div.detEx_holds}) and, by \cref{pf:galred:divone}, a
division representation $r$ inside and covering some basic representation. Then $\det r=\psi$: at the arithmetic Frobenius
elements $s\in V$ of (S1), $\det r(s^{-1})$ and $\psi(s^{-1})$ both equal $\iota^{-1}\prod_{B^{(a)}_w}b$; the inverses of these
elements are dense (Chebotarev), and both sides are continuous. Let $\rho$ be basic and let $r'$ be the division
representation inside and covering $\rho$ given by \cref{pf:galred:divone}. The characteristic polynomials of $r$ and $r'$
agree at the inverses of the Frobenius elements lying in both open dense sets $V$, outside both exceptional sets; this set
is dense, so $\galredtr r=\galredtr r'$ by continuity, and $r\cong r'$ since semisimple representations in characteristic $0$
are determined by their traces \cite[(30.16)]{CurtisReiner1962}. Being inside $\rho$ and covering $\rho$ are invariant under
isomorphism, so $r$ is inside and covers $\rho$. Now (i) is \cref{pf:galred:frobcond}, as generic Satake parameters are
multiplicatively generic.

For (ii) and (iii) we use that if $r$ is inside a semisimple $\rho$, then at a frame $r|_{D_A}$ is crystalline whenever
$\rho|_{D_A}$ is, its Hodge--Tate multisets are contained in those of $\rho$, and the characteristic polynomials of elements of
$W_A$ on $\galredWD(r|_{D_A})$ divide those on $\galredWD(\rho|_{D_A})$ (\lref{UpperTheory.Gal.Int.Div.subLoc_holds}).
(ii): $r$ is crystalline at $w$. Since $r$ covers $\rho$ and is semisimple, $\rho$ is inside the sum of $m$ copies of $r$, so
the $\tau$-Hodge--Tate multiset of $\rho$ is contained in $m$ copies of the $\tau$-Hodge--Tate multiset $J_r$ of $r$. By
hypothesis it contains $J_\lambda$, which consists of $a$ distinct integers; so $J_\lambda\subset J_r$, and as $J_r$ has $a$
elements, $J_r=J_\lambda$. (iii): $r$ is crystalline at $w$. For the given $g$
of degree $-1$, $P_{\galredWD(r|_{D_A})}(g)$ divides $P_{\galredWD(\rho|_{D_A})}(g)$, hence a power of
$P_{\galredsph_A(B)}(g)=\prod_{B^{(a)}_w}(X-b)$; the product of its roots is the eigenvalue of $g$ on the Weil--Deligne
representation of $\det r=\psi$, which is $\prod_{B^{(a)}_w}b$ by the clauses for $\psi$. By multiplicative genericity,
$P_{\galredWD(r|_{D_A})}(g)=P_{\galredsph_A(B)}(g)$. As $r|_{D_A}$ is crystalline, $\galredWD(r|_{D_A})$ is unramified without
monodromy, hence determined by the characteristic polynomial of one element of degree $-1$, and $\sfname{Phi}(r,w,B)$ follows
(\lref{UpperTheory.Gal.Int.Div.phiSeed_holds}).
\end{proof}

\begin{proof}[Proof that $\sfname{GeoOne}$ implies $\sfname{SeedFinOdd}$]
Take $Z_0=\emptyset$, and let $U$ be a seed. Write $x^2=-d$ in $E$. Since $-d$ is a square in $\Q_\ell$, $\ell$ splits in
$\Q(x)$; complex conjugation acts on $\Q(x)$ nontrivially, so if $u$ is a split place above $\ell$ and $w$ a place above
$\ell$, then $u$ or $c(u)$ lies over the same place of $\Q(x)$ as $w$; and $c(u)$ is apart from exactly the same places as $u$.
A gapped multiset $J$ is lacunary: shifting by $-j_0$ we may assume $0=j_0\le j_1\le\dots$ with $j_{i+1}>aj_i$; if
$|B|\sum_Aj=|A|\sum_Bj$ for disjoint nonempty $A,B$, let $j_m$ be the largest index used, say in $A$ (then $m\ge1$); then
$|B|\sum_Aj\ge|B|j_m>|B|\,a\,j_{m-1}\ge|A|\,|B|\,j_{m-1}\ge|A|\sum_Bj$, a contradiction. Hence $U$ has lacunary weight at $v'$
in the sense of \cref{def:galois:geoone}. The hypotheses of $\sfname{GeoOne}$ on $U$ at $w_0$ (type of $\ell$-power order) and
at $v'$ (spherical, lacunary weight) are invariant under replacing $w_0$ by $c(w_0)$ and $v'$ by $c(v')$, by conjugation of
local components (\lref{UpperTheory.Gal.Int.Div.locIso_of}). Hence for every place $w$ above $\ell$ there are $u_0\in\{w_0,c(w_0)\}$ and
$u'\in\{v',c(v')\}$, both lying over the same place of $\Q(x)$ as $w$, for which $\sfname{GeoOne}$ applies, and its conclusions
(G1)--(G3) hold for $(u_0,u')$.

Apply \cref{pf:galred:divall} to $U$ and the basic representation of (G1) for the pair attached to $w=v'$; let $r$ be the
resulting division representation. It satisfies (S1). For (S2), let $w\in Y$ be $d$-usable, prime to $\ell$, apart from $S$,
with $B$ generic; $U$ is spherical at $w$ since $U$ is a seed, so (G2) gives a basic $\rho$ killed by $B$ at $w$, and
\cref{pf:galred:divall}(i) gives $\galredFr_w(r)=B^{(a)}_w$. For (S3) and (S4), let $w\in Y$ lie above $\ell$, apart from
$w_0,v',S$; then $U$ is spherical at $w$, and with $(u_0,u')$ chosen for $w$, $w$ lies over the same place of $\Q(x)$ as
$u'$ and is apart from $u_0$ and $u'$, so (G3) provides the basic representations required in
\cref{pf:galred:divall}(ii) and (iii).
\end{proof}

\subsubsection{Congruences}\label{sec:pf:galred:cong}

The density of seeds is proved by approximating $\ell$-adically, in a space of algebraic automorphic forms, a given form by
forms in the span of representations with the required properties, and translating such approximations into polynomial
approximations of Hecke eigenvalue systems.

For a level $K$ (an open compact subgroup of $\galredU_a(\A_{E^+,f})$) and a weight $W$, let $\mathcal A(W,K)$ be the
finite-dimensional space of $K$-invariant algebraic automorphic forms of weight $W$, and for a property $\mathcal P$ of
automorphic representations let $\mathcal A_{\mathcal P}(W,K)$ be the span of the $K$-fixed parts of the representations
with property $\mathcal P$. $\mathcal P$ is \emph{dense} at $(W,K)$ if every $f\in\mathcal A(W,K)$ can be approximated
$\ell$-adically by elements of $\mathcal A_{\mathcal P}(W,K)$: for every $\varepsilon>0$ there is
$G\in\mathcal A_{\mathcal P}(W,K)$ with $f-G=e\cdot f$ pointwise for a function $e$ of $\ell$-adic size at most
$\varepsilon$. The following statements are used; (ii) and (iii) are proved in \cref{sec:pf:galred:aux}, the others in
the development.

\begin{enumerate}[label=(\roman*),nosep]
\item (\lref{UpperTheory.Gal.Gen.toPolyX_holds}) If $\mathcal P$ is dense at $(W,K)$ and a family of $\ell$-adically
  integral Hecke operators $t_j$ of level $K$ acts on the $K$-fixed vectors of every automorphic representation $U$ by
  scalars $\eta_j(U)$, then for every $U$ of weight $W$ with nonzero $K$-fixed vectors, the system $\eta(U)$ is
  polynomially approximable by the systems $\eta(U')$ of the $U'$ of weight $W$ and level $K$ with property $\mathcal P$.
\item (\cref{pf:galred:densegen}) Genericity at finitely many split places apart from $w_0$ of the Satake parameters, in
  the sense of the nonvanishing of a fixed nonzero polynomial, is dense, for every weight and level away from $w_0$.
\item (\cref{pf:galred:wtlac}) \emph{Lacunary weight families.} The eigenvalue system of $U$ of weight $\lambda$ and level
  $K\cap K'_{v'}$ ($K'_{v'}$ an open subgroup at $v'$) is polynomially approximable by eigenvalue systems of
  representations $U'$ of level $K$ whose weights $\mu$ agree with $\lambda$ at the embeddings not lying over $v'$ or
  $c(v')$, and such that $J_{\mu(\tau_0)}$ is gapped for a fixed $\tau_0$ lying over $v'$.
\item (\lref{UpperTheory.Gal.Int.Cong.ellSt_holds}) Types of $\ell$-power order exist in every neighbourhood of $1$ in
  $\GL_a(E_{w_0})$, $w_0\mid\ell$, with $m$ as large as required; combined with (iii) and the density of forms of a given type
  (\lref{UpperTheory.Gal.Gen.tyStep_holds}) this gives the statement (iii) with the additional property that $U'$ has a
  type of $\ell$-power order at $w_0$ (\lref{UpperTheory.Gal.Int.Cong.stepSt_of_ex}).
\item (\lref{UpperTheory.Gal.Gen.combine_holds}) If $\theta$ is polynomially approximable
  by $\mathcal S$, every $\eta\in\mathcal S$ is polynomially approximable by $\mathcal T$ and satisfies $Q(\eta)\ne0$ for a
  fixed integral polynomial $Q$, then $\theta$ is polynomially approximable by $\{\eta\in\mathcal T:Q(\eta)\ne0\}$; and
  polynomial approximability is preserved by forgetting coordinates.
\item (\lref{GlobalTheory.Cong.finU_holds}) If the operators have rational values, the eigenvalues lie in a finite
  extension of $\Ql$.
\end{enumerate}

\begin{proof}[Proof that $\sfname{CongFinOdd}$ holds]
Fix the data and $U$. A level $K$ away from $w_0$, an open subgroup $K_{v'}$ at $v'$ and local levels $(K_w)$ with the
required properties, such that $U$ has nonzero $K\cap K_{v'}$-invariant vectors, exist
(\lref{UpperTheory.Gal.Int.Cong.exists_level}); let $\lambda$ be the weight of $U$. Given $Z$ and $W$, choose: at each
$w\in W$ and $n\in\Z$ a Hecke operator of level $K$ supported at $w$ acting on spherical vectors by the power sum
$p_n(\cdot,w)$, normalized by a nonzero constant $c(w,n)$ to be $\ell$-adically integral
(\lref{UpperTheory.Gal.Int.Cong.exists_powOp} and \lref{UpperTheory.Gal.Int.Cong.exists_scale}); at each $z\in Z$ a
\emph{test operator} acting on spherical vectors by the value of $D_a$ at the Satake parameters, similarly normalized
(\lref{UpperTheory.Gal.Int.Cong.exists_testOp}). Together with the given family of local Hecke operators this is a
family $t$ of integral operators of level $K\cap K_{v'}$, with eigenvalue systems $\eta(\cdot)$.

By (i) and (ii), $\eta(U)$ is polynomially approximable by the systems $\eta(U_1)$ of representations $U_1$ of weight
$\lambda$ and level $K\cap K_{v'}$ which are generic at the places of $Z$. By (iv), each such $\eta(U_1)$ is polynomially
approximable by the systems $\eta(U')$ of representations $U'$ of level $K$, with lacunary weight at $v'$ relative to
$U$, the same weights as $U$ away from $v'$, and a type of $\ell$-power order at $w_0$. With $Q$ the product of the test
coordinates, which is nonzero exactly when $U'$ is generic at all places of $Z$, (v) shows that $\eta(U)$ is polynomially
approximable by the systems of such $U'$ which are generic at $Z$; forgetting the test coordinates gives the
approximation required in $\sfname{CongFinOdd}$. The $U'$ obtained are seeds of the required level: they are unramified
off $T$ and spherical at $v'$ and at the split places apart from $w_0,v'$ where $U$ is, because they have nonzero
$K$-fixed vectors, and genericity at the places of $W$ is obtained in the same way by including $W$ among the test
places. The finiteness assertion is (vi).
\end{proof}

\subsubsection{Limits}\label{sec:pf:galred:lim}

\begin{theorem}\label{pf:galred:lim}
$\sfname{SeedFinOdd}\wedge\sfname{CongFinOdd}$ implies $\sfname{LimFinOdd}$.
\formallabel{UpperTheory.Gal.Lim.limFinOdd_of}
\end{theorem}

The proof is an $\ell$-adic interpolation of Galois pseudo-characters in the manner of \cite{Taylor1991}. Write
$\galredtr\rho$ for the trace function of a representation $\rho$, and say that $r$ is \emph{linearly approximable} by a set
$\mathcal C$ of representations if for every $\varepsilon>0$ there are finitely many $\rho\in\mathcal C$ and $c_\rho$ with
$|\galredtr r(g)-\sum c_\rho\galredtr\rho(g)|\le\varepsilon$ for all $g$. We use the following results; (iii) and (v) are
proved in \cref{sec:pf:galred:aux}, the others in the development.

\begin{enumerate}[label=(\roman*),nosep]
\item (\lref{GlobalTheory.FrobOpsIdx.frobOpsIdx}) For a level $(K_w)$ and the set $W_0$ of split places $w\nmid\ell$ with
  $K_w=\GL_a(\OO_{E_w})$, there is a family of local Hecke operators indexed by $J$, with places $w_j\in W_0$, rational
  values, Frobenius elements $s_j$ at $w_j$ and exponents $n_j\ge1$, such that: every arithmetic Frobenius at a place of
  $W_0$ occurs with $n_j=1$; if $\sigma$ is spherical with Satake parameter $B$ at $w_j$ and $\rho$ satisfies
  $\iota\det(X-\rho(s_j^{-1}))=\prod_{B^{(a)}}(X-b)$, then $\iota^{-1}z_j(\sigma)=\galredtr\rho(x_j)$, $x_j=s_j^{-n_j}$; and conversely
  if $\galredtr r(x_j)=\iota^{-1}z_j(\sigma)$ for all $j$ with $w_j=w$, then $\iota\det(X-r(s^{-1}))=\prod_{B^{(a)}_w}(X-b)$ for every
  arithmetic Frobenius $s$ at $w$.
\item (\lref{UpperTheory.Gal.Lim.exists_trace_limit}) Let $\Gamma$ be a topological group, $x\colon J\to\Gamma$,
  $\theta\colon J\to L$ with $L/\Ql$ finite, $(\rho_i)_{i\in I}$ continuous representations, $\eta_i\colon J\to\Qlbar$ and
  $D_i\subset J$ with $\eta_i(j)=\galredtr\rho_i(x_j)$ for $j\in D_i$, such that $x(\bigcap_{i\in s}D_i)$ is dense for every finite
  $s\subset I$, and $\theta$ is polynomially approximable by $\{\eta_i\}$. Then there is $t\colon\Gamma\to L$ such that
  every polynomial approximation of $\theta$ by a combination $\sum c_i\eta_i$ with error $\varepsilon$ gives
  $|t-\sum c_i\galredtr\rho_i|\le\varepsilon$ on $\Gamma$, and for every $A\subset I$ with $\theta$ polynomially approximable by
  $\eta(A)$, $t$ is polynomially approximable by $\{\galredtr\rho_i:i\in A\}$ and $t(x_j)=\theta_j$ whenever $j\in D_i$ for all
  $i\in A$.
\item (\cref{pf:galred:limrep}) Let $\Gamma$ be compact, $x\colon J\to\Gamma$ with dense image,
  $\theta\colon J\to L$, and $\mathcal C$ a set of continuous representations $\Gamma\to\GL_a(\Qlbar)$ such that $\theta$ is
  polynomially approximable by the functions $j\mapsto\galredtr\rho(x_j)$, $\rho\in\mathcal C$. Then there is a continuous
  semisimple $r$ with $\galredtr r(x_j)=\theta_j$, and $\galredtr r$ is polynomially approximable by $\{\galredtr\rho\}$.
  Polynomial approximability of traces implies linear approximability (apply the definition to the coordinate polynomials
  $X_g$, $g\in\Gamma$).
\item (\lref{GlobalTheory.Cong.ker_of_linApprox}) If $r$ is semisimple and linearly approximable by $\mathcal C$, and
  every $\rho\in\mathcal C$ is trivial at $s$, then $r(s)=1$.
\item (\cref{pf:galred:closedAtL}) \emph{Closedness at $\ell$.} If $r$ is continuous, semisimple and
  linearly approximable by a set $\mathcal C$ of continuous representations which are semistable at a frame with a
  common $\tau$-Hodge--Tate multiset $J_\tau$ for every $\tau$ induced by the frame, then $r$ is de Rham with $\tau$-Hodge--Tate
  multisets $J_\tau$; if the members of $\mathcal C$ are crystalline, so is $r$. Hence
  (\lref{GlobalTheory.CongL.crysDRU_of_linApprox}) $\sfname{CrysDR}$ passes to linear limits of representations
  satisfying $\sfname{CrysDR}$ for representations with the same weights at the embeddings over $w$.
\end{enumerate}

\begin{proof}[Proof of \cref{pf:galred:lim}]
Let $U$ be unramified off $T$. Take the level $(K_w)$ of $\sfname{CongFinOdd}$, the finite set $S$ of split places where
$U$ is not spherical, and the family of Hecke operators of (i) for the set
$W_0$ of split places $w\nmid\ell$ where $U$ is spherical (at such places apart from $w_0,v'$ the level is maximal). Let
$\theta=\theta_U$. For every finite set $Y$ of places, $\sfname{SeedFinOdd}$ (with $U_0=U$ and this $S$) gives a set
$Z_0(Y)$, and for every seed $U'$ a representation $\rho(Y,U')$; let $Z(Y)=Z_0(Y)\cup\{w\in Y:\text{split},\ w\nmid\ell,\ U
\text{ spherical at }w\}$ and $W(Y)$ the set of places of $Y$ above $\ell$, apart from $w_0,v'$, where $U$ is spherical. Let
$I$ be the set of triples $i=(Y,U',\theta')$ where $U'$ is a seed of level $(K_w)$, generic at $Z(Y)\cup W(Y)$, with the same
weights as $U$ away from $w_0,v'$, and $\theta'=\theta_{U'}$; put $\rho_i=\rho(Y,U')$, $\eta_i=\theta'$ and $D_i=\{j:
\eta_i(j)=\galredtr\rho_i(x_j)\}$. By $\sfname{CongFinOdd}$ applied with $Z(Y),W(Y)$, for every $Y$ the system $\theta$
(extended by the normalized power sums at $W(Y)$) is polynomially approximable by the systems of the $i$ with first
component $Y$.

For every finite $s\subset I$, the set $x(\bigcap_{i\in s}D_i)$ is dense: by Chebotarev's theorem the inverses of Frobenius
elements at $d$-usable places outside a finite set are dense (\lref{GlobalTheory.CongL.UnitLim.dense_frob_usable_off}),
the open dense sets $V$ of (S1) for the $\rho_i$, $i\in s$,
intersect them in a dense set, and for such a Frobenius element at $w=w_j$ the seed $\rho_i$ has the Frobenius characteristic
polynomial of $\galredSat_w(U')$ (S1), so by (i) $\galredtr\rho_i(x_j)=\eta_i(j)$. The eigenvalues of $\theta$ lie in a finite
extension $L$ of $\Ql$ (the operators of (i) have rational values). Now (ii) gives a function $t$ with $t(x_j)=\theta_j$
approximated by combinations of the $\galredtr\rho_i$, and (iii) gives a continuous semisimple $r$ with $\galredtr r=t$ which is
linearly approximable by $\{\rho_i:i\in A\}$ for every $A\subset I$ for which $\theta$ is polynomially approximable by $\eta(A)$, in
particular for $A=\{i:Y_i=\{w\}\}$ for every place $w$.

(L1): let $w$ be $d$-usable, $w\nmid\ell$, $B=\galredSat_w(U)$, and $A_w=\{i:Y_i=\{w\}\}$. Every $\rho_i$, $i\in A_w$, satisfies
$\galredFr_w(\rho_i)=\galredSat_w(U_i)^{(a)}_w$ by (S2) (the seeds in $A_w$ are generic at $w\in Z(\{w\})$), in particular is trivial on
$I_A$; by (iv), $r(I_A)=1$. For the Frobenius elements, $\galredtr r(x_j)=t(x_j)=\theta_j$ for $w_j=w$, by (ii) applied to $A_w$
(for $i\in A_w$, $j\in D_i$ by (S2) and (i)), and the converse part of (i) gives $\galredFr_w(r)=B^{(a)}_w$.

(L2): for $w\mid\ell$ apart from $w_0,v'$ with $U$ spherical at $w$, the representations $\rho_i$, $i\in A_w$, satisfy
$\sfname{CrysDR}(U_i,\rho_i,w)$ by (S3), and $U_i$ has the same weights as $U$ at the embeddings over $w$; by (v), $r$
satisfies $\sfname{CrysDR}(U,r,w)$.

(L3): for such $w$ and $B=\galredSat_w(U)$, take $\mathcal C=\{\rho_i:i\in A_w\}$; these are crystalline at $w$ with the common
Hodge--Tate multisets $J_{\lambda}$ of the weights of $U$, and by (S4) their Weil--Deligne representations at $w$ have the
characteristic polynomials of $\galredsph_A(\galredSat_w(U_i))$. The joint polynomial approximation of $\theta$ and the normalized
power sums $c(w,n)p_n(U,w)$ by the corresponding data of the $i\in A_w$, together with the uniform approximation of $t$
obtained in (ii), gives the joint approximation of (L3) (\lref{UpperTheory.Gal.Lim.joint}; the power sums of an element
$\sigma\in W_A$ on $\galredsph_A(B)$ are the power sums $p_{-n\deg\sigma}(U,w)$, and the constants $c(w,n)$ are divided out).
\end{proof}

% ---------------------------------------------------------------------------
\subsection{From odd rank to \texorpdfstring{$\sfname{CoreAll}$}{CoreAll}}\label{sec:pf:galred:core}

\begin{theorem}\label{pf:galred:coreallof}
$\sfname{SeedFinOdd}\wedge\sfname{CongFinOdd}\wedge\sfname{TLiftK}$ implies $\sfname{CoreAll}$.
\formallabel{UpperTheory.Gal.Int.coreAll_of}
\end{theorem}

The proof follows the chain \eqref{eq:pf:galred:chain}; it is given at the end of this subsection.

\subsubsection{Isomorphism and transfer}

\begin{lemma}\label{pf:galred:repiso}
Let $F$ be a number field and $r,r'\colon G_F\to\GL_a(\Qlbar)$ continuous and semisimple, such that for all but finitely
many places $w$ of $F$ of prime norm there is a multiset $S$ with $\galredFr_w(r)=\galredFr_w(r')=S$. Then $r\cong r'$. In
particular two representations weakly attached to the same cuspidal $\pi$ are isomorphic, and if $r\cong r'$ then each
clause of \cref{pf:galred:clauses} holds for $r'$ whenever it holds for $r$.
\formallabel{UpperTheory.Gal.Int.Core.isoOne_holds}
\end{lemma}

\begin{proof}
The places of prime norm have Dirichlet density one, so by Chebotarev's theorem the inverses of their Frobenius elements,
outside a finite set of places, form a dense subset of $G_F$ (\lref{GlobalTheory.Published.chebotarevDegreeOne'}). At
these elements $r$ and $r'$ have the same characteristic polynomials, so $\galredtr r=\galredtr r'$ by continuity, and $r\cong
r'$ since semisimple representations in characteristic $0$ are determined by their traces
\cite[\S30]{CurtisReiner1962}. Weakly attached representations
satisfy the hypothesis. The clauses are invariant under conjugation by an element of $\GL_a(\Qlbar)$: (unr) because
characteristic polynomials are, the others because the notions of $\ell$-adic Hodge theory and Fontaine's Weil--Deligne
representation are.
\end{proof}

\subsubsection{Frobenius at \texorpdfstring{$\ell$}{l} in the limit}

\begin{theorem}[$\sfname{ClosedPhi}$]\label{pf:galred:closedphi}
Let $r\colon G_E\to\GL_a(\Qlbar)$, $a\ge1$, be continuous and semisimple, $w$ a place, $A$ a place above $w$ with a frame at
$A$, $\sigma\in W_A$, $J(\tau)$ a multiset of integers for every embedding $\tau$, $\mathcal C$ a set of continuous
representations which are crystalline at the frame with $\tau$-Hodge--Tate multiset $J(\tau)$ for every $\tau$ induced by the
frame, and $S\colon\N\to\Qlbar$ such that $r$ and $S$ are jointly approximated by $\mathcal C$ and the power sums
$j\mapsto\iota^{-1}p_j(\sigma,\galredWD(\rho|_{D_A}))$, $\rho\in\mathcal C$: for every $n$ and $\varepsilon>0$ there are
finitely many $\rho\in\mathcal C$ and $c_\rho$ with $|\galredtr r(h)-\sum c_\rho\galredtr\rho(h)|\le\varepsilon$ for all
$h\in G_E$ and $|S(j)-\sum c_\rho\iota^{-1}p_j(\sigma,\galredWD(\rho|_{D_A}))|\le\varepsilon$ for $j\le n$. Then
$p_n(\sigma,\galredWD(r|_{D_A}))=\iota(S(n))$ for all $n$.
\formallabel{UpperTheory.Gal.ClosedPhiFromClosedTrace.congPhi_holds}
\end{theorem}

\begin{proof}
For $x\in\galredWD_A$, the eigenvalues of $\rho_x(\sigma^n)=\rho_x(\sigma)^n$ are the $n$-th powers of those of
$\rho_x(\sigma)$, so $p_n(\sigma,x)$ is the trace of $\sigma^n$ on $x$, i.e.\ minus the subleading coefficient of
$P_x(\sigma^n)$. Fix $n$. The members of $\mathcal C$ are crystalline, hence semistable. The hypothesis, for the index $n$,
is the joint approximation of \cref{pf:galred:closedtr} for $g=\sigma^n$ and $\theta=S(n)$, so
\cref{pf:galred:closedtr} gives $\iota^{-1}p_n(\sigma,\galredWD(r|_{D_A}))=S(n)$.
\end{proof}

Let $\sfname{LimU}$ be $\sfname{LimFinOdd}$ with (L3) replaced by: $\sfname{Phi}(r,w,B)$ for every $w\mid\ell$ apart from
$w_0,v'$ at which $U$ is spherical with Satake parameter $B$.

\begin{lemma}\label{pf:galred:limU}
$\sfname{LimFinOdd}\wedge\sfname{ClosedPhi}$ implies $\sfname{LimU}$.
\formallabel{UpperTheory.Gal.Int.Core.limU_of_limOdd}
\end{lemma}

\begin{proof}
Given the joint approximation (L3) for $\sigma\in W_A$, $\sfname{ClosedPhi}$ gives $p_n(\sigma,\galredWD(r|_{D_A}))=
p_n(\sigma,\galredsph_A(B))$ for all $n\ge0$, and power sums determine characteristic polynomials in characteristic $0$.
\end{proof}

\subsubsection{All places above \texorpdfstring{$\ell$}{l}}

For a CM field $E$, $\sfname{FieldOK}'(\ell,E)$ means: $\sqrt{-d}\in E$ for some integer $d>0$; $[E^+:\Q]$ is even; $E/E^+$ is
unramified at all finite places; every place above $\ell$ is split; and there are three places above $\ell$ which are
pairwise apart. For a set $Q$ of places and $m\ge1$, $\sfname{CoreAt}(Q,m)$ states: every automorphic $U$ of $\galredU_m$ with
weights $\lambda_\tau$ has a continuous semisimple $r$ satisfying the four clauses at every $w\in Q$ where $U$ is spherical,
with $B=\galredSat_w(U)$, the clause (dR) being taken with the infinitesimal characters $\galredinf(\lambda_\tau)$. Let
$\sfname{CoreUd}$ be: $\sfname{CoreAt}(\{d\text{-usable places}\},m)$ for every odd $m$, every $E$ with
$\sfname{FieldOK}'(\ell,E)$ and every $d$ with $\sqrt{-d}\in E$ and $-d$ a square in $\Q_\ell$.

\begin{theorem}\label{pf:galred:joinU}
$\sfname{LimU}$ implies $\sfname{CoreUd}$.
\formallabel{UpperTheory.Gal.Int.Core.joinU_holds}
\end{theorem}

\begin{proof}
Let $U$ be automorphic on $\galredU_m$, $m$ odd. Then $U$ is unramified off a finite set $T$ of non-split places
(\lref{UpperTheory.Gal.Int.Core.autT_holds}), and $\sfname{FieldOK}'(\ell,E)$ implies $\sfname{FieldOK}(E,\ell,d)$. For
any two places $u,v$ above $\ell$ which are apart, $\sfname{LimU}$ gives $r_{u,v}$ with (L1), (L2) and $\sfname{Phi}$ at the
places above $\ell$ apart from $u,v$ where $U$ is spherical; together these are the four clauses at such places (the clause
(unr) being vacuous above $\ell$; $\sfname{Phi}$ is equivalent to (phi)). Let $w_1,w_2,w_3$ be three pairwise apart places
above $\ell$, and put $r=r_{w_2,w_1}$; it satisfies the clauses at all $d$-usable places prime to $\ell$ by (L1).
Let $w\mid\ell$ be a place where $U$ is spherical. Among $w_1,w_2,w_3$ at most one lies over the same place of $E^+$ as $w$,
so two of them, $u,v$, are apart from $w$ and from each other. The representations $r$ and $r_{u,v}$ have the same Frobenius
multisets at all $d$-usable places prime to $\ell$ where $U$ is spherical. $U$ is spherical at almost all split places
(\lref{UpperTheory.Gal.TP.Sub.satAE_holds}), and almost all places of prime norm are $d$-usable (indeed a place $w$ of prime norm $q$ which is unramified over $E^+$ is split,
since an inert place has norm a square, and $-d$ is a square in $E_w=\Q_q$ because $\sqrt{-d}\in E$). Hence $r\cong r_{u,v}$ by
\cref{pf:galred:repiso}, and the clauses at $w$ transfer from $r_{u,v}$ to $r$.
\end{proof}

\subsubsection{From odd to even rank}\label{sec:pf:galred:twoup}

For a CM field $E$, $\sfname{InertOK}(E)$ means that every non-split place of $E$ lies over a place of $E^+$ which is prime to
$2$, has norm $>3$ and is prime to the different of $E^+$. An algebraic Hecke character $\psi$ of $E$ (of infinity type
$p$, i.e.\ $\psi_\infty(x)=\prod_\tau\tau(x)^{p(\tau)}$ on totally positive units, in the sense of
\cref{pf:galred:isweight}) has an $\ell$-adic character $L(\psi)\colon G_E\to\Qlbar^\times$, continuous, with
$\galredFr_w(L(\psi))=\{\psi(\varpi_w)\}$ at places where $\psi$ is unramified, with the expected behaviour at $\ell$, and
$\psi\mapsto L(\psi)$ is injective (\cite{Serre1968}; \lref{UpperTheory.Gal.TP.Sub.charEx}).

\begin{definition}[Lifting data]\label{pf:galred:twoupdef}
Let $n$ be even and $U$ automorphic on $\galredU_n$ with weights $\lambda$. A \emph{lift} of $U$ of size $t$, with twist
$\kappa$ and auxiliary characters $\psi=(\psi_1,\dots,\psi_t)$, is an automorphic $U'$ of $\galredU_{n+t}$, with weights
$\lambda'$, such that $\kappa$ and the $\psi_i$ are algebraic Hecke characters of $E$, the Hodge--Tate multisets satisfy
$J_{\lambda'_\tau}=(J_{\lambda_\tau}+m_\tau)\cup\{-p_i(\tau)\}_i$ ($m$ the infinity type of $\kappa$, $p_i$ that of $\psi_i$), and
at every split place $w$ where $\kappa$ and the $\psi_i$ are unramified and $U$ is spherical, $U'$ is spherical and
\[
  \galredSat_w(U')^{(n+t)}_w=\{b\,\kappa(\varpi_w)^{-1}:b\in\galredSat_w(U)^{(n)}_w\}\cup\{\psi_i(\varpi_w)\}_i.
\]
$\sfname{TwoUpOK}$ states: for every CM field $E$ with $E/E^+$ unramified at all finite places, $\sqrt{-d}\in E$ for some
$d>0$, $[E^+:\Q]$ even and $\sfname{InertOK}(E)$, every even $n\ge1$, every automorphic $U$ of $\galredU_n$ and every split
place $w_0$, there are $t$ with $n+t$ odd, an algebraic Hecke character $\kappa$, and two lifts $U_1,U_2$ of $U$ of size $t$
with twist $\kappa$ and auxiliary characters $\psi^{(1)},\psi^{(2)}$ such that $\psi^{(1)}_i\ne\psi^{(2)}_j$ for all $i,j$, and
$\kappa$, $\psi^{(1)}_i$ are unramified at $w_0$.
\formallabel{UpperTheory.Gal.TP.TwoUpOK}
\end{definition}

\begin{theorem}\label{pf:galred:twoup}
$\sfname{TLiftK}$ implies $\sfname{TwoUpOK}$.
\formallabel{UpperTheory.Gal.TP.twoUpOK_of_k}
\end{theorem}

\begin{proof}
Let $E$, $n$, $U$ and $w_0$ be given. The hypotheses on $E$ provide the base data of $\sfname{TLiftK}$ (a CM type, a standard
additive character unramified at the inert places, which are prime to $2$ and of norm $>3$ by $\sfname{InertOK}$), the
hermitian space $E^n$ with its standard form and self-dual lattice, and skew-hermitian spaces $(E^{d},\delta\sum x_i\bar y_i)$
with self-dual lattices for all $d$ (\cref{pf:galred:basedata}). Let $d_0$ be the threshold given by
$\sfname{TLiftK}$ for $n$, and $t=2(d_0+n)+1$; then $n+t\ge d_0$ and $n+t$ is odd. Let $b$ be a lower bound for the weights
$\lambda_{w,k}$. There are an algebraic Hecke character $\kappa$ and two pairs $(\chi_X^{(k)},\chi_W^{(k)})$, $k=1,2$, of unitary
Hecke characters with $\chi_X^{(k)}|_{\A_{E^+}^\times}=\eta^n$, $\chi_W^{(k)}|_{\A_{E^+}^\times}=\eta^{n+t}$, archimedean exponents
$m_X^{(k)},m_W^{(k)}$ with $m_W^{(k)}\le 2b-(n+t)$, and families $\psi^{(k)}$ of $t$ algebraic Hecke characters, with
$\psi^{(1)}_i\ne\psi^{(2)}_j$, $\kappa$ and $\psi^{(1)}_i$ unramified at $w_0$, related to $(\chi_X^{(k)},\chi_W^{(k)})$ by a
dictionary which, through the unramified type II correspondence, expresses Satake parameters of lifts in terms of $\kappa$
and $\psi^{(k)}$ (\cref{pf:galred:chrpairs}, applied with the bound $2b-(n+t)$).
The weight condition $n+t+m_W^{(k)}(w)\le2\lambda_{w,j}$ of $\sfname{TLiftK}$ holds, so $\sfname{TLiftK}$ gives automorphic
representations $U_k$ of $\galredU_{n+t}$ with the weights of \cref{def:galois:tliftk}(i) and the local correspondence
(ii) at split places. At split places where everything is unramified, the explicit type II correspondence for unramified
representations computes the Satake parameters of $U_k$ as required, and the weights of $U_k$ give the Hodge--Tate
multisets (\cref{pf:galred:upof}). Thus $U_1,U_2$ are lifts as required.
\end{proof}

A set $Q$ of places is a \emph{guard} if its members are split and, for every finite set $S$ of places, the inverses of the
arithmetic Frobenius elements at the places of $Q\setminus S$ prime to $\ell$ are dense in $G_E$. If $\sqrt{-d}\in E$, the
$d$-usable places form a guard (by the argument in the proof of \cref{pf:galred:joinU} and Chebotarev's theorem).

\begin{lemma}[Subtraction]\label{pf:galred:subtract}
Let $Q$ be a guard, $n$ even, and assume $\sfname{CoreAt}(Q,m)$ for all odd $m$ and the conclusion of $\sfname{TwoUpOK}$ for
$E$ and $n$. Then $\sfname{CoreAt}(Q,n)$ holds.
\formallabel{UpperTheory.Gal.Int.Core.subtractQ_holds}
\end{lemma}

\begin{proof}
Let $U$ be automorphic on $\galredU_n$. Call $r$ \emph{$Q$-attached} to $U$ if it is continuous, semisimple, and has
$\galredFr_w(r)=\galredSat_w(U)^{(n)}_w$ at almost all $w\in Q$ prime to $\ell$ where $U$ is spherical. Two $Q$-attached
representations are isomorphic (as $Q$ is a guard; compare \cref{pf:galred:repiso}), and clauses transfer along
isomorphisms; so it suffices to construct, for every $w_0\in Q$, a $Q$-attached $r$ satisfying the four clauses at $w_0$.

Fix $w_0$ and take $t$, $\kappa$, $U_1,U_2$, $\psi^{(1)},\psi^{(2)}$ from $\sfname{TwoUpOK}$. By $\sfname{CoreAt}(Q,n+t)$ there
are representations $r_1,r_2$ of dimension $n+t$, $Q$-attached to $U_1,U_2$, satisfying the clauses at the places of $Q$. Put
$a^{(k)}_i=L(\psi^{(k)}_i)$ and write $\galredcp(r)(g)=\det(X-r(g))$. At almost all $w\in Q$, the descriptions of the Satake
parameters of $U_1$ and $U_2$ show that $\galredcp(r_1)\prod_j(X-a^{(2)}_j)$ and $\galredcp(r_2)\prod_i(X-a^{(1)}_i)$ agree at
the inverses of Frobenius elements at $w$: both are $\prod(X-\iota^{-1}c)$ over the multiset
$\{b\kappa(\varpi_w)^{-1}\}\cup\{\psi^{(1)}_i(\varpi_w)\}\cup\{\psi^{(2)}_j(\varpi_w)\}$. These elements are dense, so
\begin{equation}\label{eq:pf:galred:cpid}
  \galredcp(r_1)(g)\prod_j\bigl(X-a^{(2)}_j(g)\bigr)=\galredcp(r_2)(g)\prod_i\bigl(X-a^{(1)}_i(g)\bigr)\qquad(g\in G_E).
\end{equation}
The characters $a^{(1)}_i$ and $a^{(2)}_j$ are distinct, since $L$ is injective. We peel off the $a^{(1)}_i$ one at a time
(\lref{UpperTheory.Gal.TP.Sub.peelAll}): by \eqref{eq:pf:galred:cpid} the semisimple representations $r_1\oplus
\bigoplus_ja^{(2)}_j$ and $r_2\oplus\bigoplus_ia^{(1)}_i$ have the same characteristic polynomials, hence are isomorphic; so the
character $a^{(1)}_t$ occurs in $r_1\oplus\bigoplus_ja^{(2)}_j$, and since it differs from all $a^{(2)}_j$ it occurs in $r_1$:
$r_1\cong s_1\oplus a^{(1)}_t$ by semisimplicity, and \eqref{eq:pf:galred:cpid} holds with $(r_1,\psi^{(1)})$ replaced by
$(s_1,(\psi^{(1)}_i)_{i<t})$. After $t$ steps, $r_1\cong s\oplus\bigoplus_ia^{(1)}_i$ with $s$ continuous, semisimple, of
dimension $n$. The clauses at a place $w$ for $r_1$, with Frobenius multiset and Hodge--Tate multisets that are unions of
those of $s$ and of the $a^{(1)}_i$, imply those for $s$ (\lref{UpperTheory.Gal.TP.Sub.clSub}). Put $r=s\otimes L(\kappa)$.
At $w\in Q$ where $\kappa$ and the $\psi^{(1)}_i$ are unramified, the clauses for $r_1$ (with $\galredSat_w(U_1)$ and the
weights of $U_1$) give those for $s$ with Frobenius multiset $\{b\kappa(\varpi_w)^{-1}\}$ and Hodge--Tate multisets
$J_{\lambda_\tau}+m_\tau$, hence those for $r$ with $\galredSat_w(U)^{(n)}_w$ and $J_{\lambda_\tau}$
(\lref{UpperTheory.Gal.TP.Sub.clTw}). This applies at $w_0$ and at almost all $w\in Q$, so $r$ is $Q$-attached to $U$ and
satisfies the clauses at $w_0$.
\end{proof}

Let $\sfname{AllUd}$ be the statement: $\sfname{CoreAt}(\{d\text{-usable places}\},n)$ for every $n\ge1$, every $E$ with
$\sfname{FieldOK}'(\ell,E)$ and $\sfname{InertOK}(E)$, and every $d$ with $\sqrt{-d}\in E$ and $-d$ a square in $\Q_\ell$.

\begin{theorem}\label{pf:galred:joinAll}
$\sfname{TwoUpOK}\wedge\sfname{CoreUd}$ implies $\sfname{AllUd}$.
\formallabel{UpperTheory.Gal.Int.Core.joinAll_holds}
\end{theorem}

\begin{proof}
The $d$-usable places form a guard. For odd $n$ use $\sfname{CoreUd}$; for even $n$ apply \cref{pf:galred:subtract}, with
$\sfname{TwoUpOK}$ (whose hypotheses on $E$ are contained in $\sfname{FieldOK}'(\ell,E)\wedge\sfname{InertOK}(E)$) and
$\sfname{CoreUd}$ in all odd ranks.
\end{proof}

\subsubsection{From unitary groups to \texorpdfstring{$\GL_a$}{GL(a)}}

$\sfname{DescOffT}$ is the statement: for every CM field $E$, prime $\ell$, $a\ge1$, good conjugate self-dual $\pi$ (in scope
$\mathrm{csd}$ with $\chi=\eta^a$) and finite set $Z$ of places, there are a finite set $T$ of non-split places disjoint from
$Z$ and an automorphic $U$ of $\galredU_a$, unramified off $T$, such that every local component of $\pi$ at a split place is a
local component of $U$ there, and $\galredinf_\tau(\pi)=\galredinf(\lambda)$ for every weight $\lambda$ of $U$ at $\tau$. It is
proved in the development, from the Jacquet--Shalika bound and multiplicity one, in the style of
\cite{Labesse2011,ClozelHarrisLabesse2011} (\lref{UpperTheory.Gal.Desc.descOffT}, which states exactly this).

Let $\sfname{CoreGLu}$ be the statement: for every $E$ with $\sfname{FieldOK}'(\ell,E)$ and $\sfname{InertOK}(E)$, every $d$ with
$\sqrt{-d}\in E$ and $-d$ a square in $\Q_\ell$, and every $\pi$ in scope $\mathrm{csd}$ (with $\chi=\eta^a$), there is a continuous
semisimple $r$ satisfying the four clauses at every $d$-usable place $w$ where $\pi$ is unramified, with
$B=\galredSat_w(\pi)$.

\begin{theorem}\label{pf:galred:joinGL}
$\sfname{DescOffT}\wedge\sfname{AllUd}$ implies $\sfname{CoreGLu}$.
\formallabel{UpperTheory.Gal.Int.Core.joinGL_holds}
\end{theorem}

\begin{proof}
Descend $\pi$ to $U$ by $\sfname{DescOffT}$ (with $Z=\emptyset$), let $\lambda$ be the weights of $U$, and take $r$ from
$\sfname{AllUd}$ for $U$. Let $w$ be $d$-usable with $\galredSat_w(\pi)=B$. A local component of $\pi$ at $w$ is a local
component of $U$ at $w$, spherical with Satake parameter $B$, so $\galredSat_w(U)=B$ and the clauses for $U$ at $w$ hold; they
are the clauses for $\pi$, because $\galredinf_\tau(\pi)=\galredinf(\lambda_\tau)$.
\end{proof}

\subsubsection{Removing the conditions on the field}

Let $\sfname{RowCsd}$ be the statement: for every CM field $E$, $a\ge1$ and every regular algebraic $\pi$ polarized by
$\eta^a$ (in scope $\mathrm{all}$ with $\chi=\eta^a$), there is a continuous semisimple $r$ satisfying the four clauses at
\emph{every} place $w$ where $\pi$ is unramified. Let $\sfname{QuadLift}$ be the statement: if $M/K$ is a quadratic extension of
totally complex fields and $\pi$ is cuspidal on $\GL_a(\A_K)$, and $M/K$ is ramified at some place where $\pi$ is unramified,
then $\pi$ has a cuspidal lift to $M$ (\cref{pf:galred:lifts}). It follows from the strong form of cyclic base change
\cite{ArthurClozel1989} and Jacquet--Shalika rigidity \cite{JacquetShalika1981}
(\lref{GlobalTheory.QuadLiftBridge.quadLift}; the ramification hypothesis excludes that $\pi$ is a twist of itself by the
quadratic character of $M/K$).

For a number field $F$, a \emph{quadratic CM extension} is a quadratic extension $K/F$ inside $\Omega$ which is a CM field. We
use the following results of the development.
\begin{enumerate}[label=(\roman*),nosep]
\item (\lref{UpperTheory.Gal.Int.Core.patchOne_holds}) \emph{Patching.} Let $\mathcal B$ be a property of pairs (place of $F$,
  multiset), with at most one multiset for each place, $\mathcal S$ an infinite set of quadratic CM extensions $K/F$, and
  for each $K\in\mathcal S$ let $r_K\colon G_K\to\GL_a(\Qlbar)$ be continuous and semisimple, such that at almost all places $u$
  of $K$ of prime norm, $\galredFr_u(r_K)=B^{(a)}_u$ for the multiset $B$ with $\mathcal B(u|_F,B)$. Then there is a continuous
  semisimple $\rho\colon G_F\to\GL_a(\Qlbar)$ with $r_K\cong\rho|_{G_K}$ for all but finitely many $K\in\mathcal S$. The proof
  extends each $r_K$ to $G_F$ using the invariance of its trace under conjugation, and patches the extensions along the
  infinitely many quadratic characters (\lref{GlobalTheory.IndexTwoPatch.main}, which states: for infinitely many distinct
  quadratic characters $\eta_K$ of a topological group $\Gamma$ with open kernels and continuous semisimple representations of
  the kernels whose traces are invariant under conjugation by $\Gamma$ and agree on intersections of kernels, there is a
  continuous semisimple representation of $\Gamma$ whose trace agrees with almost all of them); the hypotheses on traces
  follow from \cref{pf:galred:repiso} applied over $K$ and over the compositum of two such fields.
\item \emph{Descent of clauses} (\lref{UpperTheory.Gal.Int.Core.downUnr_holds} and its analogues for (dR), (crys), (phi)).
  If $P$ is a lift of $\pi$ to a quadratic CM extension $K/E$, $w$ a place of $E$ whose decomposition groups lie in $G_K$, and
  the clause $(c)$ holds for $P$ and $\rho|_{G_K}$ at every place $u$ of $K$ above $w$, then it holds for $\pi$ and $\rho$ at $w$.
\end{enumerate}

A quadratic CM extension $K/E$ \emph{serves} $\pi$ at $w_0$ if the decomposition groups at $w_0$ lie in $G_K$, $K/E$ is
ramified at some place where $\pi$ is unramified, and for every lift $P$ of $\pi$ to $K$ there are a set $p$ of places of $K$
containing almost all places of prime norm and all places above $w_0$, and a continuous semisimple representation of
$G_K$ satisfying the clauses for $P$ at the places of $p$ where $P$ is unramified.

\begin{lemma}[Step down]\label{pf:galred:stepdown}
Assume $\sfname{QuadLift}$. Let $\pi$ be cuspidal on $\GL_a(\A_E)$ such that for every place $w_0$ infinitely many quadratic CM
extensions serve $\pi$ at $w_0$. Then there is a continuous semisimple $\rho$ satisfying the four clauses for $\pi$ at every
place where $\pi$ is unramified.
\formallabel{UpperTheory.Gal.Int.Core.stepDown_holds}
\end{lemma}

\begin{proof}
For $K$ serving $\pi$ at some place, $\sfname{QuadLift}$ gives a lift $P$ of $\pi$ to $K$, and the definition gives $r_K$; at
almost all places $u$ of $K$ of prime norm, $u$ has residue degree one over $E$, $\pi$ is unramified at the place below with
Satake parameter $B$, $P$ is unramified at $u$ with Satake parameter $B$, and $\galredFr_u(r_K)=B^{(a)}_u$. Let $\mathcal S$ be
the (infinite) set of such $K$. Patching (i) gives $\rho$ with $r_K\cong\rho|_{G_K}$ for almost all $K\in\mathcal S$. Given a place
$w$ where $\pi$ is unramified and a clause $c$, choose $K$ serving $\pi$ at $w$ outside the finite exceptional set, with lift
$P$ and representation $r'_K$. The representations $r'_K$ and $r_K$ are both Frobenius-compatible with $P$ at almost all
places of prime norm, so $r'_K\cong r_K\cong\rho|_{G_K}$ (\cref{pf:galred:repiso}), and the clause holds for $P$ and
$\rho|_{G_K}$ at the places of $K$ above $w$; by descent (ii) it holds for $\pi$ and $\rho$ at $w$.
\end{proof}

\begin{theorem}\label{pf:galred:joinUp}
Jacquet--Shalika rigidity, the Gelfand--Kazhdan theorem, $\sfname{QuadLift}$ and $\sfname{CoreGLu}$ imply
$\sfname{RowCsd}$.
\formallabel{UpperTheory.Gal.Int.Core.joinUp_holds}
\end{theorem}

\begin{proof}
Let $\pi$ be regular algebraic and polarized by $\eta^a$ over $E$. If $P$ is a lift of such a $\pi$ to a CM extension, then $P$
is regular algebraic and polarized by $\eta^a$ (\lref{UpperTheory.Gal.Int.Core.csdLift_holds}, from Jacquet--Shalika
rigidity and the Gelfand--Kazhdan theorem). We apply \cref{pf:galred:stepdown} three times, along towers of quadratic CM
extensions constructed in the development:
\begin{enumerate}[label=(\alph*),nosep]
\item for every CM field $F$ and finite set $S$ of places of $F$ there are infinitely many \emph{real quadratic steps}
  $K=F(\delta)$, $\delta^2$ a positive integer, which are CM fields in which the places of $S$ split (and the places of $F^+$ below
  them split in $F^+(\delta)$), and which are ramified at some place outside $S$
  (\lref{UpperTheory.Gal.UnramifiedFromSplittingCriterion.infinitely_many_real_quadratic_steps});
\item if $E\subset K_1\subset K_2$ are two such steps, split at a place above $\ell$, then every CM extension of $K_2$ has three
  pairwise apart places above $\ell$ (\lref{UpperTheory.Gal.UnramifiedFromSplittingCriterion.threeApart_holds});
\item if $K_2$ is a real quadratic step over a CM field, then for every finite set $S$ of places of $K_2$ there is a quadratic CM
  extension $K/K_2$, ramified at some place outside $S$, in which the places of $S$ split, with $\sqrt{-d}\in K$ for some $d>0$,
  $[K^+:\Q]$ even, $\sfname{InertOK}(K)$, every non-split place of $K$ unramified over $K^+$ and lying over a place outside
  $S$, and every place of $K$ above $S$ $d$-usable (\lref{UpperTheory.Gal.Int.Core.exists_quadratic_CM_ext_split_at}); $K/K^+$
  is then unramified at all finite places (\lref{UpperTheory.Gal.UnramifiedFromSplittingCriterion.allUnr_holds}).
\end{enumerate}
A step which is ramified at a place outside a set containing the places where $\pi$ is ramified is ramified at a place where
$\pi$ is unramified, so $\sfname{QuadLift}$ applies. In the innermost step
(c), over $K_2$ and for a lift $P_2$ of $\pi$ to $K_2$, the set $S$ is taken to contain the given place $w_0$, the places above
$\ell$ and the places where $P_2$ is ramified. Then $K$ satisfies $\sfname{FieldOK}'(\ell,K)$, $-d$ is a square in $\Q_\ell$,
and every non-split place of $K$ is prime to $\ell$, unramified over $K^+$ and a place where the lift $P$ of $P_2$ to $K$ is
unramified, so $P$ is good and conjugate self-dual; $\sfname{CoreGLu}$ gives a representation satisfying the clauses for $P$ at
the $d$-usable places, which include almost all places of prime norm (see the proof of \cref{pf:galred:joinU}) and the
places above $w_0$. So $K$ serves $P_2$ at $w_0$, and \cref{pf:galred:stepdown} applies over $K_2$ (for every $P_2$), then over
$K_1$ with the steps of (a) as serving extensions, then over $E$.
\end{proof}

\begin{lemma}\label{pf:galred:joinEnd}
$\sfname{RowCsd}$ implies $\sfname{Core}$.
\formallabel{UpperTheory.Gal.Int.Core.joinEnd_holds}
\formallabel{GlobalTheory.RoadW.core_of_rowStarCore}
\end{lemma}

\begin{proof}
Let $(\pi,\chi)$ be in scope $\mathrm{csd}$; then $\chi=\eta^a$ and $\sfname{RowCsd}$ gives $r$ satisfying the clauses at all
places where $\pi$ is unramified. Since $\pi$ is unramified at almost all places and almost all places are prime to $\ell$, the
clause (unr) shows that $r$ is weakly attached to $\pi$; this is $\sfname{Ex}(\mathrm{csd})$. If $r'$ is weakly attached to
$\pi$, then $r'\cong r$ by \cref{pf:galred:repiso}, and the clauses transfer from $r$ to $r'$; this is
$\sfname{SLC}(\mathrm{csd},c)$ for all $c$.
\end{proof}

\begin{theorem}\label{pf:galred:toCoreK}
$\sfname{TLiftK}\wedge\sfname{ClosedPhi}\wedge\sfname{DescOffT}\wedge\sfname{LimFinOdd}$ implies $\sfname{Core}$.
\formallabel{UpperTheory.Gal.Int.Core.toCoreK}
\end{theorem}

\begin{proof}
$\sfname{LimFinOdd}\wedge\sfname{ClosedPhi}$ gives $\sfname{LimU}$ (\cref{pf:galred:limU}), hence $\sfname{CoreUd}$
(\cref{pf:galred:joinU}); $\sfname{TLiftK}$ gives $\sfname{TwoUpOK}$ (\cref{pf:galred:twoup}), hence $\sfname{AllUd}$
(\cref{pf:galred:joinAll}); with $\sfname{DescOffT}$ we get $\sfname{CoreGLu}$ (\cref{pf:galred:joinGL}). Jacquet--Shalika
rigidity, the Gelfand--Kazhdan theorem \cite{JacquetShalika1981,GelfandKazhdan1975} and $\sfname{QuadLift}$ are proved in the
development, so \cref{pf:galred:joinUp} gives $\sfname{RowCsd}$, and \cref{pf:galred:joinEnd} gives $\sfname{Core}$.
\end{proof}

\subsubsection{From \texorpdfstring{$\sfname{Core}$}{Core} to \texorpdfstring{$\sfname{CoreAll}$}{CoreAll}}

\begin{theorem}\label{pf:galred:coreallofcore}
$\sfname{Core}$ implies $\sfname{CoreAll}$.
\formallabel{UpperTheory.Gal.Int.Inp.coreAll_of_core}
\end{theorem}

\begin{proof}
We use the package $\sfname{Inp}$ of \cref{sec:pf:galred:R1}. There are four steps.

\emph{From $\mathrm{csd}$ to $\mathrm{good}$} (\lref{GlobalTheory.GalCut.ex_good_of_ex_csd}). Let $(\pi,\chi)$ be in scope
$\mathrm{good}$, and $w_0$ a place above $\ell$ (split, by goodness). There is an algebraic Hecke character $\kappa$ of $E$,
unramified at $w_0$, such that $(\pi\otimes\kappa,\eta^a)$ is in scope $\mathrm{csd}$
(\lref{GlobalTheory.T.exists_good_conjSelfDual_twist_split_unramified}). $\sfname{Ex}(\mathrm{csd})$ gives $r'$ weakly
attached to $\pi\otimes\kappa$, and $r'\otimes L(\kappa^{-1})$, $L(\kappa^{-1})$ the $\ell$-adic character of $\kappa^{-1}$
\cite{Serre1968}, is weakly attached to $\pi$. So $\sfname{Ex}(\mathrm{good})$ holds. In the same way the four clauses pass from
level $\mathrm{csd}$ to level $\mathrm{good}$ (\lref{UpperTheory.Gal.Int.Inp.holdsS_good_unr}, \lref{UpperTheory.Gal.Int.Inp.holdsS_good_crys}, \lref{UpperTheory.Gal.Int.Inp.holdsS_good_dR_of_csd}, and \lref{UpperTheory.Gal.Int.Inp.holdsS_good_phi}),
using at $\ell$ the Weil--Deligne representations of the characters $L(\kappa)$ and the theorem that de Rham representations are
potentially semistable \cite{Berger2002}.

\emph{Existence at level $\mathrm{all}$} (\lref{UpperTheory.Gal.Int.Inp.ex_all_of_good}). Let $(\pi,\chi)$ be in scope
$\mathrm{all}$. For every place $w_0$ there are infinitely many quadratic CM extensions $K/E$ such that the decomposition groups at
$w_0$ lie in $G_K$, the places of $K$ above $w_0$ are split over $K^+$, $K/E$ is ramified at some place where $\pi$ is unramified,
and every lift of $\pi$ to $K$ is good (\lref{GlobalTheory.GalCut.infinite_qcm_split_good}). By $\sfname{QuadLift}$, $\pi$ has
a lift $P_K$ to each such $K$; $P_K$ is regular algebraic and polarized (\lref{GlobalTheory.LiftPolar.main}), hence in scope
$\mathrm{good}$, and $\sfname{Ex}(\mathrm{good})$ gives $r_K$ weakly attached to $P_K$. Patching
(\lref{GlobalTheory.CMBasePatch.main}, the analogue of (i) above) gives $\rho\colon G_E\to\GL_a(\Qlbar)$, continuous and
semisimple, with $r_K\cong\rho|_{G_K}$ for almost all $K$. Then $\rho$ is unramified at almost all places, since a representation
of $G_E$ whose restriction to $G_K$ is unramified at the places above $w$ is unramified at $w$ when $w$ is unramified in $K$ (then $I_A\subset G_K$). At a place $w\nmid\ell$ where $\pi$ is unramified with Satake parameter
$B$ and $\rho$ is unramified, choose $K$ as above for $w_0=w$, outside the exceptional set: the places $u$ of $K$ above $w$ are split
over $K^+$, $P_K$ is unramified at $u$ with Satake parameter $B$ (residue degree one), $\sfname{SLC}(\mathrm{good},\mathrm{unr})$
gives $\galredFr_u(r_K)=B^{(a)}_u$, and as $D_A\subset G_K$ this gives $\galredFr_w(\rho)=B^{(a)}_w$
(an arithmetic Frobenius at $A$ for $E$ is one for $K$, and $\rho|_{G_K}\cong r_K$). So $\rho$ is weakly attached to $\pi$.

\emph{The clauses at level $\mathrm{all}$} (\lref{UpperTheory.Gal.Int.Inp.holdsS_all_unr} and its analogues
\lref{UpperTheory.Gal.Int.Inp.holdsS_dR_all_of_good}, \lref{UpperTheory.Gal.Int.Inp.holdsS_all_crys}, and \lref{UpperTheory.Gal.Int.Inp.holdsS_all_phi}).
Let $r$ be weakly attached to $\pi$ in scope $\mathrm{all}$, and $w$ a place where $\pi$ is unramified with Satake parameter $B$.
Choose $K$ as above for $w_0=w$ and a lift $P$ of $\pi$ to $K$; $P$ is in scope $\mathrm{good}$ and $r|_{G_K}$ is weakly attached to
$P$. Since $D_A\subset G_K$, the place $u$ of $K$ below $A$ has $q_u=q_w$; at $u$,
$P$ is unramified with Satake parameter $B$; $u$ is split over $K^+$, so the clauses at level $\mathrm{good}$ hold for $P$ and
$r|_{G_K}$ at $u$, and since $D_A\subset G_K$ they are the clauses for $\pi$ and $r$ at $w$. For (unr): $r(I_A)=1$ and
$\iota\det(X-r(\sigma^{-1}))=\prod_{B^{(a)}_w}(X-b)$ for every arithmetic Frobenius $\sigma$ at $A$, because $I_A$ and $\sigma$ lie in
$G_K$ and $\sigma$ is an arithmetic Frobenius at $A$ for $K$. The clauses (dR), (crys) and (phi) are transferred in the same way from
level $\mathrm{good}$; the transfer of (phi) also uses (crys) at level $\mathrm{good}$.
\end{proof}

\begin{proof}[Proof of \cref{pf:galred:coreallof}]
$\sfname{SeedFinOdd}\wedge\sfname{CongFinOdd}$ implies $\sfname{LimFinOdd}$ (\cref{pf:galred:lim}). $\sfname{ClosedPhi}$ holds
(\cref{pf:galred:closedphi}) and $\sfname{DescOffT}$ holds
(\lref{UpperTheory.Gal.Desc.descOffT}), so with $\sfname{TLiftK}$, \cref{pf:galred:toCoreK} gives $\sfname{Core}$, and
\cref{pf:galred:coreallofcore} gives $\sfname{CoreAll}$.
\end{proof}

\begin{proof}[Proof of \cref{pf:galred:seedcong}]
The first assertion is proved in \cref{sec:pf:galred:seeds}. For the second, $\sfname{CongFinOdd}$ is the instance, for the
type condition ``type of $\ell$-power order at $w_0$'', genericity defined by $D_a$, and lacunary weights, of a statement
proved for arbitrary rank (\cref{sec:pf:galred:cong}); the lacunary weight families exist (\cref{pf:galred:wtlac}).
\end{proof}

% ---------------------------------------------------------------------------
\subsection{Proofs of the auxiliary results}\label{sec:pf:galred:aux}

This subsection proves the auxiliary results to which the preceding subsections refer: facts about places, CM
extensions and local components used in \cref{sec:pf:galred:rest}; the division theorem, the density and family
statements of the congruence argument and the pseudo-character and closedness statements of the limit argument
(\cref{sec:pf:galred:odd}); and the data for the two theta lifts of \cref{pf:galred:twoup}. In each proof the steps are
given in full, and the deeper ingredients proved in the development are stated and cited.

\subsubsection{Places of degree one}

\begin{lemma}\label{pf:galred:gnotsq}
Let $E$ be a number field, $w$ a finite place of $E$ and $g$ a prime. If there is a continuous ring isomorphism
$i\colon\Q_g\to E_w$, then $g\notin w^2$, i.e.\ $w$ is unramified over $g$.
\formallabel{DiscreteGroups.sglob.DzType.g_not_mem_sq}
\end{lemma}

\begin{proof}
Let $\varpi\in\OO_E$ have $w$-adic valuation $1$ and put $x=i^{-1}(\varpi)\in\Q_g$. A continuous ring isomorphism of
nonarchimedean local fields preserves the condition $|y|<1$ (\lref{DiscreteGroups.sglob.DzIso.norm_lt_iff}) and maps
elements of absolute value at most $1$ to such elements (\lref{DiscreteGroups.sglob.DzIso.norm_le_of_le}). Hence
$|x|_g<1$, so $|x|_g\le g^{-1}$ and $x/g\in\Z_g$, and $i(x/g)$ is integral in $E_w$. Now $\varpi=i(x)=g\cdot i(x/g)$, so
$1=\ord_w(\varpi)\ge\ord_w(g)$. If $g\in w^2$ then $\ord_w(g)\ge2$, a contradiction.
\end{proof}

\subsubsection{CM extensions for spherical lifts}

\begin{lemma}\label{pf:galred:cmlift}
Assume $\sfname{Inq}$. Let $E$ be a CM field, $(\pi,\chi)$ in scope $\mathrm{all}$ on $\GL_n$, $n\ge1$, $r$ weakly
attached to $\pi$, $A$ a place of $\Omega$ above a finite place $w$ of $E$, $\xi_1,\dots,\xi_n$ smooth characters of
$E_w^\times$, $a$ a local Artin map at $A$ and $g\in W_A$ with $\deg g\ne0$. There are a CM field $M$ with $E\subset
M\subset\Omega$, fixed by $g$, and a cuspidal character lift $P$ of $\pi$ to $M$, with $(P,\chi')$ in scope
$\mathrm{all}$ for some $\chi'$, such that $\xi_j(a(t))=1$ for all $j$ and every $t\in W_A$ of degree $0$ fixing $M$.
\formallabel{GalTw.GlobalTheory.RoadU.TorLiftC.exists_cm_lift_fix_charKer}
\end{lemma}

\begin{proof}
For a field $K\subset\Omega$ and an open normal subgroup $N$ of $G_K$, call a Galois extension $M/K$ \emph{disjoint
from $N$} if it is solvable and $G_K=G_MN$. We use the following results of the development.
\begin{enumerate}[label=(\roman*),nosep]
\item (\lref{GlobalTheory.T.exists_solvable_cm_extension_split_place}) For a CM field $K$, a place $A$ and an open
  normal subgroup $N$ of $G_K$ there is a CM extension $M_0/K$ disjoint from $N$ such that $D_A$ fixes $M_0$ and the
  decomposition group of $A$ over $M_0^+$ fixes $M_0$.
\item (\lref{GalTw.GlobalTheory.T.exists_cuspidal_charLift_of_solvable_cmExtension}; uses $\sfname{Inq}$) If $(\pi,\chi)$
  is in scope $\mathrm{all}$ over $K$, $r$ is weakly attached to $\pi$, $N$ is an open normal subgroup of $G_K$ on which
  $\galredtr r$ does not vanish, and $M/K$ is a CM extension disjoint from $N$, then $\pi$ has a cuspidal character lift
  $P$ to $M$ with $(P,\chi')$ in scope $\mathrm{all}$ for some $\chi'$.
\item (\lref{GlobalTheory.T.exists_open_subgroup_weilElement_characterKernel}) For a local Artin map $a_0$ at $A$ for a
  field $K$, finitely many smooth characters $\zeta_j$ of $K_{v}^\times$ ($v$ below $A$) and $g\in W_A(K)$, there is an
  open subgroup $V$ of $D_A(K)$ containing $g$ and the closure of the commutator subgroup, with $\zeta_j(a_0(t))=1$ for
  all $j$ and $t\in V\cap I_A(K)$.
\item (\lref{GlobalTheory.T.exists_solvable_cm_extension_abelianCompletion}) For a CM field $K$, an open normal
  subgroup $N$ of $G_K$ and an open subgroup $V$ of $D_A(K)$ containing the commutator subgroup, there is a CM extension
  $M/K$ disjoint from $N$ such that an element of $D_A(K)$ fixes $M$ if and only if it lies in $V$.
\item (\lref{GlobalTheory.T.cuspRep_characterSupportLift_trans_tower}) Character lifts compose along towers.
\end{enumerate}
Since $\galredtr r$ is continuous on the profinite group $G_E$ and equals $n\ne0$ at $1$, there is an open normal
subgroup $N_0$ of $G_E$ on which $\galredtr r$ does not vanish. Take $M_0$ by (i) for $N_0$, and a character lift $\pi_0$
of $\pi$ to $M_0$ by (ii). Let $v_0$ be the place of $M_0$ below $A$; since $D_A$ fixes $M_0$, we have $W_A(M_0)=W_A$, and
we choose a local Artin map $a_0$ at $A$ for $M_0$ (\lref{GlobalTheory.localArtinExists_holds}). Apply (iii) to
$a_0$, the characters $\zeta_j=\xi_j\circ\Nm_{(M_0)_{v_0}/E_w}$ and $g$, giving $V$. The restriction $r_0=r|_{G_{M_0}}$ is
weakly attached to $\pi_0$ (\lref{GlobalTheory.GalCut.weak_comp_res}); choose, as before, an open normal subgroup
$N_1$ of $G_{M_0}$ on which $\galredtr r_0$ does not vanish. Take $M$ by (iv) for $M_0$, $N_1$ and $V$, a character lift
$P$ of $\pi_0$ to $M$ by (ii), which is a character lift of $\pi$ by (v). Since $g\in V$, $g$ fixes $M$. If $t\in W_A$ has
degree $0$ and fixes $M$, then $t\in D_A(M_0)$ fixes $M$, so $t\in V$, and $t\in I_A(M_0)$; local Artin maps are
compatible with norms (\lref{GlobalTheory.T.localArtin_comp_norm_compat}), so
$\xi_j(a(t))=\xi_j(\Nm(a_0(t)))=\zeta_j(a_0(t))=1$.
\end{proof}

\subsubsection{Unramified abelian CM extensions}

\begin{definition}[Towers adapted to a set of places]\label{pf:galred:tower}
Let $E$ be a CM field and $S$ a finite set of finite places of $E$. A CM field $X\subset\Omega$ abelian over $E$ is
\emph{reached by a tower adapted to $S$} if there is a chain $E=X_0\subset X_1\subset\dots\subset X_k=X$ of CM fields
abelian over $E$ such that each $X_{i+1}/X_i$ is Galois of prime degree and there are a place $u\notin S$ of $E$ which
splits completely in $X_i$ and a place $A$ of $\Omega$ above $u$ whose inertia group in $G_{X_i}$ does not fix $X_{i+1}$.
\formallabel{UpperTheory.Gal.UnramifiedFromSplittingCriterion.Tower}
\end{definition}

\begin{theorem}\label{pf:galred:unrfield}
Let $E$ be a CM field, $w$ a split place of $E$, $f\ge1$ and $S$ a finite set of finite places of $E$. There is a CM field
$X\subset\Omega$, abelian over $E$ and reached by a tower adapted to $S$, such that for every place $A$ of $\Omega$ above
$w$ the inertia group $I_A$ fixes $X$ and the place of $X$ below $A$ has norm $q_w^f$, and for every non-split place $u$
of $E$ and every $A$ above $u$ the inertia group $I_A$ fixes $X$.
\formallabel{UpperTheory.Gal.UnramifiedFromSplittingCriterion.unrField_holds}
\end{theorem}

\begin{proof}
Let $F=E^+$ and $v$ the place of $F$ below $w$. A subgroup $H$ of $G_F$ is \emph{good} if it is open, contains the
commutator subgroup and contains every complex conjugation (every $c\in G_F$ acting as complex conjugation through some
embedding $\Omega\to\C$); it is \emph{split-ramified} if for every place $t$ of $F$ and every $A$ above $t$, either the
inertia group $I_A(F)$ lies in $H$ or $D_A(F)$ fixes $E$. Good and split-ramified subgroups are stable under finite
intersections and under enlargement. For good $H$ the field $X_H$ fixed by $\{x\in G_E:x|\in H\}$ is a CM field abelian
over $E$ (\lref{UpperTheory.Gal.UnramifiedFromSplittingCriterion.exFld}). Recall that $w$ is split if and only if
$D_A(F)$ fixes $E$ for one, or every, $A$ above $w$, since the decomposition group of $w$ in $\Gal(E/F)$ is the image of
$D_A(F)$. We use:
\begin{enumerate}[label=(\roman*),nosep]
\item (\lref{UpperTheory.Gal.UnramifiedFromSplittingCriterion.unrCharPP}, built on
  \lref{UpperTheory.Gal.UnramifiedFromSplittingCriterion.ordPrime}) For every prime power $f$ there is a character
  $r\colon G_F\to\C^\times$ with good and split-ramified kernel such that, for every $A$ above $v$, $I_A(F)\subset\ker r$
  and $r(\phi)$ has order exactly $f$ for every arithmetic Frobenius $\phi$ at $A$.
\item (\lref{UpperTheory.Gal.UnramifiedFromSplittingCriterion.freshChar_holds}) For every open subgroup $U$ of $G_F$,
  finite set $T$ of places of $F$ and $m\ge2$, there are a character $\psi$ of $G_F$ with good and split-ramified kernel and
  $\psi^m=1$, trivial on the decomposition groups at the places of $T$, and a place $u\notin T$ of $F$ whose decomposition
  groups lie in $U$ and such that $\psi(\tau)$ has order $m$ for some inertia element $\tau$ at each $A$ above $u$.
\item (\lref{UpperTheory.Gal.UnramifiedFromSplittingCriterion.layerOf_holds}) Let $H'\subset H$ be good subgroups with
  $H'$ of prime index in $H$, and $u$ a place of $F$ not below $S$ whose decomposition groups fix $E$ and lie in $H$, such
  that some inertia element at some $A$ above $u$ is not in $H'$. Then $X_{H'}/X_H$ is a step of a tower adapted to $S$.
\item (\lref{UpperTheory.Gal.UnramifiedFromSplittingCriterion.groupStep}) An abstract lemma on groups: let $\Gamma$
  be a group with a class of good subgroups stable under finite intersections and enlargement, a property $\mathcal P$ of
  subgroups, a subset $D$, and families $(D_u)$, $(I_u)$, $I_u\subset D_u$, of subsets indexed by a set of ``places''.
  Assume that for every good $U$ and $m\ge1$ there are a character $\psi$ with good kernel, $\psi^m=1$, $\psi|_D=1$, and a
  place $u$ with $D_u\subset U$ and $\psi(\tau)$ of order $m$ for some $\tau\in I_u$; and that $\mathcal P(H)$ implies
  $\mathcal P(H')$ for good $H'\subset H$ of prime index such that some place $u$ has $D_u\subset H$ and $I_u\not\subset H'$.
  Then for every good $H_0$ with $\mathcal P(H_0)$ and every character $r$ of finite order with good kernel there is a
  character $r'$ with good kernel, $r'=r$ on $D$, and $\mathcal P(H_0\cap\ker r')$.
\item Let $K\subset M$ be number fields, $A$ a place
  above places $v$ of $K$ and $u$ of $M$ such that $I_A(K)$ fixes $M$, $\phi$ an arithmetic Frobenius at $A$ for $K$ and
  $d\ge1$ such that $\phi^k$ fixes $M$ if and only if $d\mid k$. Then $q_u=q_v^d$. (Indeed $I_A(M)=I_A(K)$, and the image of
  the open subgroup $D_A(M)=D_A(K)\cap G_M$ in $D_A(K)/I_A(K)\cong\hat\Z$, topologically generated by $\phi$, is the open
  subgroup $e\hat\Z$ with $e$ the residue degree of $u$ over $v$; it contains $\phi^k$ exactly when $e\mid k$, so $e=d$.)
\end{enumerate}
\emph{Step 1.} There is a character $r$ as in (i) for $f$ itself: write $f$ as a product of prime powers, take characters
$r_p$ by (i) and let $r=\prod r_p$; its kernel contains $\bigcap\ker r_p$, so it is good and split-ramified, it contains
$I_A(F)$, and $r(\phi)$ has order $f$ because the $r_p(\phi)$ commute and have pairwise coprime orders.

\emph{Step 2.} Apply (iv) with $\Gamma=G_F$, the good and split-ramified subgroups as good subgroups, $\mathcal P(H)$:
``$X_H$ is reached by a tower adapted to $S$'', $D$ the union of the decomposition groups $D_A(F)$ at the places $A$ above
$v$, the places $u$ of $F$ not below $S$ whose decomposition groups fix $E$, with $D_u$, $I_u$ the unions of the
decomposition and inertia groups at the places above $u$, $H_0=G_F$ (so $X_{H_0}=E$) and $r$ from Step 1 (of finite
order, as its kernel is open). The first hypothesis of (iv) holds by (ii) with $T=\{v\}\cup\{\text{places below }S\}$ (for
$m=1$ take $\psi=1$ and the place given by (ii) for $m=2$): the place $u$ of (ii) has decomposition groups fixing $E$, because $\psi$ is nontrivial on inertia
at $u$ and $\ker\psi$ is split-ramified. The second holds by (iii). We obtain $r'$ with good and split-ramified kernel,
$r'=r$ on the decomposition groups above $v$, such that $X=X_{\ker r'}$ is reached by a tower adapted to $S$; an element
$x\in G_E$ fixes $X$ if and only if $r'(x|_F)=1$.

\emph{Step 3.} Let $A$ lie above $w$. For $t\in I_A$, $t|_F\in I_A(F)\subset D_A(F)$, so $r'(t|_F)=r(t|_F)=1$ and $t$ fixes
$X$. As $w$ is split, $D_A(F)$ fixes $E$, so by (v) (with $d=1$) $q_w=q_v$; hence an arithmetic Frobenius $\phi$ at $A$
for $E$ is one for $F$, and $\phi^k$ fixes $X$ if and only if $r(\phi)^k=1$, i.e.\ $f\mid k$. By (v), the place of $X$
below $A$ has norm $q_w^f$. Finally let $u$ be non-split, $A$ above $u$ and $t\in I_A$. As $\ker r'$ is split-ramified,
either $I_A(F)\subset\ker r'$, and then $t$ fixes $X$, or $D_A(F)$ fixes $E$, which would make $u$ split.
\end{proof}

\subsubsection{Supports under unramified base change}

For a totally complex field $K\subset\Omega$ and a cuspidal $\pi$ on $\GL_n(\A_K)$, a cuspidal $P$ on $\GL_n(\A_X)$,
$K\subset X\subset\Omega$, is a \emph{weak base change of $\pi$ along a chain of prime degree} if there are totally
complex fields $K=X_0\subset\dots\subset X_k=X$ with $X_{i+1}/X_i$ Galois of prime degree and cuspidal $P_i$ on
$\GL_n(\A_{X_i})$, $P_0=\pi$, $P_k=P$, such that the Satake parameters of $P_{i+1}$ at almost all places are those of
$P_i$, raised to the power of the residue degree. The lifts of (iii) in \cref{sec:pf:galred:R2} are of this kind.

\begin{theorem}[Supports under unramified base change]\label{pf:galred:suppbc}
Let $E$ be a totally complex number field, $\pi$ cuspidal on $\GL_n(\A_E)$, $n\ge1$, and $P$ on $\GL_n(\A_X)$ a weak base
change of $\pi$ along a chain of prime degree. Let $A$ be a place of $\Omega$ above places $v$ of $E$ and $u$ of $X$,
such that $E_v\cong\Q_g$ continuously for a prime $g$ and $I_A$ fixes $X$. Let $\tau,\tau'$ be local components of $\pi$
at $v$ and of $P$ at $u$, with $\tau$ of depth zero, and let $f$ satisfy $q_u=g^f$ and be divisible by the dimension of
every member of every supercuspidal support of $\tau$. Then there are tamely ramified characters $\xi_1,\dots,\xi_n$ of
$X_u^\times$ with $\supp(\tau')=\{\xi_1,\dots,\xi_n\}$ and
\[
  P_{\galredrc(\tau)}(t)=\prod_j\bigl(X-\xi_j(a'(t))\bigr)\qquad(t\in I_A)
\]
for every local parameter map $\galredrc$ of $E_v$ and every local Artin map $a'$ of $X_u$ at $A$. If $\tau$ is mixed, then
$\{\xi_1,\dots,\xi_n\}$ is unlinked.
\formallabel{UpperTheory.Gal.InertialTypesAlt.suppBC_outright}
\end{theorem}

\begin{proof}
For a nonarchimedean local field $k$ we write $\tau^{K_1}$ for the representation of $\GL_n(k_k)$ on the vectors of a
representation $\tau$ of $\GL_n(k)$ fixed by $K_1=\ker(\GL_n(\OO_k)\to\GL_n(k_k))$ ($k_k$ the residue field). If $k'/k$ is
finite with residue fields $\F_q\subset\F_{q'}$, Shintani's norm map relates conjugacy classes of $\GL_n(\F_{q'})$ up to
Frobenius-twisted conjugacy with conjugacy classes of $\GL_n(\F_q)$, and gives the \emph{Shintani lift} of class functions
from $\GL_n(\F_q)$ to $\GL_n(\F_{q'})$ \cite{Shintani1976}. The proof combines the following results of the development.
\begin{enumerate}[label=(\roman*),nosep]
\item (\lref{UpperTheory.Gal.InertialTypesAlt.L.inflMatch}) \emph{Matching for inflated functions.} Let $k'/k$ be an
  unramified extension of $p$-adic fields of characteristic $0$ with Frobenius $s$, and $f_b$ a class function on
  $\GL_n(k_k)$. For $\delta\in\GL_n(k')$ with regular semisimple norm $\delta s(\delta)\cdots s^{[k':k]-1}(\delta)$ and
  $\gamma\in\GL_n(k)$ with the same characteristic polynomial, the $s$-twisted orbital integral at $\delta$ of the
  inflation to $\GL_n(\OO_{k'})$ of the Shintani lift of $f_b$ equals the orbital integral at $\gamma$ of the inflation of
  $f_b$; and a regular semisimple $\gamma$ at which the latter integral is nonzero has the characteristic polynomial of a
  norm.
\item (\lref{UpperTheory.Gal.InertialTypesAlt.L.finTrace_of_matching}) \emph{Transfer of level-zero traces.} From
  (i) and the comparison of trace formulas underlying cyclic base change: in a step $X_{i+1}/X_i$ of prime degree, at a
  place which is inert and unramified, if a class function $f$ on $\GL_n$ of the residue field satisfies
  $\sum_\gamma f(\gamma)\galredtr\tau_i^{K_1}(\gamma)\ne0$ for the local component $\tau_i$ of $P_i$, then the Shintani lift
  of $f$ acts on $\tau_{i+1}^{K_1}$ by a nonzero operator.
\item (\lref{UpperTheory.Gal.InertialTypesAlt.suppAny_of}) \emph{Toral supports.} If (ii) holds, then in the
  situation of the theorem there are tame characters $\xi_j$ with $\supp(\tau')=\{\xi_j\}$ and the displayed identity on
  inertia. The proof follows the level-zero types of the local components along the chain (at the places above $v$ every
  step is split or inert and unramified, as $I_A$ fixes $X$) by (ii), using Green's parametrization of the irreducible
  characters of $\GL_n(\F_q)$ \cite{Green1955}, Shintani's description of the lifts of cuspidal characters (the lift of
  a cuspidal character of $\GL_m(\F_q)$ to $\GL_m(\F_{q^f})$, $m\mid f$, is orthogonal to all products of cuspidal
  characters of proper Levi subgroups, \lref{UpperTheory.Gal.InertialTypesAlt.killsAll_holds}), and the type law
  (L12) for the local parameter maps.
\item (\lref{UpperTheory.Gal.InertialTypesAlt.mixType_of}) \emph{Types in the mixed case.} If (ii) holds and $\tau$ is
  mixed, then $\tau'$ contains the character $\theta_1\otimes\dots\otimes\theta_c$ (through determinants of the blocks) of
  the parahoric subgroup attached to a decomposition $n=d_1+\dots+d_c$, for pairwise distinct characters $\theta_b$ of
  $k_u^\times$.
\item (\lref{UpperTheory.Gal.TypesAndIwahori.unlinked_of_typeVec}) A generic irreducible representation containing such
  a type with pairwise distinct $\theta_b$ has unlinked supercuspidal support.
\end{enumerate}
By (i) and (ii), the hypothesis of (iii) and (iv) holds. Then (iii) gives the characters $\xi_j$ and the identity on
inertia. If $\tau$ is mixed, $\tau'$ contains a type as in (iv); being a local component of a cuspidal representation,
$\tau'$ is generic (\lref{GlobalTheory.CuspRep.Generic.isGeneric_of_isLocComp}), so by (v) every support of $\tau'$, in
particular $\{\xi_j\}$, is unlinked.
\end{proof}

\subsubsection{Counting depth-zero cuspidal representations}

\begin{proposition}\label{pf:galred:dzcount}
Let $F$ be a nonarchimedean local field with residue field $k_F$, $n\ge1$, $\varpi$ a uniformizer of $F$ and
$c\in\C^\times$, and let $\mathcal L$ be the set of monic irreducible polynomials of degree $n$ in $k_F[X]$ other than
$X$. There is a map $Y$ from $\mathcal L$ to $\galredIrr(F)$, injective, such that every $Y(P)$ is a cuspidal
representation of $\GL_n(F)$ with nonzero vectors fixed by $K_1=\ker(\GL_n(\OO_F)\to\GL_n(k_F))$, whose central character
$\omega$ satisfies $\omega(\varpi)=c$.
\formallabel{PositivityTheory.Rg.DzCount.depth_zero_count_holds}
\end{proposition}

\begin{proof}
Representations of $\GL_n(k_F)$ are regarded as representations of $\GL_n(\OO_F)$ through reduction. We use:
\begin{enumerate}[label=(\alph*),nosep]
\item (\lref{PositivityTheory.Rg.DzCount.st7}) There are irreducible cuspidal representations $\tau_P$ of
  $\GL_n(k_F)$, $P\in\mathcal L$, which are pairwise non-isomorphic (Green's parametrization \cite{Green1955}).
\item (\lref{PositivityTheory.Rg.DzCount.st4}) For every irreducible representation $\tau$ of $\GL_n(k_F)$ there are an
  irreducible smooth representation $\sigma$ of $\GL_n(F)$ with a central character $\omega$, $\omega(\varpi)=c$, and an
  injective $\GL_n(\OO_F)$-equivariant map $j\colon\tau\to\sigma$.
\item (\lref{PositivityTheory.Rg.DzCount.st5}) If moreover $\tau$ is cuspidal, then $\sigma$ is cuspidal
  \cite{MoyPrasad1996}.
\item (\lref{PositivityTheory.Rg.DzCount.B.st6}) If $\tau$ is cuspidal and $\sigma$, with a central character,
  contains $\tau$ through $j$, then the space of $K_1$-fixed vectors of $\sigma$ is $j(\tau)$.
\end{enumerate}
For $P\in\mathcal L$ choose $\sigma_P,\omega_P,j_P$ by (b) for $\tau_P$, and put $Y(P)=\sigma_P$. By (c), (b) and (d),
$Y(P)$ is cuspidal, has central character $\omega_P$ with $\omega_P(\varpi)=c$, and its $K_1$-fixed vectors form
$j_P(\tau_P)\ne0$. If $Y(P)\cong Y(P')$, an isomorphism maps $K_1$-fixed vectors onto $K_1$-fixed vectors,
$\GL_n(\OO_F)$-equivariantly, and by (d) it induces an isomorphism $\tau_P\cong\tau_{P'}$; so $P=P'$ by (a).
\end{proof}

\subsubsection{Proof of the division theorem}\label{sec:pf:galred:divproof}

We prove \cref{pf:galred:division}. Let the data be as there. For a subset $\Sigma\subset M_N(K)$, the polynomial functions
on $M_N(K)$ (polynomials in the entries, with coefficients in $K$) vanishing on $\Sigma$ form an ideal $I(\Sigma)$, and the
\emph{tangent space} of $\Sigma$ at $1$ is the space of $Y\in M_N(K)$ with $d_1f(Y)=0$ for all $f\in I(\Sigma)$. Let
$\mathfrak g$ be the tangent space of $\rho(\Gamma)$ at $1$ (the Lie algebra of the Zariski closure of $\rho(\Gamma)$) and
$\mathcal A$ the $K$-subalgebra of $M_N(K)$ generated by $\mathfrak g$. A matrix has \emph{degree at most $n$} if a
nonzero polynomial of degree at most $n$ annihilates it.

\emph{Step 1: an element of $\mathfrak g$ with lacunary eigenvalues.} The space $\mathfrak g$ is cut out in $M_N(K)$ by
finitely many linear forms $d_1f_1,\dots,d_1f_m$, $f_i\in I(\rho(\Gamma))$, and by hypothesis $Z\in M_N(K')$ satisfies
them. The conditions ``$X_0$ satisfies $d_1f_i(X_0)=0$, $\det(e_i-X_0)=0$ for all $i$, and $(e_i)$ is lacunary'' form a
finite system of polynomial equations and inequations with coefficients in $K$ in the unknowns $(X_0,e)$; it has the
solution $(Z,(e_i))$ over $K'$, hence a solution $(X_0,(e_{0,i}))$ over the algebraically closed field $K$ (Hilbert's
Nullstellensatz). So $X_0\in\mathfrak g$ has the lacunary eigenvalues $e_{0,1},\dots,e_{0,n}$; they are distinct, since
lacunarity with $A=\{i\}$, $B=\{j\}$ gives $e_{0,i}\ne e_{0,j}$.

\emph{Step 2: degrees.} Every $\rho(g)$ has degree at most $n$. Indeed, $y$ has degree at most $n$ if and only if
$1,y,\dots,y^n$ are linearly dependent, i.e.\ all $(n+1)$-minors of the matrix with columns $y^0,\dots,y^n$ (as vectors
of length $N^2$) vanish; this is a closed condition on $y$, it holds for $\rho(s)$, $s\in S$, because $P_s(\rho(s))=0$, and
$S$ is dense and $\rho$ continuous.

\emph{Step 3: idempotents.} (\lref{LateralTheory.Fin.Struct.intersectionNumber_baseChange}) Since $\rho$ is semisimple,
every $\rho(g)$ has degree at most $n$ and $X_0\in\mathfrak g$ has the $n$ lacunary eigenvalues $e_{0,i}$, there are
nonzero pairwise orthogonal idempotents $p_1,\dots,p_n$ with $\sum p_i=1$ such that $\sum_ie_{0,i}p_i\in\mathfrak g$, $p_i\in
\mathcal A$, $p_i\mathcal Ap_i\subset Kp_i$, and $p_iap_j\in\mathfrak g$ for $a\in\mathcal A$ and $i\ne j$. Moreover
$\mathcal A$ is normalized by $\rho(\Gamma)$ and preserves every $\rho(\Gamma)$-stable subspace of $K^N$.

\emph{Step 4: the key property.} (\lref{LateralTheory.Fin.Key.intersectionNumber_baseChange}) For every normal
subgroup $U$ of finite index of $\Gamma$ and every $g\in\Gamma$: (a) there is $m$ in the Zariski closure of $\rho(U)$ such
that $\rho(g)m$ has $n$ distinct eigenvalues; (b) for each $i$, if $p_ia\rho(g)p_i\ne0$ for some $a\in\mathcal A$, then
$p_i\rho(g)=p_iy$ for some $y\in\mathcal A$.

\emph{Step 5: an algebra containing the image.} (\lref{LateralTheory.Fin.BigCorner.intersectionNumber_baseChange})
By Step 3 and Step 4(b) there is a $K$-subalgebra $R\subset M_N(K)$ containing $\rho(\Gamma)$ and the $p_i$, with
$p_iRp_i\subset Kp_i$, normalized by $\rho(\Gamma)$ and preserving every $\rho(\Gamma)$-stable subspace.

\emph{Step 6: $R$ is semisimple.} Let $J$ be the set of $Y\in R$ such that $aY$ is nilpotent for all $a\in R$; it is the
radical of the finite-dimensional algebra $R$, a nilpotent two-sided ideal, and it is stable under conjugation by
$\rho(\Gamma)$. Hence $W=JK^N$ is $\rho(\Gamma)$-stable. As $\rho$ is semisimple, $K^N=W\oplus W'$ with $W'$
$\rho(\Gamma)$-stable, hence $R$-stable; then $JW'\subset W\cap W'=0$, so $JK^N=JW=J^2K^N$, and by induction
$JK^N=J^kK^N=0$ for large $k$. Thus $J=0$.

\emph{Step 7: matrix units.} (\lref{LateralTheory.Fin.MatrixUnits₂.intersectionNumber_baseChange}) As $p_iRp_i\subset
Kp_i$ and $R$ has no nonzero radical, there are $E_{ij}\in R$ with $E_{ii}=p_i$, $p_iE_{ij}p_j=E_{ij}$, $p_iap_j\in KE_{ij}$
for $a\in R$, $E_{ji}\ne0$ whenever $E_{ij}\ne0$, and $E_{ij}E_{jk}=E_{ik}\ne0$ whenever $E_{ij},E_{jk}\ne0$.

\emph{Step 8: the representation.} (\lref{LateralTheory.Fin.Rep.intersectionNumber_baseChange}) From the matrix units
one obtains a representation $r\colon\Gamma\to\GL_n(K)$ and matrices $M\in M_{N\times n}(K)$, $B\in M_{n\times N}(K)$ with
$BM=1$ and $\rho(g)M=Mr(g)$ for all $g$, such that the images of all $T\in M_{N\times n}(K)$ with $\rho(g)T=Tr(g)$ for all
$g$ span $K^N$.

\emph{Step 9: consequences.} $r(g)=B\rho(g)M$, so $r$ is continuous; $M$ has rank $n$; $r$ is isomorphic to a
subrepresentation of $\rho$, hence semisimple. A polynomial $q$ annihilates $\rho(g)$ if and only if it annihilates
$r(g)$: if $q(\rho(g))=0$ then $Mq(r(g))=q(\rho(g))M=0$, so $q(r(g))=0$; if $q(r(g))=0$ then $q(\rho(g))T=Tq(r(g))=0$ for
all intertwiners $T$ as in Step 8, whose images span $K^N$, so $q(\rho(g))=0$. In particular $P_s(r(s))=0$ for $s\in S$, and
the characteristic polynomial of $r(g)$, which annihilates $r(g)$, annihilates $\rho(g)$.

\emph{Step 10: separability.} (\lref{LateralTheory.Fin.DenseAlg.intersectionNumber_baseChange}) By Step 4(a), for
every normal subgroup $U$ of finite index and every $g$ there is $u\in U$ with $\det(X-r(gu))$ separable; and then
(\lref{LateralTheory.Fin.FinDense.intersectionNumber_baseChange}) for a continuous representation of a compact group
with this property, the set of $g$ with separable characteristic polynomial is open and dense.

\emph{Step 11: characteristic polynomials and determinant.} If a monic $P$ of degree $n$ annihilates $x\in M_n(K)$ and
$\det(X-x)$ is separable, then $\det(X-x)=P$: the minimal polynomial of $x$ has the $n$ distinct eigenvalues as roots, so
it is $\det(X-x)$, and it divides $P$. Hence $\det(X-r(s))=P_s$ for $s\in S$ with $\det(X-r(s))$ separable. The set of
such $s$ is dense, as the intersection of the dense set $S$ with an open dense set, and there $\det r(s)=(-1)^nP_s(0)=
\psi(s)$; by continuity $\det r=\psi$. If a monic $P$ of degree $n$ is product-generic, annihilates $x$ and
$\det x=(-1)^nP(0)$, then the eigenvalues of $x$ are $n$ roots of $P$ with the product of all roots of $P$, so they are the
roots of $P$ and $\det(X-x)=P$; this applies to $x=r(s)$, $s\in S$, as $\det r(s)=\psi(s)=(-1)^nP_s(0)$.

\emph{Step 12: covering.} (\lref{LateralTheory.Fin.FinEmbed.intersectionNumber_baseChange}) As $\rho$ and $r$ are
semisimple and the intertwiners from $r$ to $\rho$ span, there are $T_1,\dots,T_k\in M_{n\times N}(K)$ with
$T_j\rho(g)=r(g)T_j$ and $T'_1,\dots,T'_k\in M_{N\times n}(K)$ with $\rho(g)T'_j=T'_jr(g)$, such that $\sum_jT'_jT_j=1$. The
$(nk)\times N$ matrix $M'$ obtained by stacking the $T_j$ satisfies $(r(g)\otimes1_k)M'=M'\rho(g)$ and has rank $N$, since
$\sum_jT'_jT_j=1$. So $r$ is inside $\rho$ (Step 8) and covers $\rho$. Together with Steps 9--11 this proves
\cref{pf:galred:division}.

\subsubsection{Density of generic Satake parameters}

Recall from \cref{sec:pf:galred:cong} the spaces $\mathcal A(W,K)$ and $\mathcal A_{\mathcal P}(W,K)$ and the notion of
density. A level $K$ is \emph{away from} $w_0$ if it is an open compact subgroup of $\galredU_a(\A_{E^+,f})$ of the type
used there; for a split place $w$, $K$ is \emph{hyperspecial at $w$} if its component at $w$ is $\GL_a(\OO_{E_w})$.

\begin{theorem}\label{pf:galred:densegen}
Let $E$ be a CM field, $\ell$ a prime, $\iota\colon\Qlbar\to\C$, $a\ge1$, $w_0\mid\ell$ a split place, $K$ a level away
from $w_0$, $W$ a weight, $S$ a finite set of split places apart from $w_0$ at which $K$ is hyperspecial, and $\Phi_w\in
\C[X_1,\dots,X_a]$, $w\in S$, nonzero polynomials. Then the property ``$\Phi_w(z)\ne0$ for every $w\in S$ and every Satake
parameter $z$ of $U$ at $w$'' is dense at $(W,K)$.
\formallabel{UpperTheory.Gal.Gen.denseGen_holds}
\end{theorem}

\begin{proof}
Let $\mathcal P$ be the property in question; by definition $\mathcal A_{\mathcal P}(W,K)$ is a subspace of
$\mathcal A(W,K)$. The main input is the following step
(\lref{UpperTheory.Gal.Gen.step_holds}): for every $f\in\mathcal A(W,K)$ there are $G\in\mathcal A_{\mathcal P}(W,K)$
and a function $e$ on $\galredU_a(\A_{E^+})$ with $|\iota^{-1}e(g)|^2\le|\ell|$ for all $g$ and $f-G=e\cdot f$. Its proof
uses global Hecke operators built from local ones at the places of $S$ and a Plancherel-type count of generic Satake
parameters (\lref{UpperTheory.Gal.Gen.globOp_holds} and \lref{UpperTheory.Gal.Gen.planch_holds}).

We show by induction on $m$: for every $f\in\mathcal A(W,K)$ there are $G\in\mathcal A_{\mathcal P}(W,K)$ and $e$ with
$|\iota^{-1}e(g)|\le|\ell|^{m/2}$ and $f-G=e\cdot f$. For $m=0$ take $G=0$, $e=1$. Given $G,e$ for $m$, apply the step to
$f-G\in\mathcal A(W,K)$: there are $G'\in\mathcal A_{\mathcal P}(W,K)$ and $e'$ with $|\iota^{-1}e'|\le|\ell|^{1/2}$
and $(f-G)-G'=e'(f-G)=e'e f$. Then $G+G'\in\mathcal A_{\mathcal P}(W,K)$ and $e'e$ has the bound for $m+1$. Given
$\varepsilon>0$, choose $m$ with $|\ell|^{m/2}<\varepsilon$.
\end{proof}

\subsubsection{Lacunary weight families}

We use the notation of \cref{sec:pf:galred:cong}. For an embedding $\tau\colon E\to\C$ and $x\in\galredU_a(E^+\otimes\R)$ let
$\tau(x)\in\galredU_a(\C)\subset\GL_a(\C)$ be the component of $x$ at $\tau$. A weight $W$ is an algebraic representation,
on a space $W_{\mathrm{sp}}$, of the archimedean group; an algebraic automorphic form of weight $W$ is a linear map $F$
from $W_{\mathrm{sp}}$ to functions on $\galredU_a(\A_{E^+})$, left invariant under $\galredU_a(E^+)$, with
$F(v)(gx)=F(W(x)v)(g)$ for archimedean $x$. For $k\le a$ and $z\in\C$
with $|z|=1$ let $t_k(z)=\mathrm{diag}(z,\dots,z,1,\dots,1)$ ($k$ entries $z$); an \emph{archimedean circle} at $\tau$ is a
map $z\mapsto x(z)$ into the archimedean group with $\tau(x(z))=t_k(z)$. For a dominant weight $\nu$ of $\GL_a$ let
$\chi_\nu$ be the corresponding character of the upper triangular Borel subgroup of $\GL_a(\C)$ and
$|\nu|_k=\nu_0+\dots+\nu_{k-1}$. A function $\phi$ on the unit circle has \emph{exponents $\le s$} (resp.\ $\ge s$) if it is
a finite sum of monomials $cz^j$ with $j\le s$ (resp.\ $j\ge s$).

\begin{definition}[Families]\label{pf:galred:fam}
Let $\tau\colon E\to\C$, $W$ a weight, $K$ a level, $\nu$ a dominant weight of $\GL_a$ and $d\ge0$. A \emph{$\nu$-family
of defect $d$ at $\tau$ and level $K$} is a map $B$ assigning to $y_1,y_2\in\GL_a(\C)$ a linear map $B(y_1,y_2)$ from
$W_{\mathrm{sp}}$ to functions on $\galredU_a(\A_{E^+})$ such that: (F1) all $B(y_1,y_2)(v)$ are automorphic forms, right
invariant under $K$, lying in a fixed finite-dimensional space; (F2) for archimedean $x$, the right translate of
$B(y_1,y_2)(v)$ by $x$ is $B(\tau(x)y_1,\tau(x)y_2)(W(x)v)$; (F3) $B(y_1b,y_2)=\chi_\nu(b)B(y_1,y_2)$ for upper triangular
$b$; (F4) for every archimedean circle $x(z)$ at $\tau$ for $t_k$, every $v$ and $g$, the function $z\mapsto
B(t_k(z)y_1,t_k(z)y_2)(W(x(z))v)(g)$ has exponents $\le|\nu|_k+d$, and $z\mapsto B(1,t_k(z)y_2)(W(x(z))v)(g)$ has exponents
$\ge-d$.
\formallabel{UpperTheory.Gal.Wt.Fam}
\end{definition}

\begin{proposition}\label{pf:galred:famex}
Let $E$ be a CM field, $\ell$, $\iota$, $a\ge1$, $w_0\mid\ell$ and $v'$ split places which are apart, $K$ a level away from
$w_0$ whose component at $v'$ contains $\GL_a(\OO_{E_{v'}})$, $K_{v'}$ an open subgroup of $\GL_a(E_{v'})$ and $K^\circ$
the subgroup of elements of $K$ with component at $v'$ in $K_{v'}$. Let $f\in\mathcal A(W,K^\circ)$, $\tau$ an embedding
lying over $v'$, $\mathcal X$ a set and $b\colon\galredU_a(\A_{E^+})\times\mathcal X\to W_{\mathrm{sp}}$ such that for every
algebraic form $F$ of weight $W$ the values $\iota^{-1}F(b(g,i))(g)$ are bounded. For every $\varepsilon>0$ there are $d\ge0$
and $M\colon\galredU_a(\A_{E^+})\to\GL_a(\C)$, invariant under right translation by finite adeles with trivial component at
$v'$, such that for every dominant $\nu$ there is a $\nu$-family $B$ of defect $d$ at $\tau$ and level $K$ with
\[
  \bigl|\iota^{-1}\bigl(f(b(g,i))(g)-B(M(g),M(g))(b(g,i))(g)\bigr)\bigr|\le\varepsilon\qquad\text{for all }g,i.
\]
\formallabel{UpperTheory.Gal.Wt.famEx_holds}
\end{proposition}

\begin{proof}
For a dominant $\nu$ let $\delta_\nu(h)$ be the highest weight matrix coefficient of $h\in\GL_a(\C)$ on the irreducible
representation of highest weight $\nu$. We use:
\begin{enumerate}[label=(\alph*),nosep]
\item (\lref{UpperTheory.Gal.Wt.neatLev_holds}) There is a neat level $C\supset K$ (if $\gamma g=gcx$ with $\gamma$
  rational, $c\in C$, $x$ archimedean, then $c=1$) whose component at $v'$ lies in $\GL_a(\OO_{E_{v'}})$ and such that $f$ is
  fixed by $C\cap K^\circ$.
\item (\lref{UpperTheory.Gal.Wt.ellC_holds}) As $C$ is neat, there is $\lambda\colon\galredU_a(\A_{E^+})\to C$ with
  $\lambda(\gamma g)=\lambda(g)$, $\lambda(gx)=\lambda(g)$ and $\lambda(gc)=\lambda(g)c$ for rational $\gamma$, archimedean
  $x$ and $c\in C$.
\item (\lref{UpperTheory.Gal.Wt.twist_holds} and \lref{UpperTheory.Gal.Wt.twistBd_holds}) For $c\in C$ the function
  $F_c(v)(g)=f(v)(g\lambda(g)^{-1}c)$ is an algebraic form of weight $W$ and level $C$, and there is $R>0$ with
  $|\iota^{-1}F_c(b(g,i))(g)|\le R$ for all $c,g,i$.
\item (\lref{UpperTheory.Gal.Wt.ladic_weight_embedding_over_place}) There is a homomorphism
  $\kappa\colon\galredU_a(\A_{E^+})\to\GL_a(\C)$ with $\kappa(\gamma)=\tau(\gamma)$ for rational $\gamma$, trivial on the
  archimedean group, with $\kappa(g)=1$ if and only if $g_{v'}=1$, such that $\iota^{-1}\circ\kappa$ is continuous and
  $\iota^{-1}\kappa(g)$ is integral when $g_{v'}\in\GL_a(\OO_{E_{v'}})$ (it is $g\mapsto\iota(\tau'(g_{v'}))$ for the
  embedding $\tau'$ of $E_{v'}$ inducing $\tau$).
\item (\lref{UpperTheory.Gal.Wt.pieces_holds}) For the continuous integral homomorphism $\phi=\kappa|_C$, the open
  subgroup $C'=C\cap K^\circ\supset\ker\phi$ and $\eta>0$ there are $c_1,\dots,c_Q\in C$ and polynomials $P_1,\dots,P_Q$ in
  the matrix entries such that for every $c\in C$ there is $q_0$ with $c_{q_0}^{-1}c\in C'$,
  $|\iota^{-1}(P_{q_0}(\phi(c))-1)|\le\eta$, $|\iota^{-1}P_q(\phi(c))|\le\eta$ for $q\ne q_0$, and
  $|\iota^{-1}(\phi(c_{q_0}^{-1}c)-1)|\le\eta$ entrywise.
\item (\lref{UpperTheory.Gal.Wt.deltaNear_holds}) If $\iota^{-1}h$ is integral then $|\iota^{-1}\delta_\nu(h)|\le1$, and if
  moreover $|\iota^{-1}(h-1)|\le\eta$ entrywise then $|\iota^{-1}(\delta_\nu(h)-1)|\le\eta$.
\item (\lref{UpperTheory.Gal.Wt.famOf_holds}) Let $\mathfrak n\colon\galredU_a(\A_{E^+})\to\GL_a(\C)$ satisfy
  $\mathfrak n(\gamma g)=\mathfrak n(g)$, $\mathfrak n(gc)=\mathfrak n(g)$ and $\mathfrak n(gx)=\mathfrak n(g)\tau(x)$ for
  rational $\gamma$, $c\in C$ and archimedean $x$. For forms $F_1,\dots,F_Q$ of weight $W$ and level $C$ and polynomials
  $P_q$ there is $d$ such that for all $\nu$ and $X_1,\dots,X_Q\in\GL_a(\C)$ the map
  $B(y_1,y_2)(v)(g)=\sum_qP_q(\mathfrak n(g)y_2)\,\delta_\nu(X_q\mathfrak n(g)y_1)\,F_q(v)(g)$ is a $\nu$-family of defect
  $d$ at $\tau$ and level $C$.
\end{enumerate}
Take $C,\lambda,R$ by (a)--(c), $\kappa$ by (d), and $c_q,P_q$ by (e) with $\eta=\varepsilon/R$; put $F_q=F_{c_q}$. Let
$M(g)=\tau(g_\infty)^{-1}\kappa(g)$ and $\mathfrak n(g)=\phi(\lambda(g))M(g)^{-1}$. Then $\mathfrak n$ satisfies the
hypotheses of (g): for rational $\gamma$ the factors $\tau(\gamma)$ of $\kappa(\gamma g)$ and $\tau((\gamma g)_\infty)$
cancel; for $c\in C$, $\phi(\lambda(gc))=\phi(\lambda(g))\kappa(c)$ and $\kappa(gc)=\kappa(g)\kappa(c)$; for archimedean $x$,
$\lambda$ and $\kappa$ are unchanged and $\tau((gx)_\infty)=\tau(g_\infty)\tau(x)$. $M(gh)=M(g)$ for finite $h$ with
$h_{v'}=1$. Let $d$ be given by (g), and for $\nu$ let $B$ be the family of (g) with $X_q=\phi(c_q)^{-1}$; it has level
$C\supset K$. As $\mathfrak n(g)M(g)=\phi(\lambda(g))$,
\[
  B(M(g),M(g))(v)(g)=\sum_qP_q(\phi(\lambda(g)))\,\delta_\nu(\phi(c_q^{-1}\lambda(g)))\,F_q(v)(g).
\]
Fix $g$, let $q_0$ be given by (e) for $c=\lambda(g)$. Since $c_{q_0}^{-1}\lambda(g)\in C'$ fixes $f$,
$f(v)(g)=f(v)(g(c_{q_0}^{-1}\lambda(g))^{-1})=F_{q_0}(v)(g)$. Hence $f(v)(g)-B(M(g),M(g))(v)(g)=\sum_q\beta_qF_q(v)(g)$ with
$\beta_{q_0}=1-P_{q_0}\delta_\nu=-((P_{q_0}-1)\delta_\nu+(\delta_\nu-1))$ and $\beta_q=-P_q\delta_\nu$ for $q\ne q_0$
(evaluated at $\phi(\lambda(g))$, resp.\ $\phi(c_q^{-1}\lambda(g))$). By (e) and (f), $|\iota^{-1}\beta_q|\le\eta$ for all $q$;
with (c) and the ultrametric inequality the difference at $v=b(g,i)$ is at most $\eta R=\varepsilon$.
\end{proof}

\begin{theorem}[Lacunary weight families]\label{pf:galred:wtlac}
Let $E,\ell,\iota,a,w_0,v',K,K_{v'},K^\circ$ be as in \cref{pf:galred:famex}, $\lambda=(\lambda_\tau)$ a weight, and
$(t_j)_{j\in J}$ Hecke operators of level $K^\circ$, $\ell$-adically integral for $\lambda$, supported on elements whose
components at the places not apart from $w_0$ and at the places not apart from $v'$ are trivial, and acting on the
$K^\circ$-fixed vectors of every automorphic $U$ by scalars $\eta_j(U)$. Let $\tau_0$ lie over $v'$ and let $U$ have weight
$\lambda$ and nonzero $K^\circ$-fixed vectors. Then $(\iota^{-1}\eta_j(U))_j$ is polynomially approximable by the systems
$(\iota^{-1}\eta_j(U'))_j$ of the automorphic $U'$ with nonzero $K$-fixed vectors and weights $\mu$ such that
$\mu_\tau=\lambda_\tau$ unless $\tau$ lies over $v'$ or $c(v')$, and $J_{\mu_{\tau_0}}$ is gapped.
\formallabel{UpperTheory.Gal.Wt.wtLacEx_holds}
\end{theorem}

\begin{proof}
We use, besides \cref{pf:galred:famex}:
\begin{enumerate}[label=(\alph*),nosep]
\item (\lref{GlobalTheory.T.exists_algebraicForm_weight_level_smooth_at_place}) $U$ contains a nonzero algebraic form
  $f\in\mathcal A(\lambda,K^\circ)$.
\item (\lref{UpperTheory.Gal.Gen.Loc.frame}) There are a set $\mathcal X$ and $b\colon\galredU_a(\A_{E^+})\times
  \mathcal X\to W_{\mathrm{sp}}$ such that an algebraic form $F$ of weight $\lambda$ vanishing at all $(b(g,i),g)$ is zero,
  the values $\iota^{-1}F(b(g,i))(g)$ of every $F$ are bounded, and if they are bounded by $R'$ then so are those of
  $t_jF$ (this uses the integrality of the $t_j$).
\item (\lref{UpperTheory.Gal.Wt.famProj_holds}) A $\nu$-family $B$ of defect $d$ at $\tau_0$ and level $K$ is a finite sum
  $\sum_VF_V$ over automorphic representations $V$, where $F_V$ is a $\nu$-family of defect $d$ at $\tau_0$ and level $K$ with
  values in $V$.
\item (\lref{UpperTheory.Gal.Wt.famWt_holds}) If a $\nu$-family of defect $d$ at $\tau_0$ and level $K$ is nonzero with
  values in $V$, then $V$ has nonzero $K$-fixed vectors and a weight $\mu$ with $\mu_\tau=\lambda_\tau$ for
  $\tau\notin\{\tau_0,\bar\tau_0\}$ and $|\mu_{\tau_0,i}-\nu_i|\le2d$ for all $i$.
\item (\lref{GlobalTheory.T.eigenvector_small_sub_finite_eigen_approx_inequality}) Let operators $t_j$ on a space $\mathcal V$
  and linear forms $\mathrm{ev}_x$ ($x\in\mathcal X'$) be such that $|\iota^{-1}\mathrm{ev}_x(F)|\le R'$ for all $x$ implies the
  same for $t_jF$. Let $f_1$ be a common eigenvector with eigenvalues $\theta_j$ and $|\iota^{-1}\mathrm{ev}_{x_0}(f_1)|=1$ for
  some $x_0$, and $G_1,\dots,G_k$ common eigenvectors with eigenvalue systems $\eta^{(i)}$, such that
  $|\iota^{-1}\mathrm{ev}_x(f_1+\sum_iG_i)|\le\varepsilon$ for all $x$. Then there are $c_i$ with
  $|P(\iota^{-1}\theta)-\sum_ic_iP(\iota^{-1}\eta^{(i)})|\le\varepsilon$ for all $P\in\Z[X_j:j\in J]$.
\end{enumerate}
Let $\varepsilon>0$. Take $f$ by (a) and $\mathcal X,b$ by (b); some value $z_0=f(b(g_0,i_0))(g_0)$ is nonzero, and
$f_1=z_0^{-1}f$ has $|\iota^{-1}f_1(b(g_0,i_0))(g_0)|=1$. \Cref{pf:galred:famex} gives $d$ and $M$; let $\mu^0_k=-(a+1)^k$
($0\le k<a$, a dominant weight), $N=(a+1)(4d+a)+1$, $\nu=N\mu^0$, and $B$ the $\nu$-family with
$|\iota^{-1}(f_1-B(M(g),M(g)))|\le\varepsilon$ at all evaluation points. Decompose $B=\sum_VF_V$ by (c) and put
$G_V(v)(g)=-F_V(M(g),M(g))(v)(g)$, for the $V$ with $F_V\ne0$. As $t_j$ is supported on elements $h$ with $h_{v'}=1$, and
$M(gh)=M(g)$ for those, $t_jG_V=\eta_j(V)G_V$; also $t_jf_1=\eta_j(U)f_1$. So (e) applies with the evaluation forms
$F\mapsto F(b(g,i))(g)$, and gives the approximation of $(\iota^{-1}\eta_j(U))_j$ with error $\varepsilon$ by systems of
these $V$. By (d) each such $V$ has nonzero $K$-fixed vectors and a weight $\mu$ which agrees with $\lambda$ except at
$\tau_0$ and $\bar\tau_0$, which lie over $v'$ or $c(v')$, with $|\mu_{\tau_0,k}+N(a+1)^k|\le2d$.

It remains to see that $J_{\mu_{\tau_0}}$ is gapped. Write $\mu=\mu_{\tau_0}$ and $j_k=k-\mu_k$, which increase strictly in
$k$ as $\mu$ is dominant; put $m_k=(a+1)^k$, so $Nm_k-2d\le-\mu_k\le Nm_k+2d$. For $i+1<a$,
\[
  (j_{i+1}-j_0)-a(j_i-j_0)\ge N\bigl(m_{i+1}-1-a(m_i-1)\bigr)+\bigl(1-(a-1)i\bigr)-4d(a+1)\ge N-a^2-4d(a+1)>0,
\]
since $m_{i+1}-1-a(m_i-1)=m_i+a-1\ge1$, $(a-1)i\le a^2-1$ and $N>(a+1)(4d+a)$.
\end{proof}

\subsubsection{Pseudo-characters}

\begin{definition}[Pseudo-characters]\label{pf:galred:pschar}
Let $\Gamma$ be a monoid (or a $k$-algebra), $k$ a field and $T\colon\Gamma\to k$. Define functions
$\mathrm{alt}_m(T)$ of $m$ arguments by $\mathrm{alt}_0(T)=1$ and
\[
  \mathrm{alt}_{m+1}(T)(x_0,\dots,x_m)=T(x_0)\,\mathrm{alt}_m(T)(x_1,\dots,x_m)
  -\sum_{i=1}^m\mathrm{alt}_m(T)(x_1,\dots,x_0x_i,\dots,x_m).
\]
$T$ is a \emph{pseudo-character of dimension $d$} if $T(1)=d$, $T(xy)=T(yx)$ and $\mathrm{alt}_{d+1}(T)=0$.
\formallabel{GlobalTheory.PsChar.IsPsChar}
\end{definition}

Unwinding the
recursion, $\mathrm{alt}_m(T)(x_1,\dots,x_m)=\sum_{\sigma\in S_m}\mathrm{sgn}(\sigma)\prod_{(i_1\cdots i_r)}T(x_{i_1}\cdots x_{i_r})$,
the product over the cycles of $\sigma$; for the trace of a representation on a space $V$ this is $m!$ times the trace
of $x_1\otimes\dots\otimes x_m$ on $\bigwedge^mV$, so the trace of a $d$-dimensional representation is a pseudo-character
of dimension $d$ \cite{Taylor1991}. Each of the three conditions is a polynomial identity with integer coefficients in
finitely many values of $T$.

\begin{theorem}\label{pf:galred:ssrep}
Let $\Gamma$ be a group, $k$ an algebraically closed field of characteristic $0$ and $T\colon\Gamma\to k$ a
pseudo-character of dimension $d$. There is a semisimple representation $r\colon\Gamma\to\GL_d(k)$ with $\galredtr r=T$.
\formallabel{GlobalTheory.PsChar.exists_ssRep}
\end{theorem}

\begin{proof}
We use three results of the development, for a $k$-algebra $R$ and a $k$-linear pseudo-character $S\colon R\to k$ of
dimension $d$:
\begin{enumerate}[label=(\roman*),nosep]
\item (\lref{GlobalTheory.PsChar.cayleyHamilton}) \emph{Cayley--Hamilton.} For every $x\in R$ there is a monic $p\in
  k[X]$ of degree $d$ with $S(p(x)y)=0$ for all $y\in R$.
\item (\lref{GlobalTheory.T.pseudoCharacter_eq_zero_of_isNilpotent} and \lref{GlobalTheory.T.pseudoCharacter_apply_idempotent_nat_le})
  $S(x)=0$ for nilpotent $x$; for an idempotent $e$, $S(e)$ is a natural number at most $d$, and if $S(e)=0$ then $S(ey)=0$
  for all $y$.
\item (\lref{GlobalTheory.T.finite_and_isSemisimpleRing_of_no_nil_ideals_bounded_idempotents}) A $k$-algebra in which
  every element is algebraic, which has no nonzero $x$ with $yx$ nilpotent for all $y$, and in which sets of nonzero
  pairwise orthogonal idempotents have bounded size, is finite-dimensional and semisimple.
\end{enumerate}
Extend $T$ linearly to $T'\colon k[\Gamma]\to k$. Then $T'(xy)=T'(yx)$, and $\mathrm{alt}_{d+1}(T')=0$ since
$\mathrm{alt}_{d+1}$ is multilinear in its arguments (by induction on the recursion); so $T'$ is a pseudo-character of
dimension $d$. The set $I$ of $x$ with $T'(xy)=0$ for all $y$ is a two-sided ideal (as $T'(axy)=T'(xya)$); let
$R=k[\Gamma]/I$. $T'$ induces $S\colon R\to k$, again a pseudo-character of dimension $d$, whose trace form is
nondegenerate: $S(xy)=0$ for all $y$ implies $x=0$.

$R$ is finite-dimensional and semisimple by (iii): every $x$ is algebraic, as $p(x)=0$ for $p$ from (i) by nondegeneracy;
if $yx$ is nilpotent for all $y$, then $S(xy)=S(yx)=0$ for all $y$ by (ii), so $x=0$; and if $e_1,\dots,e_m$ are nonzero
pairwise orthogonal idempotents, then $S(e_i)\ge1$ (by (ii) and nondegeneracy), $e=\sum e_i$ is idempotent and
$m\le\sum S(e_i)=S(e)\le d$.

By Wedderburn's theorem $R\cong\prod_{i=1}^tM_{n_i}(k)$ \cite{CurtisReiner1962}. A linear form on this product vanishing on
commutators is $\sum_ic_i\galredtr_i$, $\galredtr_i$ the trace of the $i$-th factor, because the commutators of $M_{n}(k)$
span the matrices of trace zero. The matrix unit $e_{11}$ of the $i$-th factor is an idempotent with $S(e_{11})=c_i$, so
$c_i=m_i\in\N$ by (ii), and $d=S(1)=\sum_im_in_i$. The algebra homomorphism $\rho_0\colon\prod M_{n_i}(k)\to M_d(k)$
sending $(x_i)$ to the block diagonal matrix with $m_i$ copies of each $x_i$ has $\galredtr\rho_0(x)=\sum m_i\galredtr x_i$.
The composite $\rho\colon k[\Gamma]\to R\cong\prod M_{n_i}(k)\to M_d(k)$ satisfies $\galredtr\rho=T'$, and its restriction
to $\Gamma$ is a representation $r\colon\Gamma\to\GL_d(k)$ with $\galredtr r=T$. The image of $\rho$ is a quotient of the
semisimple algebra $R$, hence semisimple, so $k^d$ is a semisimple module over it; as the $r(\Gamma)$-stable subspaces are
the submodules for this image, $r$ is semisimple.
\end{proof}

\begin{theorem}\label{pf:galred:limrep}
Let $\Gamma$ be a compact group, $x\colon J\to\Gamma$ a map with dense image, $L\subset\Qlbar$ a finite extension of $\Ql$,
$\theta\colon J\to L$, and $\mathcal C$ a set of continuous representations $\Gamma\to\GL_a(\Qlbar)$ such that $\theta$ is
polynomially approximable by the functions $j\mapsto\galredtr\rho(x_j)$, $\rho\in\mathcal C$. Then there is a continuous
semisimple $r\colon\Gamma\to\GL_a(\Qlbar)$ with $\galredtr r(x_j)=\theta_j$ for all $j$ such that $\galredtr r$ is polynomially
approximable by $\{\galredtr\rho:\rho\in\mathcal C\}$.
\formallabel{GlobalTheory.Cong.exists_limit_rep}
\end{theorem}

\begin{proof}
\emph{Step 1.} For $\varepsilon>0$ and $g\in\Gamma$ there is a neighbourhood $V$ of $g$ with $|\theta_i-\theta_j|\le\varepsilon$
whenever $x_i,x_j\in V$. Indeed, take $\rho_1,\dots,\rho_s\in\mathcal C$ and $c_1,\dots,c_s$ approximating $\theta$ with error
$\varepsilon$, apply this to the polynomial $X_i-X_j$, and choose $V$ such that the continuous function $h\mapsto\sum_k
c_k\galredtr\rho_k(h)$ varies by at most $\varepsilon$ on $V$; then use the ultrametric inequality.

\emph{Step 2.} There is a continuous $T\colon\Gamma\to L$ with $T(x_j)=\theta_j$: by Step 1 the values $\theta_j$, for $x_j$
in the neighbourhoods of $g$, form a Cauchy filter, which converges in the complete field $L$; let $T(g)$ be the limit.
Step 1 at $g=x_j$ gives $T(x_j)=\theta_j$, and for $V$ as in Step 1 and $h\in V$ (open), $|T(h)-T(g)|\le\varepsilon$.

\emph{Step 3.} $T$ is polynomially approximable by $\{\galredtr\rho\}$, with the same $\rho_k,c_k$ as for $\theta$ and the same
error. Let $P$ be an integral polynomial in the variables $X_{h_1},\dots,X_{h_m}$, $h_i\in\Gamma$. The function
$F(y)=P(T(y_1),\dots,T(y_m))-\sum_kc_kP(\galredtr\rho_k(y_1),\dots,\galredtr\rho_k(y_m))$ on $\Gamma^m$ is continuous; at points
$y=(x_{j_1},\dots,x_{j_m})$, which are dense, it is the corresponding expression for $\theta$ and the polynomial obtained by
renaming the variables, so $|F(y)|\le\varepsilon$ there and hence everywhere.

\emph{Step 4.} $T$ is a pseudo-character of dimension $a$: each defining identity has the form $Q(T)=0$ for an integral
polynomial $Q$ in finitely many values with $Q(\galredtr\rho)=0$ for every $a$-dimensional $\rho$; by Step 3,
$|Q(T)|=|Q(T)-\sum_kc_kQ(\galredtr\rho_k)|\le\varepsilon$ for every $\varepsilon$.

\emph{Step 5.} By \cref{pf:galred:ssrep} there is a semisimple $r$ with $\galredtr r=T$. It is continuous. Let
$R\subset M_a(\Qlbar)$ be the span of $r(\Gamma)$; as $\Qlbar^a$ is a faithful semisimple $R$-module, $R$ is semisimple, and
the trace form of $R$ is nondegenerate (the $x$ with $\galredtr(xy)=0$ for all $y\in R$ form an ideal of nilpotent elements,
since $\galredtr(x^m)=0$ for $m\ge1$ in characteristic $0$). Choose $g_1,\dots,g_m$ with $r(g_i)$ a basis of $R$ and the dual
basis $y_i=\sum_la_{il}r(g_l)$; then $r(g)=\sum_i\galredtr(r(g)y_i)r(g_i)$ and $\galredtr(r(g)y_i)=\sum_la_{il}T(gg_l)$ is
continuous in $g$. Finally $\galredtr r(x_j)=T(x_j)=\theta_j$, and Step 3 is the last assertion.
\end{proof}

\subsubsection{Closedness at \texorpdfstring{$\ell$}{l}}

Let $A$ be a place of $\Omega$ above a finite place $w$ of a number field $E$, with a frame at $A$; let $K$ be the finite
extension of $\Ql$ and $D_A\cong G_K$ the identification given by the frame, and $L_0\subset\Qlbar$ a finite extension
containing all conjugates of $K$. For a continuous $\rho\colon G_E\to\GL_a(\Qlbar)$ a \emph{model} over a finite extension
$L\supset L_0$ is a representation on $L^a$ with the same matrices. A \emph{cover} of a continuous semisimple $r$ by a set
$\mathcal C$ of continuous representations consists of: a continuous $\rho\colon G_E\to\GL_{d}(\Zl)$ and $b_1,\dots,b_d\in
M_a(\Qlbar)$, $e\in\Zl^d$ with $r(g)b_j=\sum_i\rho(g)_{ij}b_i$ and $\sum_ie_ib_i=1$ (a lattice in the $\Ql$-span of $r(G_E)$
under left multiplication); for every $n$ a continuous $\rho'_n\colon G_E\to\GL_{d'_n}(\Zl)$ with a map $\Pi_n$ from $\rho'_n$ to
$\rho$ which is equivariant and surjective modulo $\ell^n$, compatible in $n$; and for every $n$ finitely many members
$\rho_{n,1},\dots,\rho_{n,m}$ of $\mathcal C$ such that $\rho'_n$ is a lattice, in the same sense, for $\bigoplus_k\rho_{n,k}$,
together with coefficients $c_{n,k}$ approximating the traces read through $\Pi_n$ up to $\ell^{-n}$.

\begin{theorem}[Closedness at $\ell$]\label{pf:galred:closedAtL}
Let $r\colon G_E\to\GL_a(\Qlbar)$, $a\ge1$, be continuous and semisimple, $A$ a place above $w$ with a frame at $A$, $J(\tau)$
a multiset of integers for every embedding $\tau\colon E\to\Qlbar$, and $\mathcal C$ a set of continuous representations
which are semistable at the frame with $\tau$-Hodge--Tate multiset $J(\tau)$ for every $\tau$ induced by the frame, such that
$r$ is linearly approximable by $\mathcal C$. Then $r$ is de Rham at the frame with $\tau$-Hodge--Tate multiset $J(\tau)$ for
every induced $\tau$, and if every member of $\mathcal C$ is crystalline at the frame, so is $r$.
\formallabel{GlobalTheory.LimitTop.closedAtL}
\end{theorem}

\begin{proof}
We use:
\begin{enumerate}[label=(\roman*),nosep]
\item (\lref{GlobalTheory.TraceConjugationCover.cover_exists}) A continuous semisimple representation of a compact group
  which is linearly approximable by a set of continuous representations has a cover by it.
\item (\lref{GlobalTheory.SstModel.main} and \lref{GlobalTheory.Pst.IsJumpsAt.count_eq}) Semistability of $\rho|_{D_A}$ and its
  Hodge--Tate multiplicities may be computed on any model, through $D_A\cong G_K$.
\item (\lref{GlobalTheory.CovClose.closed_of_cover}) Let $r'$ be a continuous representation of $G_K$ with a cover by a set
  of representations all of whose models over extensions of $L_0$ are semistable with Hodge--Tate multiplicities given by
  $J$. Then every model of $r'$ is de Rham with these Hodge--Tate multiplicities, and it is crystalline if all models of the
  members of the set are. (It is proved with Breuil--Kisin modules, using frames of height bounded in terms of $J$,
  uniformly in $n$, for the lattice representations of the cover; compare \cite{Berger2004,Kisin2008}.)
\item (\lref{GlobalTheory.CompactRepField.exists_finExt_entries_mem} and \lref{GlobalTheory.DeRhamRead.exists_model_of_mem}) A
  continuous representation of a compact group into $\GL_a(\Qlbar)$ has its entries in a finite extension $L$ of $\Ql$,
  which may be taken to contain any given finite extension, and then has a model over $L$.
\end{enumerate}
By (i), $r$ has a cover by $\mathcal C$; restricting along the continuous embedding $G_K\cong D_A\subset G_E$ gives a cover
of $r|_{D_A}$ by the restrictions of the members of $\mathcal C$. By (ii) the hypothesis of (iii) holds, with $J$ transported
to the embeddings of $K$ through the frame. Take a model $r_1$ of $r$ over a finite extension $L\supset L_0$ by (iv),
containing a given finite extension $L'$; by (iii), $r_1|_{D_A}$ is de Rham with the Hodge--Tate multiplicities given by
$J$, and crystalline if all members of $\mathcal C$ are crystalline. An embedding $\tau$ induced by the frame comes from an
embedding of $K$ into some such $L'$, so the $\tau$-Hodge--Tate multiset of $r$ is $J(\tau)$.
\end{proof}

For $g\in W_A$ and a representation $\rho$ of $G_E$ put $t_g(\rho)=\iota^{-1}\galredtr(g\mid\galredWD(\rho|_{D_A}))$, the
image under $\iota^{-1}$ of minus the subleading coefficient of $P_{\galredWD(\rho|_{D_A})}(g)$.

\begin{theorem}\label{pf:galred:closedtr}
In the situation of \cref{pf:galred:closedAtL}, assume that the members of $\mathcal C$ are semistable and crystalline at
the frame with $\tau$-Hodge--Tate multisets $J(\tau)$, and let $g\in W_A$ and $\theta\in\Qlbar$ be such that $r$ and
$\theta$ are jointly approximated by $\mathcal C$ and $t_g$: for every $\varepsilon>0$ there are finitely many
$\rho\in\mathcal C$ and $c_\rho$ with $|\galredtr r(h)-\sum c_\rho\galredtr\rho(h)|\le\varepsilon$ for all $h\in G_E$ and
$|\theta-\sum c_\rho t_g(\rho)|\le\varepsilon$. Then $t_g(r)=\theta$.
\formallabel{UpperTheory.Gal.ClosedPhiFromClosedTrace.ClosedTr_holds}
\end{theorem}

\begin{proof}
By \cref{pf:galred:closedAtL} (joint approximation implies linear approximation), $r$ is crystalline at the frame with
$\tau$-Hodge--Tate multisets $J(\tau)$. Let $\alpha(g)\in\Z$ be such that $g$ acts on the maximal unramified extension of
$\Ql$ as the $\alpha(g)$-th power of the arithmetic Frobenius, $n_1=\max(-\alpha(g),0)$ and $n_2=\max(\alpha(g),0)$. We use:
\begin{enumerate}[label=(\roman*),nosep]
\item (\lref{UpperTheory.Gal.ClosedPhiFromClosedTrace.coverJ_holds}) There is a cover of $r$ by $\mathcal C$ whose
  coefficients also approximate $\theta$: $|\theta-\sum_kc_{n,k}t_g(\rho_{n,k})|\le\ell^{-n}$ for all $n$.
\item (\lref{GlobalTheory.BSP.nonempty_bk_of_finExt} and \lref{GlobalTheory.BSP.exists_isPer}) $K$ has Breuil--Kisin data (a
  Breuil--Kisin ring with an Eisenstein element and its map to $A_{\inf}$) and an associated period.
\item (\lref{UpperTheory.Gal.ClosedPhiFromClosedTrace.frEx_holds}) If the members of $\mathcal C$ are semistable with
  multiplicities $J$ and crystalline, then there are $N$ and crystalline Breuil--Kisin frames of height at most $N$ for the
  lattice representations $\rho'_m$ of the cover and for $\rho$, which are linked: for every $n'$ there is $m\ge n'$ such that
  $\Pi_m$ relates the frames of $\rho'_m$ and $\rho$ modulo $\ell^{n'}$, up to a fixed power of $\ell$.
\item (\lref{UpperTheory.Gal.ClosedPhiFromClosedTrace.fixOf_holds}) The $|\alpha(g)|$-th power of the Frobenius fixes the
  maximal unramified subfield of $K$.
\item (\lref{UpperTheory.Gal.ClosedPhiFromClosedTrace.opEx_holds}) From (iii) and (iv): there are $s\in A_{\mathrm{cris}}$,
  invertible in $B_{\mathrm{cris}}$, and matrices $V$, $V'_m$ over $A_{\mathrm{cris}}$ such that $V\varphi^{n_2}(\Lambda)=s
  \varphi^{n_1}(\Lambda)$ for a matrix $\Lambda$ of crystalline periods of $\rho$, similarly for $V'_m$ and $\rho'_m$, and for
  every $M$ there is $m\ge M$ with $V\Pi_m\equiv\Pi_mV'_m\pmod{\ell^M}$.
\item (\lref{UpperTheory.Gal.ClosedPhiFromClosedTrace.tie_holds}) For a representation $r'$ crystalline at the frame,
  read through a lattice representation as in a cover, and $(s,V)$ as in (v), the vector $V\cdot e$ and the traces of the
  $b_i$ determine $t_g(r')$: $\sum_i(Ve)_i\,\lambda(\galredtr b_i)=s\,\lambda(t_g(r'))$ for every linear map $\lambda$ from
  $\Qlbar$ to the maximal unramified extension of $\Ql$. This applies to $r$ and to every $\rho_{m,k}$.
\item (\lref{UpperTheory.Gal.ClosedPhiFromClosedTrace.trLim_holds}) The identities of (vi) for $r$ and for the
  $\rho_{m,k}$, the congruences of (v) and the approximations of (i) imply $t_g(r)=\theta$.
\end{enumerate}
Take the cover by (i), restrict it to $D_A\cong G_K$, take the data of (ii)--(v), and conclude by (vi) and (vii).
\end{proof}

\subsubsection{Data for the two theta lifts}

We use the notation of \cref{def:galois:tliftk,pf:galred:twoupdef}. For an algebraic Hecke character $\psi$ of $E$ we write
$\psi(\varpi_w)$ for its value at a uniformizer at a place $w$ where it is unramified, and its \emph{infinity type} is the
function $p$ on embeddings $\tau\colon E\to\C$ with $\psi_\infty(x)=\prod_\tau\tau(x)^{p(\tau)}$ on the units in question.
Fix a CM type $(\tau_v)_{v\mid\infty}$ of $E$; the archimedean exponent of a unitary Hecke character is taken with respect
to it.

\begin{proposition}\label{pf:galred:basedata}
Let $E$ be a CM field with $E/E^+$ unramified at all finite places and $\sfname{InertOK}(E)$. There are a CM type, a
standard additive character $\psi_0$ of $\A_{E^+}/E^+$, an element $\delta\in E$ with $\bar\delta=-\delta$ and a fractional
ideal $\mathfrak a$ of $E$ such that: every place of $E^+$ inert in $E$ is prime to $2$, has norm $>3$, and $\psi_0$ is
unramified at it; for every $n$ the hermitian space $(E^n,\sum_ix_i\bar y_i)$ has a self-dual lattice; and for every $d$
the skew-hermitian space $(E^d,\delta\sum_ix_i\bar y_i)$ is definite at the CM type and $\mathfrak a^d$ is a self-dual
lattice in it. These are the base data required by $\sfname{TLiftK}$.
\formallabel{UpperTheory.Gal.TP.Sub.baseNG}
\end{proposition}

\begin{proof}
By \lref{UpperTheory.Gal.TP.Skew.skewLine_holds}, applied to the CM extension $E/E^+$, which is unramified at all
finite places, there are a CM type, a standard additive character $\psi_0$ (with empty set of exceptional places), $\delta$
and $\mathfrak a$ such that $\bar\delta=-\delta$, $\delta$ is positive for the CM type, the standard hermitian spaces have
self-dual lattices, and the skew-hermitian spaces $(E^d,\delta\sum x_i\bar y_i)$ have the self-dual lattices
$\mathfrak a^d$. Let $v$ be a place of $E^+$ inert in $E$. By $\sfname{InertOK}(E)$, $v$ is prime to $2$, has norm $>3$ and
is prime to the different of $E^+$, and a standard additive character is unramified at the places prime to the different
(\lref{PositivityTheory.Mix.psiUnr_of_not_le}).
\end{proof}

\begin{proposition}\label{pf:galred:chrpairs}
Let $E$ be a CM field with $E/E^+$ unramified at all finite places, $n$ even, $t$ odd, $(\tau_v)$ a CM type, $b\in\Z$ and
$w_0$ a split place. There are an algebraic Hecke character $\kappa$ of $E$ with infinity type $m$ and, for $k=1,2$, unitary
Hecke characters $\chi^{(k)}_X,\chi^{(k)}_W$ of $E$ with $\chi^{(k)}_X|_{\A_{E^+}^\times}=\eta^n$,
$\chi^{(k)}_W|_{\A_{E^+}^\times}=\eta^{n+t}$ and archimedean exponents $m^{(k)}_X,m^{(k)}_W$ with $m^{(k)}_W(v)\le b$ for all
$v$, and algebraic Hecke characters $\psi^{(k)}_0,\dots,\psi^{(k)}_{t-1}$ with infinity types $p^{(k)}_i$, such that
$\psi^{(1)}_i\ne\psi^{(2)}_j$ for all $i,j$, $\kappa$ and the $\psi^{(1)}_i$ are unramified at $w_0$, and for each $k$ the
following \emph{dictionary} holds (we omit the index $k$): for every $v\mid\infty$ and $0\le i<t$,
\begin{gather*}
  2m(\tau_v)=m_W(v)-m_X(v)+t,\qquad 2m(\bar\tau_v)=m_X(v)-m_W(v)+t,\\
  2p_i(\tau_v)=m_X(v)+2i-(n+2t-2),\qquad 2p_i(\bar\tau_v)=-m_X(v)+2i-(n+2t-2),
\end{gather*}
and at every finite place $w$ where $\kappa$ and all $\psi_i$ are unramified, $\chi_X$ and $\chi_W$ are unramified and
\[
  \kappa(\varpi_w)=\chi_W(\varpi_w)\chi_X(\varpi_w)^{-1}q_w^{-t/2},\qquad
  \psi_i(\varpi_w)=\chi_X(\varpi_w)\,q_w^{(n+2t-2)/2-i}.
\]
\formallabel{UpperTheory.Gal.TP.Chr.chrPairs_holds}
\end{proposition}

\begin{proof}
By \lref{UpperTheory.Gal.TP.Chr.chrEx_holds}, applied with all exponents equal to $-1$, there are a unitary character
$\chi$ of $\A_E^\times/E^\times$, unramified at $w_0$, and algebraic Hecke characters $\psi(k,s)$ ($k\in\N$, $s\in\Z$) such
that: for every $k$, $\chi^k$ is a Hecke character with archimedean exponent $-k$ at every $v$ whose restriction to
$\A_{E^+}^\times$ is $\eta^j$ for every $j\equiv k\pmod2$; $\psi(k,s)$ is unramified exactly where $\chi^k$ is, with
$\psi(k,s)(\varpi_w)=\chi(\varpi_w)^kq_w^{-s/2}$; for $k+s$ even, $\psi(k,s)$ has infinity type $(s-k)/2$ at $\tau_v$ and
$(s+k)/2$ at $\bar\tau_v$; and $\psi(k,s)=\psi(k',s')$ implies $k=k'$.

Choose an odd $K\ge1$ with $-K\le b$, and for $k=1,2$ put $j=k-1$, $s_i=2i-(n+2t-2)$ and
\begin{gather*}
  \kappa=\psi(K,t),\qquad\chi^{(k)}_X=\chi^{2j},\qquad\chi^{(k)}_W=\chi^{K+2j},\\
  m^{(k)}_X=-2j,\qquad m^{(k)}_W=-(K+2j),\qquad\psi^{(k)}_i=\psi(2j,s_i).
\end{gather*}
As $K+t$ is even, $\kappa$ is algebraic with $2m(\tau_v)=t-K=m_W-m_X+t$ and $2m(\bar\tau_v)=t+K=m_X-m_W+t$. As $2j\equiv n$ and
$K+2j\equiv n+t\pmod2$, the restrictions are $\eta^n$ and $\eta^{n+t}$, and $m_W\le-K\le b$. As $s_i$ is even,
$\psi^{(k)}_i$ is algebraic with $2p_i(\tau_v)=s_i-2j=m_X+2i-(n+2t-2)$ and $2p_i(\bar\tau_v)=s_i+2j=-m_X+2i-(n+2t-2)$. If
$\kappa$ and $\psi_0$ are unramified at $w$, then so are $\chi^K$ and $\chi^{2j}$, hence $\chi_X$ and $\chi_W$, and
$\kappa(\varpi_w)=\chi(\varpi_w)^Kq_w^{-t/2}=\chi_W(\varpi_w)\chi_X(\varpi_w)^{-1}q_w^{-t/2}$,
$\psi_i(\varpi_w)=\chi(\varpi_w)^{2j}q_w^{-s_i/2}=\chi_X(\varpi_w)q_w^{(n+2t-2)/2-i}$. The $\psi^{(1)}_i=\psi(0,s_i)$ differ
from the $\psi^{(2)}_{i'}=\psi(2,s_{i'})$. Since $\chi$ is unramified at $w_0$, so are $\chi^K$ and $\chi^0$, hence $\kappa$
and the $\psi^{(1)}_i$.
\end{proof}

\begin{proposition}\label{pf:galred:upof}
Let $E$ be a CM field with base data as in \cref{pf:galred:basedata}, $1\le n<t$, $\chi_X,\chi_W$ unitary Hecke characters
with archimedean exponents $m_X,m_W$, $U$ automorphic on $\galredU_n$ with weights $\lambda_\tau$, and $U'$ automorphic on
$\galredU_{n+t}$ satisfying (i) and (ii) of \cref{def:galois:tliftk} with $d=n+t$. Let $\kappa,m,\psi_i,p_i$ satisfy the
dictionary of \cref{pf:galred:chrpairs} with $\chi_X,\chi_W,m_X,m_W$. Then $U'$ is a lift of $U$ of size $t$ with twist
$\kappa$ and auxiliary characters $(\psi_i)$ in the sense of \cref{pf:galred:twoupdef}.
\formallabel{UpperTheory.Gal.TP.Sub.upOf}
\end{proposition}

\begin{proof}
\emph{Weights.} Every embedding $E\to\C$ is $\tau_v$ or $\bar\tau_v$ for a unique $v$, and these are distinct. For a
dominant $\mu\in\Z^a$ let $\mu^\vee_k=-\mu_{a-1-k}$; then $J_{\mu^\vee}=(a-1)-J_\mu$. Let $\lambda'_v$ be the weight of $U'$
at $\tau_v$ given by (i), and put $\lambda'_{\tau_v}=\lambda'_v$, $\lambda'_{\bar\tau_v}=(\lambda'_v)^\vee$. The weights of an
automorphic representation of a definite unitary group at conjugate embeddings are dual to each other
(\lref{GlobalTheory.DefiniteBC.hasWeightAt_conjugate}) and unique (\lref{GlobalTheory.DefU.hasWeightAt_unique});
so $U'$ has weight $\lambda'_\tau$ at every $\tau$, and $\lambda_{\bar\tau_v}=\lambda_{\tau_v}^\vee$. Let $\lambda=\lambda_{\tau_v}$,
$\lambda'=\lambda'_v$. By (i), for $k<n$, $k-\lambda'_k=(k-\lambda_k)+(t+m_W-m_X)/2=(k-\lambda_k)+m(\tau_v)$, and for
$k=n+t-1-i$, $0\le i<t$, $k-\lambda'_k=n+t-1-i-(n+m_X)/2=-p_i(\tau_v)$. So
$J_{\lambda'}=(J_\lambda+m(\tau_v))\cup\{-p_i(\tau_v)\}$. Then
$J_{\lambda'^\vee}=(n+t-1)-J_{\lambda'}=(J_{\lambda^\vee}+t-m(\tau_v))\cup\{n+t-1+p_i(\tau_v)\}$, and by the dictionary
$t-m(\tau_v)=m(\bar\tau_v)$ and $n+t-1+p_i(\tau_v)=-p_{t-1-i}(\bar\tau_v)$; this is the identity at $\bar\tau_v$.

\emph{Satake parameters.} Let $w$ be split, with $\kappa$ and the $\psi_i$ unramified at $w$, and $U$ spherical at $w$ with
Satake parameter $B$; write $q=q_w$, $x=\chi_X(\varpi_w)$, $y=\chi_W(\varpi_w)$. By the dictionary $\chi_X,\chi_W$ are
unramified at $w$. By the explicit type II correspondence for unramified representations
(\lref{UpperTheory.Gal.TP.Loc.liftSatU_holds}; here $2n<n+t$), $U'$ is spherical at $w$ with Satake parameter
$B'=\{bxy^{-1}:b\in B\}\cup\{xq^{(t-1)/2-i}:0\le i<t\}$. Multiplying by $q^{(n+t-1)/2}$,
\[
  B'^{(n+t)}_w=\{b'\,q^{t/2}xy^{-1}:b'\in B^{(n)}_w\}\cup\{xq^{(n+2t-2)/2-i}\}_i
  =\{b'\kappa(\varpi_w)^{-1}:b'\in B^{(n)}_w\}\cup\{\psi_i(\varpi_w)\}_i
\]
by the dictionary. The $\psi_i$ are algebraic of infinity type $p_i$, so $U'$ is a lift as claimed.
\end{proof}

% ---------------------------------------------------------------------------
\subsection{Remarks on the formal statements}\label{sec:pf:galred:remarks}

\begin{remark}\label{pf:galred:rem}
\begin{enumerate}[label=(\alph*),nosep]
\item The statement proved is $\sfname{Row0}(\mathrm{recT})$, which is stronger than $\sfname{Row0b}(\mathrm{recT})$ in the
  tame clause: compatibility on inertia holds at all places $w$ with $E_w\cong\Q_g$, without the condition $g>16a^2$. The
  condition of $\sfname{Row0b}$ is not used.
\item In the formal statements the hypothesis ``$\sfname{GeoOne}$'' enters only through $\sfname{SeedFinOdd}$, and in the proof of
  $\sfname{SeedFinOdd}$ the auxiliary set $Z_0$ of places of genericity is empty: the seeds need not be generic at auxiliary
  places, genericity is required only at the places where exact Frobenius compatibility is claimed.
\item The clauses (dR), (crys) and (phi) are quantified over frames; frames inducing a place $A$ exist only above $\ell$, so these
  clauses are vacuous at places prime to $\ell$ (as intended), and (unr) is vacuous above $\ell$.
\item $\galredWD_\iota(r|_{W_A})$ and $\galredWD(r|_{D_A})$ have default values (the trivial representation of dimension $a$,
  resp.\ the trivial class) when the corresponding construction does not apply; the dimension clause of $\sfname{Away}$ holds for
  the default value as well.
\item The definition of a polarized pair requires a nondegenerate invariant pairing between $\pi\otimes\psi^{-1}$ and $\pi$ with
  respect to complex conjugation; the proof of \cref{pf:galred:isweight} uses only the resulting relation between infinitesimal
  characters.
\end{enumerate}
\end{remark}
